\documentclass[reqno,centertags,11pt,draft]{amsart}
%\usepackage[letterpaper,margin=1.3in]{geometry}
\usepackage{amsmath,amsthm,amsfonts,amssymb,enumerate}
%\usepackage{mathptmx}
\usepackage[bookmarksopen=true,final]{hyperref}
\usepackage{mathrsfs}
\usepackage[nobysame]{amsrefs}
%\RequirePackage[british]{babel}
\usepackage{hyperref}
\usepackage[pdftex]{graphicx,color}
%\DeclareGraphicsRule{.pdftex}{pdf}{*}{}
\usepackage{graphicx}
\usepackage{marvosym}
\usepackage{scalerel}
\usepackage{bbm}
\usepackage{bm}
%\usepackage{mathabx}
\usepackage{setspace}
\usepackage{tikz}
%\setstretch{1.25}
\usepackage[foot]{amsaddr}
\usepackage[normalem]{ulem}
\usepackage{cases}

\usepackage{fancyhdr} 
\fancyhf{} 
\cfoot{\thepage}
\pagestyle{plain} 
\setlength{\footskip}{40pt}

\usepackage{ulem} 

\usepackage{arydshln}


\usepackage{xcolor}
%\usepackage{titlesec}
%\titleformat{\section}[block]{\LARGE\bfseries}{}{1em}{}
%\titleformat{\subsection}[block]{\large\bfseries}{}{1em}{}


\usepackage[a4paper,tmargin=4cm,bmargin=4.5cm,lmargin=3.2cm,rmargin=3.2cm]{geometry}
%\usepackage{enumitem}
%\setlength\parindent{0pt}

%%%%%%Theorem types%%%%%%
\newtheorem{thm}{Theorem}[section]
\newtheorem{lem}[thm]{Lemma}
\newtheorem{lemma}[thm]{Lemma}
\newtheorem{cor}[thm]{Corollary}
\newtheorem{prop}[thm]{Proposition}
\newtheorem{rem}[thm]{Remark}
\newtheorem{defn}[thm]{Definition}
\newtheorem{problem}{Problem}
\renewcommand*{\theproblem}{\Alph{problem}}
\newtheorem{ex}[thm]{Example}

%%%%%%Your own commands%%%%%%

\newcommand{\R}{\mathbb{R}}
\newcommand{\C}{\mathbb C}
\newcommand{\Q}{\mathbb{Q}}
\newcommand{\Z}{\mathbb{Z}}
\newcommand{\N}{\mathbb{N}}
\newcommand{\T}{\mathbb{\partial\mathbb{D}}}
\newcommand{\D}{\mathbb{D}}
\newcommand{\K}{{\bf K}}
%\newcommand{\supp}{\operatorname{supp}}
\renewcommand{\t}{\tau}
\renewcommand{\S}{{\mathscr{S}}}
\newcommand{\dd}{\,\mathrm{d} }
\newcommand{\ddd}{\,\text{\rm{\mbox{\dj}}}}
\renewcommand{\d}{\partial}
\newcommand{\supp}{\operatorname{supp}}
\newcommand{\sspan}{\operatorname{span}}

\DeclareMathOperator{\sgn}{sgn}
\DeclareMathOperator{\curl}{curl}
\DeclareMathOperator{\divergence}{div}
\DeclareMathOperator{\diam}{diam}

\newcommand{\abs}[1]{\left|#1\right|}
\newcommand{\set}[1]{\left\{#1\right\}}
\newcommand{\brkt}[1]{\left(#1\right)}
\newcommand{\jap}[1]{\big\langle #1\big\rangle}
\newcommand{\norm}[1]{\left\Vert#1\right\Vert}


%%%%%%Shortcuts for displayed equations%%%%%%

% Math \m
\newcommand{\m}[1]{\begin{equation*}
\begin{split}
#1
\end{split}
\end{equation*}}
% Numbered Math \nm
\newcommand{\nm}[2]{\begin{equation}\label{#1}
\begin{split}
#2
\end{split}
\end{equation}}
% Math Array \ma
\newcommand{\ma}[1]{\begin{align*}
#1
\end{align*}}
% Numbered Math Array \nma
\newcommand{\nma}[2]{\begin{equation}\label{#1}
\begin{aligned}
#2
\end{aligned}
\end{equation}}
% Multiply Numbered Math Array \mnma
\newcommand{\mnma}[1]{\begin{align}
#1
\end{align}}


%%%%%%Colours%%%%%%

\newcommand{\red}[1]{{\color{red}#1}}
\newcommand{\blue}[1]{{\color{blue}#1}}
\newcommand{\white}[1]{{\color{white}#1}}

\newcommand{\rk}[1]{{\color{red}#1}}
\newcommand{\mv}[1]{\marginpar{\color{blue}#1}}

\newcommand{\rkno}[1]{{{\color{red}\sout{#1}}}}


\begin{document}
\title{
Szeg\H{o} mapping and Hermite--Pad\'{e} polynomials for multiple orthogonality on the unit circle
%Two-Point Hermite--Pad\'e Approximation and Multiple Orthogonal Laurent Polynomials on the Unit Circle
}
%\author{Rostyslav Kozhan and Marcus Vaktn\"as}
\date{\today}
\author{Rostyslav Kozhan$^1$}
\email{kozhan@math.uu.se}
\author{Marcus Vaktn\"as$^{1}$}
\email{marcus.vaktnas@math.uu.se}
\address{$^{1}$Department of Mathematics, Uppsala University, S-751 06 Uppsala, Sweden}
\begin{abstract}
We investigate generalized Laurent multiple orthogonal polynomials on the unit circle satisfying simultaneous orthogonality conditions with respect to $r$ probability measures or linear functionals on the unit circle. We show that these polynomials can be characterized as solutions of a general two-point Hermite--Pad\'e approximation problem.

We derive Szeg\H{o}-type recurrence relations, establish compatibility conditions for the associated recurrence coefficients, and obtain Christoffel--Darboux formulas as well as Heine-type determinantal representations.

Furthermore, by extending the Szeg\H{o} mapping and the Geronimus relations, we relate these Laurent multiple orthogonal polynomials to multiple orthogonal polynomials on the real line, thereby making explicit the connection between multiple orthogonality on the unit circle and on the real line.



% \rkno{
% We investigate a general two-point Hermite--Pad\'e approximation problem and show that its solutions are Laurent polynomials satisfying simultaneous orthogonality conditions with respect to  $r$ linear moment functionals $L_1,\ldots,L_r$. In the special case of approximating Carath\'{e}odory functions, these linear functionals correspond to integration with respect to probability measures on the unit circle.

% For the resulting  multiple orthogonal Laurent polynomials, we establish Szeg\H{o}-type recurrence relations, the Christoffel--Darboux formulas, compatibility relations for the recurrence coefficients, and Heine-type determinantal representations. 

% We relate these objects to multiple orthogonal polynomials on the real line by generalizing the Szeg\H{o} mapping and Geronimus relations, thereby obtaining explicit connection between multiple orthogonal polynomials on the real line and multiple orthogonal Laurent  polynomials on the unit circle.
% }
\end{abstract}
\maketitle

\section{Introduction}

While multiple orthogonality on the real line has been extensively developed over the past decades, its analogue on the unit circle, introduced by M\'{i}nguez and Van Assche~\cite{MOPUC1}, still remains a relatively new area of study. Nevertheless, recent work~\cite{MOPUC2,KVMOPUC,KVMLOPUC1,KNik,HueMan} indicates that multiple orthogonality on the unit circle exhibits a rich and intricate structure that warrants a unified and systematic treatment.

In this paper, we introduce a generalized notion of Laurent multiple orthogonal polynomials, which provides a natural framework for the generalized Hermite--Pad\'{e} problem and for connecting with the Szeg\H{o} mapping and Geronimus relations, linking the theory to multiple orthogonal polynomials on the real line. Before presenting our main results, we summarize some foundational aspects of the theory that will be needed throughout the paper. In particular, we summarize key facts about orthogonal polynomials on the real line (Section~\ref{ss:intro1}) and on the unit circle  (Section~\ref{ss:intro2}), their connections to Pad\'e and Hermite--Pad\'e approximation, and the Szeg\H{o} mapping  (Section~\ref{ss:intro3}). A concise overview of our main results is provided in Section~\ref{ss:main}.


%We begin by recalling several classical constructions that will be needed throughout the paper. In particular, we summarize basic facts concerning orthogonal polynomials on the real line and on the unit circle, their connections to Pad\'{e} and Hermite--Pad\'{e} approximation, and the Szeg\H{o} mapping. A summary of our main results is given in Section~\ref{ss:main}.




% \blue{While multiple orthogonality on the real line has been well developed over the last few decades, the unit circle case introduced in \cite{MOPUC1} remains new. However, recent interest in the topic \cite{KVMOPUC,KVMLOPUC1,KNik,HueMan} reveals much more may be said about multiple orthogonality on the unit circle. 

% The goal of the paper was to put our results from the papers \cite{KVMOPUC} and \cite{KVMLOPUC1} under a general framework, by studying multiple orthogonality on the unit circle derived from a two-point Hermite--Pad\'e approximation problem containing every situation studied in previous papers on the topic. 

% The success of the approach is illustrated though the identification of Szeg\H{o} mapping relations in the multiple setting, along with Geronimus relations that follow from extended Szeg\H{o} recurrence relations of \cite{KVMOPUC}.}

% We begin by recalling several classical constructions needed later in the paper. We summarize basic facts about orthogonal polynomials on the real line and on the unit circle, together with their connections to Pad\'{e} and Hermite--Pad\'{e} approximation and the Szeg\H{o} mapping. Our main results are then summarized in Section~\ref{ss:main}.

\subsection{Orthogonality on the real line and Hermite--Pad\'{e} approximation}\label{ss:intro1}
\hfill\\

We start by reviewing the basic theory of orthogonal and multiple orthogonal polynomials on the real line, together with the associated Pad\'{e} and Hermite--Pad\'{e} approximation problems. 
Let $\mu$ be a probability measure supported on the real line $\R$ with all  moments 
$$
c_j = \int_\R x^j d\mu(x), \qquad j\in\N=\{0,1,2,\ldots\},
$$
finite. The Cauchy--Stieltjes transform of $\mu$ is the function, analytic on $\C\setminus\R$, given by
\begin{equation}
    \label{eq:m}
m(z) = \int_{\R} \frac{d\mu(x)}{x-z} = -\sum_{j=0}^\infty c_j z^{-j-1}, \qquad z\to\infty. 
\end{equation}
The monic (i.e., with leading term 1) orthogonal polynomials %$(P_n(x))_{n = 0}^\infty$ 
of $\mu$ are then defined by requiring $\deg{P_n} = n$ and
\begin{equation}\label{eq:oprl1}
    \int_\R P_n(x)x^{k} d\mu(x) = 0, \qquad k = 0,\dots,n-1.
\end{equation}
They satisfy the three-term recurrence relation
\begin{equation}
    x P_{n}(x) = P_{n+1}(x) + b_n P_n(x) + a_n P_{n-1}(x),
\end{equation}
where $a_n>0$, $n\ge 1$, ($a_0:=0$) and $b_n\in\R$ are called the Jacobi coefficients.
See~\cite{SzegoBook,SimonL2} for more details.

One of the classical settings in which orthogonal polynomials naturally appear is Pad\'{e} approximation, see, e.g.,~\cite{Baker} for a comprehensive treatment.
Given a function $m(z)$ as in~\eqref{eq:m}, the Pad\'{e} problem seeks polynomials $P_n$ and $Q_n$ of degrees at most $n$ and $n-1$, respectively, such that
\begin{equation}\label{eq:padeR}
P_n(z) m(z) +Q_n(z) = \mathcal{O}(z^{-n-1}), \qquad z \to \infty.
\end{equation}
It can be shown that all the %polynomial 
solutions $P_n$ of %~\eqref{eq:padeR} with $\deg P_n = n$ 
this problem are exactly the orthogonal polynomials defined by~\eqref{eq:oprl1}.

The Hermite--Pad\'{e} approximation problem concerns simultaneous approximation of Cauchy--Stieltjes transforms, generalizing~\eqref{eq:padeR}. %generalizes the classical Pad\'{e}  approximation~\eqref{eq:padeR} to simultaneous approximation. %by seeking to approximate several functions simultaneously.
Let $m_1(z), \ldots, m_r(z)$ be the Cauchy--Stieltjes transforms of $r$ measures $\mu_1, \ldots, \mu_r$ on $\R$. For a multi-index $\bm{n} = (n_1, \ldots, n_r) \in \N^r$, we look for polynomials $P_{\bm{n}}$ and $Q_{\bm{n},j}$  of degree $\le |\bm{n}|:=n_1+\ldots+n_r$ and $\le |\bm{n}|-1$ satisfying
\begin{equation}\label{eq:hermitePadeR}
P_{\bm{n}}(z) m_j(z) + Q_{\bm{n},j}(z) = \mathcal{O}(z^{-n_j - 1}),
\qquad z \to \infty, \quad j = 1,\ldots,r.
\end{equation}

One of the fundamental results in this theory is that a polynomial $P_{\bm{n}}$ %of degree $\le |\bm{n}|$ 
solves~\eqref{eq:hermitePadeR} if and only if it satisfies simultaneous orthogonality conditions 
\begin{equation}\label{eq:MOPRL}
    \int_\R P_{\bm{n}}(x) x^{k}  d\mu_j(x) = 0,
    \qquad k = 0,\ldots,n_j-1, \quad j = 1,\ldots,r,
\end{equation}
with respect to the system of measures 
$\bm{\mu}=(\mu_1,\ldots,\mu_r)$. Such polynomials are called (type II) multiple orthogonal polynomials of $\bm\mu$.
%These polynomials also appear as the polynomial solutions to a Hermite--Pad\'{e} approximation problem, which generalizes~\eqref{eq:padeR}. Given $r$ Cauchy--Stieltjes transforms $m_1(z),\ldots,m_r(z)$ of $\mu_1,\ldots,\mu_r$, ~\eqref{eq:m}, the Hermite--Pad\'{e} problem seeks polynomials $P_{\bm{n}}$ and $Q_{\bm{n},j}$  of degree $|\bm{n}|$ and $|\bm{n}|-1$ satisfying
%\begin{equation}\label{eq:hermitePadeR}
%P_{\bm{n}}(z) m_j(z) + Q_{\bm{n},j}(z) = \mathcal{O}(z^{-n_j - 1}),
%\qquad z \to \infty, \quad j = 1,\ldots,r.
%\end{equation}
This is a very well developed area of research, see, e.g.,~\cites{Aptekarev, Ismail, Applications,Nikishin} and references therein.

\subsection{Orthogonality on the unit circle and two-point Hermite--Pad\'{e} approximation}\label{ss:intro2}
\hfill\\

%\smallskip

Let us now introduce the parallel theory on the complex unit circle~\cites{OPUC1,OPUC2,GeronimusBook,SzegoBook}. We start with  a probability measure supported on  $\T=\{z\in\C: |z|=1\}$, and denote its moments 
$$
c_j = \int_\T w^{-j} d\mu(w), \qquad j\in\Z.
$$
The Carath\'{e}odory function of $\mu$ is defined to be the analytic on $\C\setminus\T$ function $F(z)$ given by
\begin{equation}\label{eq:F}
    F(z) = \int_\T \frac{w+z}{w-z} d\mu(w)
    =
    \begin{cases}
         c_{0} + 2\sum_{k = 1}^{\infty}c_{k}z^k, & z\to0,
         \\
        -c_{0} - 2\sum_{k = 1}^{\infty}c_{-k}z^{-k}, & z\to\infty.
    \end{cases}
\end{equation}
The monic orthogonal polynomials $(\Phi_n(z))_{n = 0}^\infty$ of $\mu$ are defined by requiring $\deg{\Phi_n} = n$ and
\begin{equation}\label{eq:opuc1}
    \int_\T \Phi_n(w)w^{-k} d\mu(w) = 0, \qquad k = 0,\dots,n-1.
\end{equation}
They satisfy the Szeg\H{o} recurrence relations given by
\begin{align}
\label{eq:scalar szego 1}
\Phi_{n+1}(z) & = z\Phi_n(z) + {\alpha}_{n+1}\Phi^*_{n}(z),
\\
\label{eq:scalar szego 2}
\Phi^*_{n+1}(z) & =  \Phi^*_n(z) + \bar\alpha_{n+1}z\Phi_{n}(z),
\end{align}
where the Szeg\H{o} dual polynomials $\Phi_n^*(z) = z^n\overline{\Phi_n(1/\bar z)}$ can be defined via $\Phi_n^*(0)=1$ and
\begin{equation}\label{eq:Phi^*}
    \int_\T \Phi^*_n(w)w^{-k} d\mu(w) = 0, \qquad k = 1,\dots,n.
\end{equation}
The recurrence coefficients $\alpha_n$, $n\ge 1$, belong to open unit disc $\D = \set{z \in \C : \abs{z} < 1}$, and are known as the Verblunsky coefficients of $\mu$ (also referred to in the literature as Schur, Szeg\H{o}, or Geronimus coefficients).\footnote{We warn the reader that here we are choosing $\alpha_n = \Phi_n(0)$ instead of the nowadays more commonly used $\alpha_n = \overline{\Phi_{n+1}(0)}$ %Note that we are using $\alpha_{n+1}$ in place where it is traditional (nowadays) to use $-\bar\alpha_n$ 
(see the historical discussion in~\cite[p.10]{OPUC1}).}

The associated Pad\'{e} approximation problem takes the following {\it two-point} form, see, e.g.~\cite{Baker,Bultheel,JNT,PeherstorferSteinbauer} and references therein.
Given a function $F(z)$ as in~\eqref{eq:F}, we seek polynomials $\Phi_n$ and $\Psi_n$ of degrees at most $n$, such that
\begin{alignat}{3}
\label{eq:padeT1}
    & \Phi_{{n}}(z)F(z) + \Psi_{{n}}(z) = \mathcal{O}(z^{n}),  \qquad&& z \rightarrow 0, 
    \\
    \label{eq:padeT2}
    & \Phi_{{n}}(z)F(z) + \Psi_{{n}}(z) = \mathcal{O}(z^{-1}), \qquad && z \rightarrow \infty.
\end{alignat}
All polynomial solutions $\Phi_n$ of~\eqref{eq:padeT1}--\eqref{eq:padeT2} %with $\deg \Phi_n = n$
are constant multiples of the orthogonal polynomials defined by~\eqref{eq:opuc1}.

In~\cite{MOPUC1}, M\'{i}nguez and Van Assche considered the Hermite--Pad\'{e} approximation problem of finding 
polynomials $\Phi_{\bm{n}}(z),\Psi_{\bm{n},1}(z),\ldots,\Psi_{\bm{n},r}(z)$ of degree $\le|\bm{n}|$ that satisfy
\begin{alignat}{3}
    \label{eq:hermitePadeT1}
    & \Phi_{\bm{n}}(z)F_j(z) + \Psi_{\bm{n},j}(z) = \mathcal{O}(z^{n_j}), \qquad  && z \rightarrow 0, 
    \\
    \label{eq:hermitePadeT2}
    & \Phi_{\bm{n}}(z)F_j(z) + \Psi_{\bm{n},j}(z) = \mathcal{O}(z^{-1}), \qquad  &&z \rightarrow \infty,
\end{alignat}
for $r$ Carath\'{e}odory functions~\eqref{eq:F} $F_1,\ldots,F_r$ associated with measures $\mu_1,\ldots,\mu_r$  on $\T$, and showed that $\Phi_{\bm{n}}(z)$ %with $\deg\Phi_{\bm{n}}\le |\bm{n}|$ 
solves~\eqref{eq:hermitePadeT1}--\eqref{eq:hermitePadeT2} if and only if it
%$\deg\Phi_{\bm{n}}\le |\bm{n}|$ 
fulfills the simultaneous orthogonality relations
\begin{equation}\label{eq:MOPUC}
    \int_\T \Phi_{\bm{n}}(w) w^{-k}  d\mu_j(w) = 0,
    \qquad k = 0,\ldots,n_j-1, \quad j = 1,\ldots,r.
\end{equation}
The polynomials $\Phi_{\bm{n}}(z)$ are therefore  called (type II) multiple orthogonal polynomials on the unit circle. See~\cites{MOPUC1,MOPUC2,KVMOPUC} for further results on the properties of $\Phi_{\bm{n}}$. %, including existence and uniqueness properties, recurrence formulas, CD formulas, Riemann--Hilbert problem, among others.

\subsection{Szeg\H{o} mapping and Geronimus relations}\label{ss:intro3}
\hfill\\

%Recall from Szeg\H{o}'s classical work, that, 
Finally, we recall the Szeg\H{o} mapping and the Geronimus relations, which provide a fundamental connection between the theories of orthogonal polynomials on the real line and on the unit circle. Given any probability measure $\gamma$ supported on the real interval $[-2,2]$, one can define a  probability measure $\mu=\operatorname{Sz}(\gamma)$ to be the probability measure on the unit circle $\T$, that is invariant under the reflection $e^{i\theta}\mapsto e^{-i\theta}$ and satisfies
\begin{equation}\label{eq:szego mapIntro}
    \int_{\T} g(2\cos\theta) d\mu(e^{i\theta}) = \int_{-2}^2 g(x) d\gamma(x),
\end{equation}
for all measurable functions $g$ on $[-2,2]$. Conversely,  any probability measure $\mu$ on $\T$ that is invariant under $e^{i\theta}\mapsto e^{-i\theta}$ arises as   $\mu=\operatorname{Sz}(\gamma)$ for a unique probability measure $\gamma$ on $[-2,2]$  determined by~\eqref{eq:szego mapIntro}. Szeg\H{o} showed~\cite{Szego} that the orthogonal polynomials $P_n$ on the real line with respect to $\gamma$ and the orthogonal polynomials $\Phi_n$ on the unit circle with respect to $\mu$ are connected by the identities
\begin{equation}
     P_n(z+z^{-1}) = \frac{1}{1+\alpha_{2n}} \big( z^{-n} \Phi_{2n}(z) + z^{n} \Phi_{2n}(1/z) \big) 
     %\\
     %&  =   z^{-n+1} \Phi_{2n-1}(z) + z^{n-1} \Phi_{2n-1}(1/z) .
\end{equation}
Later, Geronimus~\cite{GeronimusRelations,GeronimusBook} derived explicit formulas relating the recurrence coefficients of $\gamma$ and  of $\mu$:
\begin{alignat}{3}
\label{eq:GerIntro1}
    a^2_n &  = (1+\alpha_{2n-2})(1-\alpha_{2n-1}^2)(1-\alpha_{2n}),
    \qquad && n\ge 1,
    \\
\label{eq:GerIntro2}
    b_n & =(1-\alpha_{2n})\alpha_{2n-1} - (1+\alpha_{2n})\alpha_{2n+1}, \qquad && n\ge 0,
\end{alignat}
with the convention $\alpha_{-1} = 0$. %, $\alpha_0 = 1$.
%These relations provide an important bridge between the theories of orthogonal polynomials on the real line and on the unit circle.


\subsection{Summary of main results}\label{ss:main}
\hfill\\

In this paper we introduce the Laurent multiple orthogonal polynomials defined by
\begin{align}
\label{eq:generalized hermite pade orthogonality intro}
    & \int_\T \Phi_{\bm{n};\bm{m}}(w)w^{-k}d\mu_j(w) = 0, \qquad k = -m_j,\dots,n_j-1, \qquad j = 1,\dots,r,
    \\
    \label{eq:IntroSpanNEW}
    & \Phi_{\bm{n};\bm{m}} \in \operatorname{span}\set{z^{-\abs{\bm{m}}},z^{-\abs{\bm{m}}+1},\dots,{z^{\abs{\bm{n}}}}}.
\end{align}
Here $\bm{n}$ and $\bm{m}$ are two arbitrary multi-indices in $\N^r$. In particular, if $\bm{m}=\bm{0}$, then $\Phi_{\bm{n};\bm{0}}$ are the 
%type II 
multiple orthogonal polynomials $\Phi_{\bm{n}}$ of M\'{i}nguez--Van Assche~\cite{MOPUC1}.

The motivation for studying this class of polynomials is two-fold. First, we show (see Sections~\ref{ss:general type II hermite pade}--\ref{ss:HPI}) that a Laurent polynomial $\Phi_{\bm{n};\bm{m}}(z)$ satisfying~\eqref{eq:IntroSpanNEW} fulfills the orthogonality conditions~\eqref{eq:generalized hermite pade orthogonality intro} if and only if it solves the two-point Hermite--Pad\'{e} approximation problem 
\begin{alignat}{3}
    \label{eq:intro generalized two point hermite pade eq 1}
    & \Phi_{\bm{n};\bm{m}}(z)F_j(z) + \Psi_{\bm{n};\bm{m},j}(z) = \mathcal{O}(z^{n_j}),  \qquad &&z \rightarrow 0, \\
    \label{eq:intro generalized two point hermite pade eq 2}
    & \Phi_{\bm{n};\bm{m}}(z)F_j(z) + \Psi_{\bm{n};\bm{m},j}(z) = \mathcal{O}(z^{-m_j-1}), \qquad &&z \rightarrow \infty,
\end{alignat}
for suitable Laurent polynomials $\Psi_{\bm{n};\bm{m},j}$. Here $F_j$ are the Carath\'{e}odory functions associated with the measures $\mu_1,\ldots,\mu_r$.
% Here $\bm{n}$ and $\bm{m}$ are two arbitrary multi-indices in $\N^r$. In Section~\ref{ss:general type II hermite pade} we show that $\Phi_{\bm{n};\bm{m}}(z)$ in~\eqref{eq:IntroSpan} solves ~\eqref{eq:intro generalized two point hermite pade eq 1}--\eqref{eq:intro generalized two point hermite pade eq 2} if and only if it satisfies
% \begin{equation}\label{eq:generalized hermite pade orthogonality intro}
%     \int_\T \Phi_{\bm{n};\bm{m}}(w)w^{-k}d\mu_j(w) = 0, \qquad k = -m_j,\dots,n_j-1, \qquad j = 1,\dots,r.
% \end{equation}
In the scalar case $r=1$, one readily verifies that $\Phi_{n;m}(z) = z^{-m} \Phi_{n+m}(z)$.
%by simply multiplying~\eqref{eq:padeT1} and~\eqref{eq:padeT2} by replacing $n$ with $n+m$ and multiplying each equation by $z^{-m}$.
If $r\ge 2$, such a trivial relationship with $\Phi_{\bm{n}}$ in~\eqref{eq:MOPUC}, of course, no longer holds. 
%properties of our Laurent polynomials $\Phi_{\bm{n};\bm{m}}$ in~\eqref{eq:generalized hermite pade orthogonality intro} no longer have such a trivial relationship with the properties $\Phi_{\bm{n}}$ in~\eqref{eq:MOPUC} are no longer so trivially related. 

Our second motivation for introducing the Laurent framework~\eqref{eq:generalized hermite pade orthogonality intro} is that it naturally accommodates the Szeg\H{o} mapping and the Geronimus relations. This allows us to connect our results with the theory of multiple orthogonal polynomials on the real line, a connection that is not available in the multiple polynomial setting~\eqref{eq:MOPUC} considered in~\cite{MOPUC1,MOPUC2,KVMOPUC}. These connections are established in Theorems~\ref{thm:szego mapping type II} and~\ref{thm:Geronimus} in Section~\ref{ss:SzegoMap}.

To achieve this goal, we first develop the underlying framework, introducing the relevant notions and establishing several structural results that will be used throughout the paper. In Section~\ref{ss:SzegoNEW}, we show that the polynomials $\Phi_{\bm{n};\bm{m}}$ satisfy a family of nearest-neighbour recurrence relations, which can be viewed as a natural extension of the classical Szeg\H{o} recurrences for orthogonal polynomials on the unit circle. In Section~\ref{ss:CC} we derive a system of nearest-neighbour compatibility relations satisfied by the recurrence coefficients and the orthogonal polynomials. Section~\ref{ss:extra} is devoted to a collection of identities that follow from these relations; although somewhat technical in nature, they provide the main tools required for the proof of the Geronimus relations. In Section~\ref{ss:Heine} we establish analogues of the Heine determinantal formulas for the Laurent polynomials $\Phi_{\bm{n};\bm{m}}$, as well as for all associated recurrence coefficients. Section~\ref{ss:CD} contains the Christoffel--Darboux formula.
The results of Sections~\ref{ss:SzegoNEW}, \ref{ss:CC}, and~\ref{ss:CD} follow the same pattern as for the multiple orthogonal polynomials $\Phi_{\bm{n}}$, with the main distinction that $\Phi_{\bm{n};\bm{m}}$ involves two independently varying multi-indices. 
%We also show an analogue of the Christoffel--Darboux formula. 
The corresponding proofs follow closely arguments developed in our~\cite{KVMOPUC}, with appropriate modifications to accommodate the present more general framework. 
%This constitutes Sections~\ref{ss:SzegoNEW} and~\ref{ss:CD}. In Section~\ref{ss:CC} we establish the so-called compatibility relations of the recurrence coefficients, and in Section~\ref{ss:Heine} the Heine-type determinantal formulas for the Laurent polynomials $\Phi_{\bm{n};\bm{m}}$ and all the recurrence coefficients.


%Apart from the Hermite--Pad\'{e} problem~\eqref{eq:intro generalized two point hermite pade eq 1}--\eqref{eq:intro generalized two point hermite pade eq 2}, our second motivation for studying the Laurent multiple orthogonal polynomials on the unit circle was in extending the Szeg\H{o} mapping~\cite{Szego} to multiple setting.

% A major benefit of the Laurent framework~\eqref{eq:generalized hermite pade orthogonality intro} is that it enables a natural connection, via the Szeg\H{o} mapping and Geronimus relations, to the theory of multiple orthogonal polynomials on the real line --- something not available in the multiple polynomial~\eqref{eq:MOPUC} approach of~\cite{MOPUC1,MOPUC2,KVMOPUC}.

% %One of the big differences/benefits/reason to study Laurent orthogonal polynomials~\eqref{eq:generalized hermite pade orthogonality intro} compared to the multiple orthogonal polynomials $\Phi_{\bm{n}}$ from~\eqref{eq:MOPUC} in the approach of~\cite{MOPUC1,MOPUC2,KVMOPUC}, is the possibility to connect this theory to the multiple orthogonal polynomials on the real line via the Szeg\H{o} mapping and Geronimus relations. 


% It is natural to seek an analogue of these results in the multiple orthogonality setting. 
% However, such an extension cannot be formulated in terms of the usual multiple orthogonal polynomials $\Phi_{\bm{n}}$ from~\eqref{eq:MOPUC}; instead, one must work with the Laurent analogues~\eqref{eq:generalized hermite pade orthogonality intro}. Indeed, given $\mu_j=\operatorname{Sz}(\gamma_j)$, we establish that
% \begin{align}\label{eq:SzegoPolynomialRelation}
%           P_{\bm{n}}(z+z^{-1}) & = \frac{1}{(1+\alpha_{\bm{n};\bm{n}})} \Big( \Phi_{\bm{n};\bm{n}}(z) + \Phi_{\bm{n};\bm{n}}(1/z) \Big) 
%           \\
%           \label{eq:SzegoPolynomialRelation2}
%           %& = \Phi_{\bm{n};\bm{m}}(z) + \Phi_{\bm{n};\bm{m}}(1/z)
%           & = \Phi_{\bm{n};\bm{n}-\bm{e}_j}(z) + \Phi_{\bm{n};\bm{n}-\bm{e}_j}(1/z),
%     \end{align}
%     where $P_{\bm{n}}$ are the monic multiple orthogonal polynomials~\eqref{eq:MOPRL} for the system $\bm{\gamma}=(\gamma_1,\ldots,\gamma_r)$, $\Phi_{\bm{n};\bm{m}}$ are the monic Laurent multiple orthogonal polynomials~\eqref{eq:generalized hermite pade orthogonality intro} for $\bm{\mu}=(\mu_1,\ldots,\mu_r)$, and $\alpha_{\bm{n},\bm{m}}$ is the $z^{-|\bm{m}|}$ coefficient of $\Phi_{\bm{n};\bm{m}}$, playing the role of a generalized Verblunsky coefficient. We also obtain a generalization of the Geronimus relations that express the nearest-neighbour recurrence coefficients of $\bm{\gamma}$ and the nearest-neighbour recurrence coefficients for $\bm{\mu}$.
%     See Section~\ref{ss:SzegoMap} for more details.

We also note that in this work we adopt the broader setting, where orthogonality with respect to measures is replaced by orthogonality with respect to moment linear functionals. This generalization does not introduce any additional technical complications, and all statements and proofs extend naturally to it.  Readers who prefer the classical setting may simply interpret each functional as integration with respect to a measure. 

% In this formulation, one quickly notices that to preserve the proper duality, the definition of type I multiple orthogonal polynomials must be slightly modified compared to~\cite{MOPUC1,KVMOPUC}. Specifically, when working with linear functionals, the type I polynomials should be taken as polynomials in $z^{-1}$, rather than in $z$ (see ~\eqref{eq:moprlIspan},~\eqref{eq:moprlI*span} with $\bm{m}=\bm{0}$). This provides yet another motivation for our Laurent polynomial framework. See Remarks~\ref{rem:annoying1},~\ref{rem:annoying2} for more details. 



Laurent orthogonal polynomials have also been considered in several very recent works~\cites{KVMLOPUC1,KNik,HueMan}. 
The polynomials in~\cite{HueMan} correspond to the mixed multiple orthogonality setting. While their work studies recurrence relations, Christoffel--Darboux formula, and spectral transformations, only for very specific multi-indices do their polynomials coincide with ours, because their focus is restricted to step-line indices. Furthermore, because they do not introduce the II${}^*$ and I${}^*$ families, their formulas take a substantially different form.
The polynomials $\varphi_{2\bm{n}}$ in~\cite{KVMLOPUC1,KNik} coincide with the Laurent polynomials $\Phi_{\bm{n};\bm{n}}$ studied here, and normality (existence and uniqueness) for $\varphi_{\bm{n}}$ was established for broad classes of measures, including Angelesco, AT, and Nikishin systems on the unit circle. Notice that this establishes existence and uniqueness of $\Phi_{\bm{n};\bm{n}}$ and of the Hermite--Pad\'{e} approximants for ~\eqref{eq:intro generalized two point hermite pade eq 1}--\eqref{eq:intro generalized two point hermite pade eq 2} in the case $\bm{n}=\bm{m}$. The extent to which these results generalize to arbitrary indices $(\bm{n};\bm{m})$ is currently an interesting open problem. Another natural question concerns the location of zeros, which is closely related to the issue of normality, as indicated by the ideas in~\cite{KVInterlacing}. We leave these questions for future work.

\subsubsection*{Acknowledgements}
R.K. is grateful to  Maxim Derevyagin, Brian Simanek, and Sergey Denisov for useful discussions regarding Szeg\H{o}'s mapping. 


%Laurent orthogonal polynomials are also the topic of a few other papers~\cites{KVMLOPUC1,KNik,HueMan} that were announced at the same time. The $\varphi_{2\bm{n}}$ polynomials studied in our~\cite{KVMLOPUC1,KNik} coincide with the Laurent polynomials $\Phi_{\bm{n};\bm{n}}$ from this paper. At such locations \cites{KVMLOPUC1,KNik} established normality, that is, existence and uniqueness of $\varphi_{\bm{n}}$, for wide classes of measures, called Angelesco, AT, and Nikishin. It would be very interesting to find examples of measures that have existence and uniqueness of every $\Phi_{(\bm{n};\bm{m})}$. The Laurent polynomials considered in~\cite{HueMan} are for a more general mixed multiple orthogonality setting but restricted to the step-line indices. They investigate recurrence relations, CD formulas, and Christoffel/Geronimus transforms. Only for very special type of indices their polynomials coincide with ours, so the overlap of their results with our Sections~\ref{ss:SzegoNEW},~\ref{ss:CD} is minimal. Furthermore, since they do not work with $II^*$ and $I^*$ polynomials, their formulas take entirely different form. 

%%%%%%%%%%%%%%%%%%%%%%%%%%%%%%%%%%%%%%%%%


%We say that $(\bm{n};\bm{m}) \in \N^r \times \N^r$ is normal if $\Phi_{\bm{n},\bm{m}}$ is uniquely determined up to multiplication by a non-zero constant and $\deg{z^{\abs{\bm{m}}}\Phi_{\bm{n};\bm{m}}} = \abs{\bm{n}}+\abs{\bm{m}}$. When $\bm{m} = \bm{0}$ this corresponds to the usual definition of normality. For multiple orthogonal polynomials on the real line, Angelesco systems and AT systems provide the two most well-known classes of systems with all indices normal, see e.g. \cite{Nikishin, Ismail}. In \cite{KVMLOPUC1}, we identified unit circle analogues of Angelesco system and AT systems, and proved that indices of the form $(\bm{n};\bm{n})$ are normal (note that $\Phi_{\bm{n};\bm{n}}$ is $\phi_{2\bm{n}}$ in the notation of \cite{KVMLOPUC1}).  In \cite{KNik}, the same was proven for a unit circle analogue of Nikishin systems for the case $r=2$. %\rkno{It is still not clear if normality also holds for $(\bm{n};\bm{m})$ when $\bm{n} \neq \bm{m}$.} 
%{It is an interesting open problem whether $(\bm{n};\bm{m})$ with $\bm{n} \neq \bm{m}$ are always normal for Angelesco, AT, and Nikishin systems on the unit circle.}

%More generally, for a linear functional $L$~\eqref{eq:functional} on Laurent polynomials, we may work with the formal power series
%\begin{equation}\label{eq:formal power series}
%    F^{(0)}(z) = c_{0} + 2\sum_{k = 1}^{\infty}c_{k}z^k, \quad F_j^{(\infty)}(z) = -c_{0} - 2\sum_{k = 1}^{\infty}c_{-k}z^{-k}
%\end{equation}
%where $c_k$ are the moments
%\begin{equation}
%    c_k = L[w^{-k}], \qquad k \in \Z.
%\end{equation}
%We may formally write 
%\begin{equation}\label{eq:generalized caratheodory function}
%    F(z) = L[(w+z)(w-z)^{-1}]
%\end{equation} to represent the pair $(F^{(0)},F^{(\infty)})$.% 
%Then \eqref{eq:functional} is equivalent \cite{PeherstorferSteinbauer, JNT} to 
%\begin{align}\label{eq:generalized two point pade eq 1}
%    & \Phi_{\bm{n}}(z)F_j^{(0)}(z) + \Psi_{\bm{n}}(z) = \mathcal{O}(z^{n}),  &z& \rightarrow 0, \\\label{eq:generalized two point pade eq 2}
 %   & \Phi_{\bm{n}}(z)F_j^{(\infty)}(z) + \Psi_{\bm{n}}(z) = \mathcal{O}(z^{-1}), &z& \rightarrow \infty,
%\end{align}
%for some polynomials $\Psi_n$, which will necessarily be given by
%\begin{equation}
 %   \Psi_n(z) = L\Big[\frac{w+z}{w-z}(\Phi_n(w)-\Phi_n(z))\Big] + L[\Phi_n(w)].
%\end{equation}
%$\Psi_n$ are usually referred to as the \rkno{orthogonal polynomials of second kind} \rk{second kind polynomials}. 

%Nearly all papers on multiple orthogonal polynomials concern measures supported on the real line (see \cites{Aptekarev, Ismail, Applications,Nikishin} for an introduction to the subject). Until very recently, only two papers about multiple unit circle measures had been written \cite{MOPUC1,MOPUC2}. Along with this paper, \cite{HueMan, KVMLOPUC1, KVMOPUC, KNik} contain recent developments. Each of \cite{HueMan, KVMLOPUC1, KNik} discuss different variations of \eqref{eq:generalized hermite pade orthogonality intro}, rather than the orthogonalities \eqref{eq:multiple orthogonality unit circle} that were the topic of \cite{KVMOPUC, MOPUC1, MOPUC2}. 

% In Section 2, we discuss the orthogonality relations in detail, in the more general context of orthogonality with respect to linear functionals. In Section 3, we discuss type II multiple orthogonal Laurent polynomials on the unit circle in \eqref{eq:generalized hermite pade orthogonality intro}. In section 4 we prove the same equivalence for type I multiple orthogonal Laurent polynomials on the unit circle. In Section 5 we discuss recurrence relations. Our generalized Szeg\H{o} recurrence relations in \cite{KVMOPUC} naturally extend to these Laurent polynomials. For the recurrence relations in \cite{MOPUC2}, their main result extends to this setting as well. Lastly, in section 6 we extend relations for the Szeg\H{o} mapping \cite{Szego, OPUC2} in the one measure case to indices of the form $(\bm{n};\bm{n})$, as well as nearby indices. This establishes a connection between multiple orthogonal polynomials on the real line and the unit circle, and provides another motivation for studying multiple orthogonality for  Laurent polynomials. 


\section{Orthogonality Relations and Normality}\label{ss:normality}

Let $r$ sequences $(c_{k,j})_{k\in\Z}$, $j=1,\ldots,r$, of arbitrary complex numbers be given; we refer to them as moments. We define the corresponding moment functionals $L_j$ as the linear maps on the space of Laurent polynomials with complex coefficients, uniquely determined by
% Let us assume that we are given $r$ sequences $\{c_{k,j}\}_{k\in\Z}$, $j=1,\ldots,r$, of arbitrary complex numbers which will be referred to as the moments. Define the corresponding moment functionals $M_j$ to be the linear map on the space of all Laurent polynomials (with complex coefficients), uniquely determined by 
\begin{equation}\label{eq:moments linear functional}
         L_j[w^{-k}] = c_{k,j} , \qquad k\in\Z, \qquad j = 1,\dots,r. 
    \end{equation}

Of special importance is the case when a functional $L_j$ is induced by a probability measure on the unit circle $\T$
\begin{equation}\label{eq:functionals given by measures}
        L_j[w^{-k}] = \int_{\T} w^{-k} d\mu_j(w), \qquad 
        k \in \Z. % \qquad j = 1,\dots,r, 
\end{equation}

%We want to investigate orthogonal Laurent polynomials.... 

In what follows let $\Z$ be the set of integers, $\N$ be the set of non-negative integers. Given two vectors $\bm{u},\bm{v}\in \Z^r$, we write $\bm{u}\ge \bm{v}$ if $u_j\ge v_j$ for all $j=1,\ldots,r$.
%, and we write $\bm{u}\gneq \bm{v}$ if $\bm{u}\ge \bm{v}$ but $\bm{u}\ne \bm{v}$. 
Finally, we denote $|\bm{v}| = v_1+\ldots + v_r$. Note that we do {\it not} take the absolute value of $v_j$, even though $v_j$ may be negative; this is done on purpose. In particular, $\bm{u}\ge \bm{v}$ implies $|\bm{u}|\ge |\bm{v}|$.

Given two $\Z^r$ indices $\bm{n}$, $\bm{m}$, let us denote 
\begin{equation}
     \mathfrak{C}_{2r}
     =
     \left\{
     (\bm{n};\bm{m}) \in\Z^r\times\Z^r : 
     n_j+m_j\ge 0 \mbox{ for all } j
     \right\}
\end{equation}
%$(\bm{n};\bm{m})\in \mathfrak{C}_{2r}$ to denote the cone $\bm{n}+\bm{m}\ge \bm{0}$ (i.e., $n_j+m_j\ge 0$ for all $j$).
for the remainder of the text.
    
For $(\bm{n};\bm{m})\in\mathfrak{C}_{2r}$, define the $(|\bm{n}|+|\bm{m}|)\times (|\bm{n}|+|\bm{m}|)$ matrix
\begin{equation}\label{eq:T}
    T_{\bm{n};\bm{m}} =\left( \begin{array}{ccc}
        c_{\abs{\bm{m}}-m_1,1} & \cdots & c_{-\abs{\bm{n}}-m_1+1,1} \\
        c_{\abs{\bm{m}}-m_1+1,1} & \cdots & c_{-\abs{\bm{n}}-m_1+2,1} \\
        \vdots & \ddots & \vdots \\
        c_{\abs{\bm{m}}+n_1-1,1} & \cdots & c_{-\abs{\bm{n}}+n_1,1} 
        \\[0.2em]
        %\hline
        \hdashline[0.5pt/2pt]
        \\[-1.2em]
        & \vdots \\[0.2em]
        %\hline
        \hdashline[0.5pt/2pt]
        \\[-1.0em]
        c_{\abs{\bm{m}}-m_r,r} & \cdots & c_{-\abs{\bm{n}}-m_r+1,r} \\
        c_{\abs{\bm{m}}-m_r+1,r} & \cdots & c_{-\abs{\bm{n}}-m_r+2,r} \\
        \vdots & \ddots & \vdots \\
        c_{\abs{\bm{m}}+n_r-1,r} & \cdots & c_{-\abs{\bm{n}}+n_r,r}
    \end{array}\right). 
\end{equation}
If $n_j=-m_j$ for some $j$, then the corresponding $j$-th block in the above matrix is empty. Furthermore, if $\bm{n}=-\bm{m}$, then we formally take $T_{\bm{n};-\bm{n}} = 1$.

We say that $(\bm{n};\bm{m})\in \mathfrak{C}_{2r}$ %(with $\bm{n}+\bm{m}\ge \bm{0}$) 
is normal if $\det T_{\bm{n};\bm{m}} \ne 0$. It should not be surprising that this notion is related to the uniqueness of appropriately normalized Laurent orthogonal polynomials, as described next.

\begin{prop}\label{prop:normality}
Let $\bm{L} = (L_1,\dots,L_r)$ be a system of linear functionals and $(\bm{n};\bm{m})\in \mathfrak{C}_{2r}$, with $\bm{n}\ne -\bm{m}$. 
\begin{itemize}
    \item[{(i)}] $\det T_{\bm{n};\bm{m}} \ne 0$ if and only if there exists a unique Laurent polynomial $\Phi_{\bm{n};\bm{m}}$ of the form
    \begin{equation}\label{eq:moprlIIspan}
    \Phi_{\bm{n};\bm{m}}(z) = z^{\abs{\bm{n}}} + \ldots + \alpha_{\bm{n};\bm{m}}z^{-\abs{\bm{m}}},
    \end{equation}  
    that satisfies 
    \begin{equation}\label{eq:moprlII}
    L_j[\Phi_{\bm{n};\bm{m}}(w)w^{-k}] = 0, %\quad k = -m_j,\dots,n_j-1, \quad j = 1,\dots,r.
     \quad -m_j \le k \le  n_j-1, \quad 1\le j \le r.
    \end{equation}
\item[{(ii)}] $\det T_{\bm{n};\bm{m}} \ne 0$ if and only if there exists a unique Laurent polynomial $\Phi^*_{\bm{n};\bm{m}}$ of the form
    \begin{equation}\label{eq:moprlII*span}
    \Phi_{\bm{n};\bm{m}}^*(z) = \beta_{\bm{n};\bm{m}}z^{\abs{\bm{n}}} + \ldots + z^{-\abs{\bm{m}}}
\end{equation} 
    that satisfies 
    \begin{equation}\label{eq:moprlII*}
        L_j[\Phi_{\bm{n};\bm{m}}^*(w)w^{-k}] = 0, %\quad k = -m_j+1,\dots,n_j, \quad j=1,\ldots,r.
        \quad -m_j+1\le k \le n_j, \quad 1\le j \le r.
\end{equation}
     \item[{(iii)}] $\det T_{\bm{n};\bm{m}} \ne 0$ if and only if there is a unique vector $\bm{\Xi}_{\bm{n};\bm{m}} = (\Xi_{\bm{n};\bm{m},1},\dots,\Xi_{\bm{n};\bm{m},r})$ of Laurent polynomials of the form
    \begin{equation}\label{eq:moprlIspan}
    \Xi_{\bm{n};\bm{m},j}(z) \in \operatorname{span}\big\{z^k\big\}_{k = -n_j}^{m_j-1},
    \end{equation}
    that satisfies 
    \begin{subnumcases}{\sum_{j = 1}^r L_j[\Xi_{\bm{n};\bm{m},j}(w)w^{-k}] =}  
            \label{eq:moprlI}
            0, & %k  = -\abs{\bm{n}}+1,\dots,\abs{\bm{m}}-1,
            $-\abs{\bm{n}}+1\le k \le \abs{\bm{m}}-1$,
            \\
            \label{eq:moprlInorm}
            1, & $k  =-\abs{\bm{n}}$.
        %0, \quad k = -\abs{\bm{n}}+1,\dots,\abs{\bm{m}}-1.
    \end{subnumcases}
    % \begin{equation}\label{eq:moprlI}
    %     \sum_{j = 1}^r L_j[\Xi_{\bm{n};\bm{m},j}(w)w^{-k}] = 
    %     \begin{cases}
    %         0, & %k  = -\abs{\bm{n}}+1,\dots,\abs{\bm{m}}-1,
    %         -\abs{\bm{n}}+1\le k \le \abs{\bm{m}}-1,
    %         \\
    %         1, & k  =-\abs{\bm{n}}.
    %     \end{cases}
    %     %0, \quad k = -\abs{\bm{n}}+1,\dots,\abs{\bm{m}}-1.
    % \end{equation}
    \item[{(iv)}] $\det T_{\bm{n};\bm{m}} \ne 0$ if and only if there is a unique vector $\bm{\Xi}^*_{\bm{n};\bm{m}} = (\Xi^*_{\bm{n};\bm{m},1},\dots,\Xi^*_{\bm{n};\bm{m},r})$ of Laurent polynomials of the form
    \begin{equation}\label{eq:moprlI*span}
    \Xi^*_{\bm{n};\bm{m},j}(z) \in \operatorname{span}\big\{z^k\big\}_{k = -n_j+1}^{m_j},
    \end{equation}
    that satisfies 
    \begin{subnumcases}{\sum_{j = 1}^r L_j[\Xi^*_{\bm{n};\bm{m},j}(w)w^{-k}] =}         
            \label{eq:moprlI*}
            0, & %k  = -\abs{\bm{n}}+1,\dots,\abs{\bm{m}}-1,
            $-\abs{\bm{n}}+1\le k \le \abs{\bm{m}}-1$,
            \\
            \label{eq:moprlI*norm}
            1, & $k  =\abs{\bm{m}}$.
        %0, \quad k = -\abs{\bm{n}}+1,\dots,\abs{\bm{m}}-1.
    \end{subnumcases}
    % \begin{equation}\label{eq:moprlI*}
    %     \sum_{j = 1}^r L_j[\Xi^*_{\bm{n};\bm{m},j}(w)w^{-k}] = 
    %     \begin{cases}
    %         0, & %k  = -\abs{\bm{n}}+1,\dots,\abs{\bm{m}}-1,
    %         -\abs{\bm{n}}+1\le k \le \abs{\bm{m}}-1,
    %         \\
    %         1, & k  =\abs{\bm{m}}.
    %     \end{cases}
    %     %0, \quad k = -\abs{\bm{n}}+1,\dots,\abs{\bm{m}}-1.
    % \end{equation}
\end{itemize}
\end{prop}
\begin{proof}
    Treating the coefficients of $\Phi_{\bm{n},\bm{m}}$ in~\eqref{eq:moprlIIspan} as unknowns, write the system of equations \eqref{eq:moprlII}. Its $(\abs{\bm{n}}+\abs{\bm{m}}) \times (\abs{\bm{n}}+\abs{\bm{m}})$ coefficient matrix is $T_{\bm{n};\bm{m}}$ which proves (i). 
%\begin{equation}\label{eq:T}
%    T_{\bm{n};\bm{m}} = \begin{pmatrix}
%        L_1[w^{-\abs{\bm{m}}}w^{m_1}] & \cdots & L_1[w^{\abs{\bm{n}}-1}w^{m_1}] \\
%        \vdots & \ddots & \vdots \\
%        L_1[w^{-\abs{\bm{m}}}w^{-n_1+1}] & \cdots & L_1[w^{\abs{\bm{n}}-1}w^{-n_1+1}] \\
%        \hline
%        & \vdots & \\
%        \hline
%        L_r[w^{-\abs{\bm{m}}}w^{m_r}] & \cdots & L_r[w^{\abs{\bm{n}}-1}w^{m_r}] \\
%        \vdots & \ddots & \vdots \\
%        L_r[w^{-\abs{\bm{m}}}w^{-n_r+1}] & \cdots & L_r[w^{\abs{\bm{n}}-1}w^{-n_r+1}]
%    \end{pmatrix}, 
%\end{equation}

    Similar arguments work in other cases, leading to the same coefficient matrix $T_{\bm{n};\bm{m}}$ in (ii) and to the transpose of $T_{\bm{n};\bm{m}}$ in (iii)  and (iv).
\end{proof}
\begin{rem}
    Observe that in (iii) and (iv) we are using $\bm{n}$-indices with the negative sign and $\bm{m}$-indices with the positive sign, as opposed to the cases  (i) and (ii). This is to ensure duality, i.e., that uniqueness of $\Phi_{\bm{n};\bm{m}}$  is equivalent to the uniqueness of $\bm{\Xi}_{\bm{n};\bm{m}}$, as the proposition above shows. 
    %Observe that in {\it (iii)} and {\it (iv)} we are using $\bm{n}$-indices for negative powers and $\bm{m}$-indices for positive powers, as opposed to the cases {\it (i)} and {\it (ii)}. This is to ensure duality, i.e., that uniqueness of $\Phi_{\bm{n};\bm{m}}$  is equivalent to the uniqueness of $\bm{\Xi}_{\bm{n};\bm{m}}$, as the proposition above shows. 
\end{rem}
\begin{rem}\label{rem:r=1} 
In the case of orthogonality with respect to  one probability measure   (so that $r=1$, $\bm{n} = n \in \Z$, $\bm{m} = m \in \Z$, $n+m\ge0$), it is easy to see that 
\begin{alignat*}{2}
    \Phi_{n;m}(z) &= z^{-m}\Phi_{n+m}(z), \qquad  &\Phi_{n;m}^*(z) &= z^{-m}\Phi^*_{n+m}(z),
    \\
    {\Xi}_{n;m}(z) & = \tfrac{1}{\kappa_{n+m-1}} z^{-n}\Phi^*_{n+m-1}(z),
    \qquad 
    &{\Xi}^*_{n;m}(z) &= \tfrac{1}{\kappa_{n+m-1}} z^{-n+1} \Phi_{n+m-1}(z),
\end{alignat*}
% $\Phi_{n;m}(z) = z^{-m}\Phi_{n+m}(z)$, 
% $\Phi_{n;m}^*(z) = z^{-m}\Phi^*_{n+m}(z)$
% and ${\Xi}_{n;m}(z) = \tfrac{1}{\kappa_{n+m-1}} z^{m-1}(\Phi^*_{n+m-1}(z))^\sharp$, ${\Xi}^*_{n;m}(z) = \tfrac{1}{\kappa_{n+m-1}} z^{m}(\Phi_{n+m-1}(z))^\sharp$, where the ${}^\sharp$ operation is defined in~\eqref{eq:sharp1} below and 
where $\kappa_j = L[\Phi_j(w)w^{-j}] = L[|\Phi_j(w)|^2]$. 
%$\Phi_{n;m}^*(z) = z^n\overline{\Phi_{n+m}(1/\bar{z})}$
%In particular, we see  that $\Phi_n^*(z) = z^n\overline{\Phi_n(1/\bar{z})}$ (holds only if $r=1$ and $\mu$ being positive!).
\end{rem}

The Laurent polynomials described in   (i), (ii),  (iii), and (iv) will be called type II, type II$^*$, I, and I$^*$, respectively. For the case of indices of the form $(\bm{n};-\bm{n})$ (for any $\bm{n}\in\Z^r$), which are all normal, it is natural to take $\Phi_{\bm{n};-\bm{n}}(z) = \Phi^*_{\bm{n};-\bm{n}}(z) = 1$, $\bm{\Xi}_{\bm{n};-\bm{n}}(z) = \bm{\Xi}^*_{\bm{n};-\bm{n}}(z) = \bm{0}$. %$T_{\bm{0},\bm{0}} = 1$, 
%and treat $(\bm{0};\bm{0})$ as a normal multi-index.

If $(\bm{n};\bm{m})$ is normal, we will write $\alpha_{\bm{n};\bm{m}}$ for the $z^{-\abs{\bm{m}}}$-coefficient of $\Phi_{\bm{n};\bm{m}}(z)$ and $\beta_{\bm{n};\bm{m}}$ for the $z^{\abs{\bm{n}}}$-coefficient of $\Phi_{\bm{n};\bm{m}}^*(z)$. These should be viewed as the generalized Verblunsky % (or, more generally, Baxter) 
recurrence coefficients, see Section~\ref{ss:SzegoNEW}.
%, where $\Phi_{\bm{n};\bm{m}}$ and $\Phi_{\bm{n};\bm{m}}^*$ will be assumed to be the normalized polynomials \eqref{eq:moprlIIspan} and \eqref{eq:moprlII*span}. %In the case when each $L_j$ is given by integration with respect to a unit circle measure, we have the following relation between $\alpha$'s and $\beta$'s.


%; this is to ensure that the two problems are dual (i.e., normal simultaneously, see Proposition~\ref{prop:laurent normality type I}).

%\begin{prop}\label{prop:typeIMeasuresSymmetry}
%    Suppose $L_j[w^n] = \int w^n d\mu_j(w)$ for some probability measure supported on the unit circle, for $n \in \Z$ and $j = 1,\dots,r$. Then we have $\bm{\Xi}_{\bm{n};\bm{m}}^\sharp = \bm{\Xi}_{\bm{m};\bm{n}}^*$.
%\end{prop}

%In the case when each $L_j$ is given by integration with respect to a unit circle measure, there is more we can say about these Laurent polynomials.

Let us introduce the following notation. Given a function $f(z)$ of a complex variable, let
\begin{equation}\label{eq:sharp1}
        f^\sharp(z) = \overline{f(1/\bar{z})}.
\end{equation}
Similarly, given a moment functional $L$ with $L[w^{-k}] = c_k$, let
\begin{equation}\label{eq:sharp2}
        L^\sharp[w^{-k}] = \overline{L[w^k]} = \bar{c}_{-k}, \qquad k\in\Z.
\end{equation}
Equivalently,
\begin{equation}\label{eq:sharp3}
    L^\sharp[f(z)] = \overline{L[f^\sharp(z)]}.
\end{equation}


\begin{prop}\label{thm:reversal vs star polynomials}
    $(\bm{n};\bm{m})\in \mathfrak{C}_{2r}$ %, $\bm{n}+\bm{m}\ge \bm{0}$, 
    is normal with respect to a system $\bm{L}=(L_1,\ldots,L_r)$ if and only if $(\bm{m};\bm{n})$ is normal with respect to $\bm{L}^\sharp =(L_1^\sharp,\ldots,L_r^\sharp)$. In this case,
\begin{align}
\label{eq:symm1}
    &\Phi_{\bm{n};\bm{m}}(z,\bm{L}^\sharp)  = \Phi^*_{\bm{m};\bm{n}}(z,\bm{L}) ^\sharp,
    \\
    \label{eq:symm2}
    &\Phi^*_{\bm{n};\bm{m}}(z,\bm{L}^\sharp) = \Phi_{\bm{m};\bm{n}}(z,\bm{L})^\sharp,
    \\
    \label{eq:symm3}
    &\bm\Xi_{\bm{n};\bm{m}}(z,\bm{L}^\sharp)  = \bm\Xi^*_{\bm{m};\bm{n}}(z,\bm{L}) ^\sharp,
    \\
    \label{eq:symm4}
    &\bm\Xi^*_{\bm{n};\bm{m}}(z,\bm{L}^\sharp) = \bm\Xi_{\bm{m};\bm{n}}(z,\bm{L})^\sharp,
    \\
    \label{eq:symm5}
    &\alpha_{\bm{n};\bm{m}}(\bm{L}^\sharp) = \bar{\beta}_{\bm{m};\bm{n}}(\bm{L}),
    \\
    \label{eq:symm6}
    &\beta_{\bm{n};\bm{m}}(\bm{L}^\sharp) = \bar{\alpha}_{\bm{m};\bm{n}}(\bm{L}).
    \end{align}

    Furthermore, if each $L_j$ satisfies $c_{k,j} = \bar{c}_{-k,j}$ for all $k\in\N$ (e.g., if it is induced by a probability measure \eqref{eq:functionals given by measures} on the unit circle $\T$), then $(\bm{n};\bm{m})$ is normal if and only if $(\bm{m};\bm{n})$ is normal, and 
    \begin{equation}\label{eq:symm7}
        \Phi_{\bm{n};\bm{m}}^\sharp = \Phi_{\bm{m};\bm{n}}^*,
        \qquad
        \bm{\Xi}_{\bm{n};\bm{m}}^\sharp = \bm{\Xi}_{\bm{m};\bm{n}}^*, 
        \qquad 
        \alpha_{\bm{n};\bm{m}} = \bar{\beta}_{\bm{m};\bm{n}}
    \end{equation}
    hold.
\end{prop}
% \begin{prop}\label{thm:reversal vs star polynomials}
%     Suppose each of the moment functionals $L_j$ is induced by a probability measure 
%     % \begin{equation}\label{eq:functionals given by measures}
%     %     L_j[w^k] = \int_{\T} w^k d\mu_j(w), \qquad j = 1,\dots,r, \qquad k \in \Z.
%     % \end{equation} 
%     \eqref{eq:functionals given by measures}
%     %for probability measures $\mu_j$ supported 
%     on the unit circle $\T$. Then $(\bm{n};\bm{m})$ is normal if and only if $(\bm{m};\bm{n})$ is normal. In this situation we have
%     \begin{equation}
%         \bm{\Xi}_{\bm{n};\bm{m}}^\sharp = \bm{\Xi}_{\bm{m};\bm{n}}^*, \quad \Phi_{\bm{n};\bm{m}}^\sharp = \Phi_{\bm{m};\bm{n}}^*, \quad \alpha_{\bm{n};\bm{m}} = \bar{\beta}_{\bm{m};\bm{n}}.
%     \end{equation}
%     %where $f^\sharp$ denotes the function
% \end{prop}
% \begin{proof}
%     This follows immediately from the equality
%     \begin{equation}\label{eq:TSymmetry}
%         \int_{\T} \overline{f(1/\bar{w})} w^{-k} d\mu(w) = \overline{\int_\T f(w) w^k d\mu(w)}, \qquad k\in\Z.
%     \end{equation}
% \end{proof}
\begin{proof}
    It is easy to see from~\eqref{eq:T} that $\det T_{\bm{n};\bm{m}}(\bm{L}) = \pm \overline{\det T_{\bm{m};\bm{n}}(\bm{L}^\sharp)}$, which implies the normality claim.  
    %Equalities~\eqref{eq:symm1}--\eqref{eq:symm4} follow from the definitions,~\eqref{eq:moprlIIspan}--\eqref{eq:moprlI} and~\eqref{eq:sharp1}--\eqref{eq:sharp2}. 
    Equalities~\eqref{eq:symm1}--\eqref{eq:symm4} follow from the definitions and~\eqref{eq:sharp3}.
    Equalities~\eqref{eq:symm1}--\eqref{eq:symm2} then imply~\eqref{eq:symm5}--\eqref{eq:symm6}. 

    In the special case when the moments $c_{k,j}$ in~\eqref{eq:moments linear functional} satisfy $c_{k,j} = \bar{c}_{-k,j}$ (in particular this holds if $L_j$ corresponds to a positive measure on $\T$), then $L_j = L_j^\sharp$. Combining this with~\eqref{eq:symm1},~\eqref{eq:symm3}, and~\eqref{eq:symm5} gives \eqref{eq:symm7}.
\end{proof}
\begin{rem}\label{rem:annoying1}
    The type II and I polynomials  $\Phi_{\bm{n}}$, $\Phi_{\bm{n}}^*$, $\bm{\Lambda}_{\bm{n}}$, $\bm{\Lambda}_{\bm{n}}^*$ from~\cite{KVMOPUC}  are  related to our Laurent polynomials via
\begin{alignat}{3}
    \label{eq:mopucII}
    &\Phi_{\bm{n}}(z) = \Phi_{\bm{n};\bm{0}}(z),\qquad &&  \Phi_{\bm{n}}^*(z) = \Phi^*_{\bm{n};\bm{0}}(z), \\ 
    \label{eq:mopucI}
    &\bm{\Lambda}_{\bm{n}}(z) = z^{-1} \bm{\Xi}^*_{\bm{0};\bm{n}}(z), \qquad &&  \bm{\Lambda}_{\bm{n}}^*(z) = \bm{\Xi}_{\bm{0};\bm{n}}(z).
\end{alignat}
    The same polynomials $\Phi_{\bm{n}}$, $\bm{\Lambda}_{\bm{n}}$  were also studied earlier in~\cite{MOPUC1,MOPUC2}.
\end{rem}
\begin{rem}\label{rem:annoying2}
    While the relations~\eqref{eq:mopucII} are natural, those in~\eqref{eq:mopucI} may at first appear somewhat unexpected. This can be explained by two factors. First, the polynomials $\bm\Lambda_{\bm{n}}$ analyzed in~\cite{MOPUC1,MOPUC2,KVMOPUC} for systems of positive measures are not the most natural objects in the setting of general linear functionals $L_j$ since their uniqueness  (normality) is {\it not} equivalent to that of $\Phi_{\bm{n}}$, but instead to that of $\Phi_{\bm{n}}$ corresponding to the functionals $L_j^\sharp$ (see~\eqref{eq:sharp2}).
    %The ``correct'' choice should have been either $\bm{\Xi}_{\bm{n};\bm{0}}(z)$ or $z\bm{\Xi}_{\bm{n};\bm{0}}(z)$. 
    This explains the reversed order of indices $(\bm{0};\bm{n})$ in ~\eqref{eq:mopucI}. Second, the additional factor $z^{-1}$ is, in fact, natural, as it leads to a more aesthetically consistent formulation of the Hermite--Pad\'{e} problem~\eqref{eq:HPtypeI1swapped}--\eqref{eq:HPtypeI2swapped}, the type I recurrence relations~\eqref{eq:SzegoI1}--\eqref{eq:SzegoI2} (compare with~\eqref{eq:SzegoII1}--\eqref{eq:SzegoII2}), and the Christoffel--Darboux formula~\eqref{eq:cd formula1}--\eqref{eq:cd formula2} (compare with~\cite[Thm 6.1]{KVMOPUC}). 
\end{rem}


\section{A General Two-Point Hermite--Pad\'e Problem of Type II}\label{ss:general type II hermite pade}

%In this section let us discuss a different type of two-point Hermite--Pad\'{e} problem. Here we allow the degree of approximations at $0$ and $\infty$ to be arbitrary non-negative integers. 

In this section we discuss the Hermite--Pad\'{e} approximation problem associated with the Laurent polynomials defined in the previous section.  

To each linear functional $L_j$, $j=1,\dots,r$, given in~\eqref{eq:moments linear functional}, we assign the formal power series
\begin{equation}
\label{eq:FGeneral}
    F_j^{(0)}(z) = c_{0,j} + 2\sum_{k = 1}^{\infty}c_{k,j}z^k, \quad F_j^{(\infty)}(z) = -c_{0,j} - 2\sum_{k = 1}^{\infty}c_{-k,j}z^{-k}.
\end{equation}
Since
\begin{equation}\label{eq:kernel}
        \frac{w+z}{w-z}
        =
        \begin{cases}
            1+ 2\sum_{k=1}^\infty \frac{z^k}{w^k},\qquad  & \mbox{if } |z|<|w|, \\
            -1 - 2\sum_{k=1}^\infty \frac{w^k}{z^k}, \qquad & \mbox{if } |z|>|w|, \\
        \end{cases}
\end{equation} 
the pair $(F_j^{(0)},F_j^{(\infty)})$ can formally be viewed as the generalized Carath\'{e}odory function
\begin{equation}\label{eq:generalized caratheodory function}
    F_j(z) = L_j[(w+z)(w-z)^{-1}].
\end{equation} 
In the case when $L_j$ is associated with a probability measure $\mu_j$ as in~\eqref{eq:functionals given by measures}, then~\eqref{eq:FGeneral} are the power series expansions of the Carath\'{e}odory function $F_j(z)$ \eqref{eq:F} of $\mu_j$ at $0$ and at $\infty$. 

Let two multi-indices $\bm{n} = (n_1,\dots,n_r)\in\N^r$, $\bm{m} = (m_1,\dots,m_r)\in\N^r$ %, with $\bm{n}+\bm{m}\ge\bm{0}$,
be given%, not simultaneously zero
. Consider the generalized two-point Hermite--Pad\'e problem of finding the Laurent polynomial
\begin{equation}\label{eq:spanFull}
    \Phi_{\bm{n};\bm{m}}
    %,\Psi_{\bm{n};\bm{m},1},\ldots,\Psi_{\bm{n};\bm{m},r}
    \in \operatorname{span}\big\{z^k\big\}_{k = -\abs{\bm{m}}}^{\abs{\bm{n}}},
\end{equation}
that simultaneously satisfies
\begin{alignat}{3}
\label{eq:generalized two point hermite pade eq 1}
    & \Phi_{\bm{n};\bm{m}}(z)F_j^{(0)}(z) + \Psi_{\bm{n};\bm{m},j}(z) = \mathcal{O}(z^{n_j}),  \qquad &&z \rightarrow 0, \\\label{eq:generalized two point hermite pade eq 2}
    & \Phi_{\bm{n};\bm{m}}(z)F_j^{(\infty)}(z) + \Psi_{\bm{n};\bm{m},j}(z) = \mathcal{O}(z^{-m_j-1}), \qquad &&z \rightarrow \infty,
\end{alignat}
for some Laurent polynomials $\Psi_{\bm{n};\bm{m},1},\ldots,\Psi_{\bm{n};\bm{m},r} \in \operatorname{span}\big\{z^k\big\}_{k = -\abs{\bm{m}}}^{\abs{\bm{n}}}$.
%simultaneously for each $j=1,\ldots,r$. 

Note that here we take $\bm{n}\ge\bm{0}$ and $\bm{m}\ge \bm{0}$, which is more restrictive than in the previous section. This is related to the fact that $F^{(0)}$ and $F^{(\infty)}$ start from the $z^0$-term, and therefore the problem~\eqref{eq:spanFull}, \eqref{eq:generalized two point hermite pade eq 1}, \eqref{eq:generalized two point hermite pade eq 2} can become ill-posed if $n_j<0$ or $m_j<0$. 

%{We can treat ~\eqref{eq:generalized two point hermite pade eq 1}--\eqref{eq:generalized two point hermite pade eq 2} purely algebraically as formal power series, without worrying about the convergence.}
 
%Here we think of $(\bm{n};\bm{m})$ as one whole multi-index $(n_1,\dots,n_r,m_1,\dots,m_r)$, and 
%Here we use the index $\bm{n}$ to keep track of the indices in \eqref{eq:generalized two point hermite pade eq 1}, and $\bm{m}$ to keep track of the indices in \eqref{eq:generalized two point hermite pade eq 2}. 
%{Note that the Hermite-Pad\'e problem associated with $\varphi_{\bm{n}}(z)$ considered in Section~\ref{ss:HPphiEven} with all $n_j$ even, coincides with the current problem ~\eqref{eq:spanFull}, \eqref{eq:generalized two point hermite pade eq 1}, \eqref{eq:generalized two point hermite pade eq 2} associated with $\varphi_{\bm{n}/2;\bm{n}/2}(z)$. However the the Hermite-Pad\'e problems in Sections~\ref{ss:HPphiOdd1}, ~\ref{ss:HPphiOdd2}, and~\ref{ss:HPphiFull} are {\it not} special cases of~\eqref{eq:spanFull}, \eqref{eq:generalized two point hermite pade eq 1}, \eqref{eq:generalized two point hermite pade eq 2}.}

%Note that \eqref{eq:two point hermite pade eq 1}-\eqref{eq:two point hermite pade eq 2} is just \eqref{eq:generalized two point hermite pade eq 1}-\eqref{eq:generalized two point hermite pade eq 2} when $\bm{m} = \bm{n}$, and \eqref{eq:two point hermite pade 1 old}-\eqref{eq:two point hermite pade 2 old} is \eqref{eq:generalized two point hermite pade eq 1}-\eqref{eq:generalized two point hermite pade eq 2} when $\bm{m} = \bm{0}$. In other words, the Laurent polynomial $\Phi_{\bm{n};\bm{m}}(z) = z^{-\abs{\bm{m}}}Q_{\bm{n};\bm{m}}(z)$ satisfies $\Phi_{\bm{n};\bm{n}} = \Psi_{2\bm{n}}$ and $\Phi_{\bm{n};\bm{0}} = \Phi_{\bm{n}}$. Theorem \ref{thm:hermite pade orthogonality} also generalizes to \eqref{eq:generalized two point hermite pade eq 1}-\eqref{eq:generalized two point hermite pade eq 2} through the following result.

\begin{thm}\label{thm:generalized hermite pade orthogonality}
    Any Laurent polynomial $\Phi_{\bm{n};\bm{m}}$ that solves ~\eqref{eq:spanFull}, \eqref{eq:generalized two point hermite pade eq 1}, \eqref{eq:generalized two point hermite pade eq 2} 
    %is a type II multiple orthogonal Laurent polynomials with respect to $\bm{L}$ satisfying 
    automatically satisfies the orthogonality conditions~\eqref{eq:moprlII},
    %\begin{equation}\label{eq:moprlII}
    %    L_j[\varphi_{\bm{n};\bm{m}}(w)w^{-k}] = 0, \qquad k = -m_j,\dots,n_j-1.
    %\end{equation}
    %for all $j=1,\ldots,r$,
    and $\Psi_{\bm{n};\bm{m},j}$ is given by
    \begin{equation}\label{eq:psiGeneral}
        \Psi_{\bm{n};\bm{m},j}(z) = L_j\Big[\frac{w+z}{w-z}\big(\Phi_{\bm{n};\bm{m}}(w) - \Phi_{\bm{n};\bm{m}}(z)\big)\Big] + L_j[\Phi_{\bm{n};\bm{m}}(w)].
   \end{equation}
   Conversely, any Laurent polynomial $\Phi_{\bm{n},\bm{m}} \in \operatorname{span}\set{z^k}_{k = -\abs{\bm{m}}}^{\abs{\bm{n}}}$ that satisfies \eqref{eq:moprlII} solves the Hermite--Pad\'e problem ~\eqref{eq:spanFull}, \eqref{eq:generalized two point hermite pade eq 1}, \eqref{eq:generalized two point hermite pade eq 2}  together with the Laurent polynomial \eqref{eq:psiGeneral}.
\end{thm}
\begin{rem}
    Combining this theorem with the results in~\cite{KVMLOPUC1,KNik}, we can see that up to a normalization constant, the solution to the Hermite--Pad\'{e} problem ~\eqref{eq:spanFull}, \eqref{eq:generalized two point hermite pade eq 1}, \eqref{eq:generalized two point hermite pade eq 2}  with $\bm{m}=\bm{n}$ is unique if $\bm{L}$ is an Angelesco, AT, or Nikishin system (with $r=2$) on the unit circle.
\end{rem}
\begin{proof}
    %The proof follow similar lines to the original result of Peherstorfer--Steinbauer~\cite{PeherstorferSteinbauer}. 
    Let us write $\Phi_{\bm{n};\bm{m}}$ and $\Psi_{\bm{n};\bm{m},j}$ in the form
    \begin{align}\label{eq:coefficients of phi and psi}
        & \Phi_{\bm{n};\bm{m}}(z) =
        \kappa_{-\abs{\bm{m}}}z^{-\abs{\bm{m}}}+\dots + \kappa_{\abs{\bm{n}}}z^{\abs{\bm{n}}} , 
        %\sum_{k=-\abs{\bm{m}}}^{\abs{\bm{n}}} \kappa_{k}z^{k},
        \\ 
        \label{eq:coefficients of y}
        & \Psi_{\bm{n};\bm{m},j}(z) =
        \lambda_{-\abs{\bm{m}},j}z^{-\abs{\bm{m}}} + \dots +\lambda_{\abs{\bm{n}},j}z^{\abs{\bm{n}}}.
        %\sum_{k=-\infty}^\infty \lambda_{k,j}z^{k},
    \end{align}
    where there are finitely many non-zero $\lambda_{k,j}$ coefficients, so that~\eqref{eq:coefficients of y} is in fact a Laurent polynomial.
    Then we get
    \begin{align*}
        & \Phi_{\bm{n};\bm{m}}(z)F_j^{(0)}(z) = a_{-\abs{\bm{m}},j}z^{-\abs{\bm{m}}} + \dots + a_{n_j-1,j}z^{n_j-1} + \mathcal{O}(z^{n_j}), \quad z \rightarrow 0, \\
        & \Phi_{\bm{n};\bm{m}}(z)F_j^{(\infty)}(z) = b_{\abs{\bm{n}},j}z^{\abs{\bm{n}}} + \dots + b_{-m_j,j}z^{-m_j} + \mathcal{O}(z^{-m_j-1}), \quad z \rightarrow \infty,
    \end{align*}
    where $a_{-\abs{\bm{m}},j},\dots,a_{n_j-1,j}$ and $b_{-m_j,j},\dots,b_{\abs{\bm{n}}-1,j}$ are given by
    \begin{align}
        \label{eq:coefficients of series type II}
         & a_{k,j} = \kappa_kc_{0,j} + 2\kappa_{k-1}c_{1,j} + \dots + 2\kappa_{-\abs{\bm{m}}} c_{\abs{\bm{m}}+k,j},  %&k& = -\abs{\bm{m}},\dots,n_j-1, 
         \\ 
         \label{eq:coefficients of series 2}
         & b_{k,j} = -\kappa_kc_{0,j} - 2\kappa_{k+1} c_{-1,j} - \dots - 2 \kappa_{\abs{\bm{n}}} c_{-\abs{\bm{n}}+k,j}.%, &k& = -m_j,\dots,\abs{\bm{n}}.
    \end{align}
    For \eqref{eq:generalized two point hermite pade eq 1}--\eqref{eq:generalized two point hermite pade eq 2} to hold we must have
    \begin{align}\label{eq:coefficients type II eq 1}
        \lambda_{k,j} = a_{k,j}, \qquad &k = -\abs{\bm{m}},\dots,n_j-1, \\\label{eq:coefficients type II eq 2}
        \lambda_{k,j} = b_{k,j}, \qquad &k = -m_j,\dots,\abs{\bm{n}}.
    \end{align}
    This implies 
    \begin{equation}
        a_{k,j} = b_{k,j}, \qquad k = -m_j,\dots,n_j-1,
    \end{equation}
    and by \eqref{eq:moments linear functional} and \eqref{eq:coefficients of series type II}--\eqref{eq:coefficients of series 2}, this turns into \eqref{eq:moprlII}. 
    
    From \eqref{eq:coefficients type II eq 1}--\eqref{eq:coefficients type II eq 2} we also see that $\Psi_{\bm{n};\bm{m},j}$ is uniquely determined by $\Phi_{\bm{n};\bm{m}}$. To compute $\Psi_{\bm{n};\bm{m},j}$, write $\widetilde{\Psi}_{\bm{n};\bm{m},j}$ for the right hand side of \eqref{eq:psiGeneral}, and define $R_{\bm{n};\bm{m},j}$ by
    \begin{equation}
        R_{\bm{n};\bm{m},j}(z) = L_j\Big[\frac{w+z}{w-z}\Phi_{\bm{n};\bm{m}}(w) \Big].
    \end{equation}
    Again, this should be treated as the two formal power series
    \begin{align}\label{eq:secondKindFunctionTypeII 1}
    R^{(0)}_{\bm{n};\bm{m},j}(z) 
        & =  L_j[\Phi_{\bm{n};\bm{m}}(w)]+ 2\sum_{k=1}^\infty z^k L_j[w^{-k}\Phi_{\bm{n};\bm{m}}(w)],  & z & \to 0,  
        %\mbox{if } & |z|<|w|, 
        \\ \label{eq:secondKindFunctionTypeII 2}
    R^{(\infty)}_{\bm{n};\bm{m},j}(z) & =  -L_j[\Phi_{\bm{n};\bm{m}}(w)] - 2\sum_{k=1}^\infty z^{-k} L_j[w^{k}\Phi_{\bm{n};\bm{m}}(w)], & z & \to \infty. 
    % & \mbox{if } |z|>|w|. \\
    \end{align}
    By \eqref{eq:moprlII}
    we get
    \begin{align}\label{eq:R asymptotics type II eq 1}
        R^{(0)}_{\bm{n};\bm{m},j}(z) 
        & =  \mathcal{O}(z^{n_j}),  & z & \to 0, \\ \label{eq:R asymptotics type II eq 2}
        R^{(\infty)}_{\bm{n};\bm{m},j}(z) & = -L_j[\Phi_{\bm{n};\bm{m}}(w)] + \mathcal{O}(z^{-m_j-1}),  & z & \to \infty.
    \end{align}
    The constant term in~\eqref{eq:R asymptotics type II eq 2} is only relevant if $n_j=0$ and otherwise can be removed by orthogonality~\eqref{eq:moprlII}. 
    
    Note that $\frac{w+z}{w-z}(\Phi_{\bm{n};\bm{m}}(w) - \Phi_{\bm{n};\bm{m}}(z))$ is a Laurent polynomial in $z$ (as well as in $w$). By~\eqref{eq:kernel},
    % \begin{equation}\label{eq:kernel}
    %     \frac{w+z}{w-z}
    %     =
    %     \begin{cases}
    %         1+ 2\sum_{k=1}^\infty \frac{z^k}{w^k},\qquad  & \mbox{if } |z|<|w| \\
    %         -1 - 2\sum_{k=1}^\infty \frac{w^k}{z^k}, \qquad & \mbox{if } |z|>|w| \\
    %     \end{cases}
    % \end{equation} 
    this polynomial coincides with the $z$-power series at $0$
    % $$
    % \Big(1+ 2\sum_{k=1}^\infty \frac{z^k}{w^k} \Big) \Big(\varphi_{\bm{n};\bm{m}}(w) - \varphi_{\bm{n};\bm{m}}(z)\Big)
    % $$
    \begin{equation}
        \Big(1+ 2\sum_{k=1}^\infty \frac{z^k}{w^k} \Big) \Phi_{\bm{n};\bm{m}}(w) -  \Big(1+ 2\sum_{k=1}^\infty \frac{z^k}{w^k} \Big)  \Phi_{\bm{n};\bm{m}}(z),
    \end{equation}
    as well as with the $z$-power series at $\infty$
    % $$
    % \Big(-1- 2\sum_{k=1}^\infty \frac{w^k}{z^k} \Big) \Big(\varphi_{\bm{n};\bm{m}}(w) - \varphi_{\bm{n};\bm{m}}(z)\Big)
    % $$
    \begin{equation}
       \Big(-1- 2\sum_{k=1}^\infty \frac{w^k}{z^k} \Big) \Phi_{\bm{n};\bm{m}}(w) - \Big(-1- 2\sum_{k=1}^\infty \frac{w^k}{z^k} \Big)  \Phi_{\bm{n};\bm{m}}(z). 
    \end{equation}
    Therefore we can apply the $L_j$ functional to obtain
    \begin{align}\label{eq:psiVsRGeneral}
        \widetilde{\Psi}_{\bm{n};\bm{m},j}(z) & = {R}^{(0)}_{\bm{n};\bm{m},j}(z)  -  \Phi_{\bm{n};\bm{m}}(z) F^{(0)}_j(z)  + L_j[\Phi_{\bm{n},\bm{m}}(w)], 
        \\
        \widetilde{\Psi}_{\bm{n};\bm{m},j}(z) & = {R}^{(\infty)}_{\bm{n};\bm{m},j}(z)  -  \Phi_{\bm{n};\bm{m}}(z) F^{(\infty)}_j(z) + L_j[\Phi_{\bm{n},\bm{m}}(w)].
    \end{align}
    By \eqref{eq:R asymptotics type II eq 1}-\eqref{eq:R asymptotics type II eq 2} we get that $\widetilde{\Psi}_{\bm{n};\bm{m},j}$ satisfies equalities~\eqref{eq:generalized two point hermite pade eq 1}--\eqref{eq:generalized two point hermite pade eq 2} which implies
    $\widetilde{\Psi}_{\bm{n};\bm{m},j} = \Psi_{\bm{n};\bm{m},j}$. 

    Conversely, given $\Phi_{\bm{n};\bm{m}}$ and the corresponding $\Psi_{\bm{n};\bm{m},j}$ as in~\eqref{eq:psiGeneral}, define $R_{\bm{n};\bm{m},j}$ as in~\eqref{eq:secondKindFunctionTypeII 1}--\eqref{eq:secondKindFunctionTypeII 2}. As above, we show that it satisfies~\eqref{eq:R asymptotics type II eq 1} and ~\eqref{eq:R asymptotics type II eq 2}, which then becomes ~\eqref{eq:generalized two point hermite pade eq 1}, \eqref{eq:generalized two point hermite pade eq 2}.
\end{proof}

Let us  also introduce the related Hermite--Pad\'e problem
\begin{alignat}{3}
\label{eq:generalized two point hermite pade star eq 1}
    & \Phi_{\bm{n};\bm{m}}^*(z)F_j^{(0)}(z) -\Psi_{\bm{n};\bm{m},j}^*(z) = \mathcal{O}(z^{n_j+1}), \qquad &&z \rightarrow 0, \\\label{eq:generalized two point hermite pade star eq 2}
    & \Phi_{\bm{n};\bm{m}}^*(z)F_j^{(\infty)}(z) - \Psi_{\bm{n};\bm{m},j}^*(z) = \mathcal{O}(z^{-m_j}), \qquad &&z \rightarrow \infty,
\end{alignat}
for Laurent polynomials 
\begin{equation}\label{eq:spanFull*}
 \Phi_{\bm{n};\bm{m}}^*, \Psi_{\bm{n};\bm{m},1}^*, \ldots, \Psi_{\bm{n};\bm{m},r}^* \in \operatorname{span}\big\{z^k\big\}_{k = -\abs{\bm{m}}}^{\abs{\bm{n}}}.
\end{equation}
\begin{thm}
    Any Laurent polynomial $\Phi_{\bm{n};\bm{m}}^*$ that solves~\eqref{eq:generalized two point hermite pade star eq 1}, \eqref{eq:generalized two point hermite pade star eq 2}, \eqref{eq:spanFull*} must satisfy the orthogonality relations \eqref{eq:moprlII*}, and $\Psi_{\bm{n};\bm{m},j}^*$ is given by
    \begin{equation}\label{eq:psiGeneral*}
        \Psi_{\bm{n};\bm{m},j}^*(z) = L_j\Big[\frac{w+z}{w-z}\big(\Phi_{\bm{n};\bm{m}}^*(z) - \Phi_{\bm{n};\bm{m}}^*(w) \big)\Big] + L_j[\Phi_{\bm{n};\bm{m}}^*(w)].
    \end{equation}
     Conversely, any Laurent polynomial $\Phi^*_{\bm{n},\bm{m}} \in \operatorname{span}\set{z^k}_{k = -\abs{\bm{m}}}^{\abs{\bm{n}}}$ that satisfies \eqref{eq:moprlII*} solves the Hermite--Pad\'e problem ~\eqref{eq:generalized two point hermite pade star eq 1}, \eqref{eq:generalized two point hermite pade star eq 2}, \eqref{eq:spanFull*}   together with the Laurent polynomial \eqref{eq:psiGeneral*}.
\end{thm}

\begin{proof}
   Use the same methods as in the proof of Theorem~\ref{thm:generalized hermite pade orthogonality}.
    The minus sign in~\eqref{eq:psiGeneral*} arises from the fact that, when expanding~\eqref{eq:secondKindFunctionTypeII 1} and~\eqref{eq:secondKindFunctionTypeII 2} for $\Phi_{\bm{n};\bm{m}}^*$, the additional term $L_j[\Phi_{\bm{n};\bm{m}}^*(w)]$ vanishes for $R^{(\infty)}$ but remains for $R^{(0)}$.
\end{proof}

% Note that the coefficient matrix of this system is the same as the coefficient matrix of the system of equations \eqref{eq:moprlII*} for Laurent polynomials on the form 
% \begin{equation}
%     \varphi_{\bm{n};\bm{m}}^*(z) = \beta_{\bm{n};\bm{m}}z^{\abs{\bm{n}}} + \dots + z^{-\abs{\bm{m}}}.
% \end{equation} 
% Hence normality in \eqref{eq:moprlII*} is equivalent to normality in \eqref{eq:moprlII}. 


\section{A General Two-Point Hermite-Pad\'e Problem of Type I}\label{ss:HPI}

Let again $\bm{n}\in\N^r$ and $\bm{m}\in\N^r$ be given. %, not simultaneously zero
We now consider the two-point Hermite--Pad\'e problem which can be viewed as dual to the one in the previous section~\eqref{eq:spanFull}, \eqref{eq:generalized two point hermite pade eq 1}, \eqref{eq:generalized two point hermite pade eq 2}. 

Let $F_j^{(0)}$, $F_j^{(\infty)}$ be the formal series~\eqref{eq:FGeneral} associated with linear functionals $L_j$ in~\eqref{eq:moments linear functional}. As before, if $L_j$ corresponds to integration with respect to a probability measure as in~\eqref{eq:functionals given by measures}, then $F_j^{(0)}$ and $F_j^{(\infty)}$ are the power series expansions of the Carath\'{e}odory function $F_j(z)$ \eqref{eq:F} of $\mu_j$ at $0$ and at $\infty$.
We are interested in the solutions to the following approximation problem:
\begin{alignat}{3}
\label{eq:HPtypeI1swapped}
    & \sum_{j = 1}^r \Xi_{\bm{n};\bm{m},j}(z)F^{(0)}_j(z) + \Upsilon_{\bm{n};\bm{m}}(z) = \mathcal{O}(z^{\abs{\bm{m}}}), \qquad && z \rightarrow 0, \\
\label{eq:HPtypeI2swapped}
    & \sum_{j = 1}^r \Xi_{\bm{n};\bm{m},j}(z)F_j^{(\infty)}(z) + \Upsilon_{\bm{n};\bm{m}}(z) = \mathcal{O}(z^{-\abs{\bm{n}}}), \qquad && z \rightarrow \infty,
\end{alignat}
for Laurent polynomials 
\begin{equation}
    \label{eq:HPtypeIspanSwapped}
    \Xi_{\bm{n};\bm{m},j}(z) \in \operatorname{span}\big\{z^{k}\big\}_{k=-n_j}^{m_j-1},
    \qquad
     \Upsilon_{\bm{n};\bm{m}}(z) \in \operatorname{span}\big\{z^{k}\big\}_{k=-N}^{M-1},
\end{equation}
%$\Xi_{\bm{n};\bm{m},j}^*$ will be in $\operatorname{span}\set{z^k}_{k=-m_j}^{n_j-1}$, solving the Hermite-Pad\'e problem
where $N:=\max\{n_1,\ldots,n_r\}$, $M:=\max\{m_1,\ldots,m_r\}$.


%for some Laurent polynomial $\Upsilon_{\bm{n};\bm{m}}^*$. 


\begin{thm}
    If the Laurent polynomials %vector 
    $\bm{\Xi}_{\bm{n};\bm{m}} = (\Xi_{\bm{n};\bm{m},1},\dots,\Xi_{\bm{n};\bm{m},r})$ and $\Upsilon_{\bm{n};\bm{m}}$ solve~\eqref{eq:HPtypeI1swapped}, \eqref{eq:HPtypeI2swapped}, \eqref{eq:HPtypeIspanSwapped},  then $\bm{\Xi}_{\bm{n};\bm{m}}$ satisfies the orthogonality relations \eqref{eq:moprlI}, and $\Upsilon_{\bm{n};\bm{m}}$ is given by
    \begin{equation}\label{eq:upsilon}
        \Upsilon_{\bm{n};\bm{m}}(z) =  \sum_{j = 1}^r L_j\Big[\frac{w+z}{w-z}\big(\Xi_{\bm{n};\bm{m},j}(w)-\Xi_{\bm{n};\bm{m},j}(z)\big)\Big] + {\delta_{\bm{n};\bm{m}} \sum_{j = 1}^r L_j[\Xi_{\bm{n};\bm{m},j}(w)]},
    \end{equation}
   {where the constant $\delta_{\bm{n};\bm{m}}$ is given by
   \begin{equation}\label{eq:delta}
       \delta_{\bm{n};\bm{m}} = \begin{cases}
           1, \qquad & \bm{m}=\bm{0} \mbox{ and } \bm{n}\neq\bm{0}, \\
           -1, \qquad & \bm{n}=\bm{0} \mbox{ and } \bm{m}\neq\bm{0}, \\
           0, \qquad &  \mbox{otherwise}.
       \end{cases}
   \end{equation}}

    Conversely, any vector $\bm{\Xi}_{\bm{n};\bm{m}} = (\Xi_{\bm{n};\bm{m},1},\dots,\Xi_{\bm{n};\bm{m},r})$ with $\Xi_{\bm{n};\bm{m},j} \in \operatorname{span}\set{z^k}_{k = -n_j}^{m_j-1}$, that satisfies %the orthogonality relations
    ~\eqref{eq:moprlI}, solves the Hermite--Pad\'{e} problem~\eqref{eq:HPtypeI1swapped}, \eqref{eq:HPtypeI2swapped}, \eqref{eq:HPtypeIspanSwapped} together with the Laurent polynomial~\eqref{eq:upsilon}.
\end{thm}


%%%%%%%%%%%%%%%%%%%%%%%%%%%%%%%%%%%%%%%%%%%%%%%%%%%%%%%%%%%%%%


\begin{proof}
    Write $\Xi_{\bm{n};\bm{m},j}$ and $\Upsilon_{\bm{n};\bm{m}}$ in the form 
    \begin{align*}\label{eq:lamdba}
         \Xi_{\bm{n};\bm{m},j} (z)&=  \kappa_{-n_j,j}z^{-n_j} + \kappa_{-n_j+1,j}z^{-n_j+1} + \ldots + \kappa_{m_j-1,j}z^{m_j-1}, \\ 
        \Upsilon_{\bm{n};\bm{m}}(z) & =\lambda_{-N}z^{-N} + \lambda_{-N+1}z^{-N+1} + \ldots + \lambda_{M-1}z^{M-1} .
    \end{align*}
    We put $\kappa_{k,j} = 0$ if $k \ge m_j$ or $k\le -n_j-1$, and %it will later become clear that 
    $\lambda_{k} = 0$ if $k \ge M$ or $k \le -N -1$. We then have 
    \begin{alignat*}{3}
        & \Xi_{\bm{n};\bm{m},j} F_j^{(0)} = a_{-n_j,j}z^{-n_j} + \ldots + a_{|\bm{m}|-1,j}z^{|\bm{m}|-1} + \mathcal{O}(z^{|\bm{m}|}), \quad &&z \rightarrow 0, \\
        & \Xi_{\bm{n};\bm{m},j} F_j^{(\infty)} = b_{m_j-1,j}z^{m_j-1} + \ldots + b_{-|\bm{n}|+1,j}z^{-|\bm{n}|+1} + \mathcal{O}(z^{-|\bm{n}|}), \quad  &&z \rightarrow \infty,
    \end{alignat*}
    where we have
    % \begin{align}\label{eq:coefficients of series 1 I}
    %     a_{k,j} & = \kappa_{k,j} c_{0,j}+ 2\kappa_{k-1,j}c_{1,j} + \ldots + 2\kappa_{-n_j+1,j} c_{n_j+k-1,j}, \\ \label{eq:PadeToOrthoEq I w lambda 1} \lambda_k & = a_{k,1} + \ldots + a_{k,r},
    % \end{align}
    \begin{equation}
        \label{eq:coefficients of series 1 I}
        a_{k,j}  = \kappa_{k,j} c_{0,j}+ 2\kappa_{k-1,j}c_{1,j} + \ldots + 2\kappa_{-n_j,j} c_{n_j+k,j}
        %\\ \label{eq:PadeToOrthoEq I w lambda 1} \lambda_k & = a_{k,1} + \ldots + a_{k,r},
    \end{equation}
    for $k = -n_j,\ldots,|\bm{m}|-1$, and $a_{k,j}=0$ for $k <-n_j$, and
    % \begin{align}\label{eq:coefficients of series 2 I}
    %     b_{k,j} & = -\kappa_{k,j} c_{0,j} - 2\kappa_{k+1,j} c_{-1,j} - \ldots - 2 \kappa_{m_j,j} c_{-m_j+k,j}, \\ \label{eq:PadeToOrthoEq I w lambda 2}\lambda_k & = b_{k,1} + \ldots + b_{k,r}.
    % \end{align}
    \begin{equation}
        \label{eq:coefficients of series 2 I}
        b_{k,j} = -\kappa_{k,j} c_{0,j} - 2\kappa_{k+1,j} c_{-1,j} - \ldots - 2 \kappa_{m_j-1,j} c_{-m_j+k+1,j}
        %\\ \label{eq:PadeToOrthoEq I w lambda 2}\lambda_k & = b_{k,1} + \ldots + b_{k,r}.
    \end{equation}
    for $k = -|\bm{n}|+1,\ldots,m_j-1$, and $b_{k,j} =0$ for $ k>m_j-1$.
    
    % \begin{alignat*}{2}
    %     % a_{k,j}  & = \kappa_{k,j} c_{0,j}+ 2\kappa_{k-1,j}c_{1,j} + \ldots + 2\kappa_{-n_j,j} c_{n_j+k,j},
    %     %   \quad &\mbox{if }& -n_j\le k\le |\bm{m}|-1,\\
    %     % a_{k,j}  & = 0, & \mbox{if }& k <-n_j
    %     % \\
    %     a_{k,j}  & = 
    %     \begin{cases}
    %         \kappa_{k,j} c_{0,j}+ 2\kappa_{k-1,j}c_{1,j} + \ldots + 2\kappa_{-n_j,j} c_{n_j+k,j},
    %        &\mbox{if }-n_j\le k\le |\bm{m}|-1,\\
    %       0,&\mbox{if }  k< n_j,
    %     \end{cases}
    %     \\
    %     b_{k,j}  & = 
    %     \begin{cases}
    %         -\kappa_{k,j} c_{0,j} - 2\kappa_{k+1,j} c_{-1,j} - \ldots - 2 \kappa_{m_j-1,j} c_{-m_j+k+1,j},
    %        &\mbox{if }-|\bm{n}|+1\le k \le m_j-1,\\
    %       0,&\mbox{if }  k>m_j-1.
    %     \end{cases}
    % \end{alignat*}
    
    %Let us put $a_{k,j} = 0$ for any $k<-n_j+1$, $b_{k,j} = 0$ for any $k>m_j$. 

    %Assume that $|\bm{n}|-1 \ge N =  \max\{n_j\}$ and $|\bm{m}|-1 \ge M =  \max\{m_j\}$, that is $\bm{n}$ has at least two non-zero components, and the same is true of $\bm{m}$. 
    To get \eqref{eq:HPtypeI1swapped} and \eqref{eq:HPtypeI2swapped} we then necessarily need
    \begin{alignat}{3}
        \label{eq:lambda1}
        \lambda_k + \sum_{j=1}^r a_{k,j} & = 0, \qquad && k\le |\bm{m}|-1,
        \\
        \label{eq:lambda2}
        \lambda_k + \sum_{j=1}^r b_{k,j} & = 0, \qquad && k\ge |\bm{n}|-1.
    \end{alignat}
    In particular, this implies
    \begin{equation}
        \label{eq:PadeToOrthoEq I}
          \sum_{j=1}^r a_{k,j} = \sum_{j=1}^r b_{k,j},  \quad k = -|\bm{n}|+1,\dots,|\bm{m}|-1,
    \end{equation}
    which, combined with~\eqref{eq:coefficients of series 1 I},  \eqref{eq:coefficients of series 2 I} and~\eqref{eq:moments linear functional}, turns into~\eqref{eq:moprlI}. From \eqref{eq:lambda1} and \eqref{eq:lambda2}, we also see that $\Upsilon_{\bm{n};\bm{m}}$ is uniquely determined by $\bm{\Xi}_{\bm{n};\bm{m}}$. Observe that \eqref{eq:lambda1} and \eqref{eq:lambda2} is consistent with $\lambda_{k} = 0$ if $k \ge M$ (since in this case $b_{k,j}=0$ for all $j$) or $k \le -N -1$ (in this case $a_{k,j}=0$ for all $j$).
    
    %We also necessarily need
    % \begin{align}
    %      \label{eq:PadeToOrthoEq1 I}
    %      & -\lambda_{k} = \sum_{j=1}^r a_{k,j}, & k &= -N,\dots,-|\bm{n}|, \\ 
    %      \label{eq:PadeToOrthoEq2 I}
    %      & -\lambda_{k,j} =  \sum_{j=1}^r a_{k,j} = \sum_{j=1}^r b_{k,j},  & k &= -|\bm{n}|+1,\dots,|\bm{m}|-1, \\ 
    %      \label{eq:PadeToOrthoEq3 I}
    %      & -\lambda_{k,j} = \sum_{j=1}^r b_{k,j},  & k &= |\bm{m}|,\dots,M.
    % \end{align}
    %\begin{equation}
     %\label{eq:PadeToOrthoEq I w lambda}
      %   -\lambda_{k} = \sum_{j=1}^r b_{k,j},  \quad k = -N+1,\dots,M, 
    %\end{equation}
    %assuming that $M\le |\bm{m}|-1$. The case $M = |\bm{m}|$ (that is, all the components of $\bm{m}$, except for potentially one of them, are zero) seem to be special in that $\Upsilon$ is then {\it not} uniquely defined. 
    
    
    %If $M\le |\bm{m}|-1$ then~\eqref{eq:PadeToOrthoEq I w lambda} determines $\Upsilon_{\bm{n};\bm{m}}$ uniquely, given $\bm{\Xi}_{\bm{n};\bm{m}}$.   
    % \begin{equation}
    %      \int \big(\kappa_{|\bm{\ell}|}z^{|\bm{\ell}|} + \dots + \kappa_{-|\bm{\ell}|}z^{-|\bm{\ell}|}\big)z^{-k}d\mu_j(z) = 0, \qquad k = -\ell_j,\dots,\ell_j-1.
    % \end{equation}
    Now consider the formal power series
    \begin{align}
    \label{eq:secondKindFunctionTypeI 1}
    R^{(0)}_{\bm{n};\bm{m}}(z) 
    & = \sum_{j=1}^r \Big( L_j[\Xi_{\bm{n};\bm{m},j}(w)]+ 2\sum_{k=1}^\infty z^k L_j[\Xi_{\bm{n};\bm{m},j}(w) w^{-k}] \Big) , % & z & \to0, %\mbox{if } |z|<|w|, 
    \\
    \label{eq:secondKindFunctionTypeI 2}
     R^{(\infty)}_{\bm{n};\bm{m}}(z)   & =     - \sum_{j=1}^r \Big( L_j[\Xi_{\bm{n};\bm{m},j}(w)] +
            2\sum_{k=1}^\infty z^{-k} L_j[\Xi_{\bm{n};\bm{m},j}(w) w^{k}] \Big). %& z & \to\infty. %\mbox{if } |z|>|w|. \\
    \end{align}

    Using \eqref{eq:moprlI*}, we get 
    \begin{alignat}{3}\label{eq:R0 I}
        R_{\bm{n};\bm{m}}^{(0)}(z) & = \sum_{j=1}^r  L_j[\Xi_{\bm{n};\bm{m},j}(w)] + 
        \mathcal{O}(z^{|\bm{m}|}), \qquad && z \to 0,
        \\
        \label{eq:Rinfty I}
        R_{\bm{n};\bm{m}}^{(\infty)}(z)
        & = 
         - \sum_{j=1}^r  L_j[\Xi_{\bm{n};\bm{m},j}(w)] +
         \mathcal{O}(z^{-|\bm{n}|}), \qquad && z \to \infty.
     \end{alignat}
    The constant terms in~\eqref{eq:R0 I} are only relevant if $\bm{n}=\bm{0}$ and otherwise can be removed. Similarly, the constant terms
    in~\eqref{eq:Rinfty I} are only relevant if $\bm{m}=\bm{0}$. 

    Now denote $\widetilde{\Upsilon}_{\bm{n};\bm{m}}(z)$ to be the right-hand side of~\eqref{eq:upsilon}. It is clear that it is a Laurent polynomial in $\operatorname{span}\big\{z^k\big\}_{k = -N}^{M-1}$. Expanding algebraically the right-hand side of~\eqref{eq:upsilon} into power series, we get  
    \begin{align}
        \widetilde{\Upsilon}_{\bm{n};\bm{m}}(z)
        & =
        R^{(0)}_{\bm{n};\bm{m}}(z) - \sum_{j=1}^r \Xi_{\bm{n};\bm{m},j}(z) F_j^{(0)}(z) +\delta_{\bm{n};\bm{m}} \sum_{j = 1}^r L_j[\Xi_{\bm{n};\bm{m},j}(w)],
        \\
        \widetilde{\Upsilon}_{\bm{n};\bm{m}}(z)
        & =
        R^{(\infty)}_{\bm{n};\bm{m}}(z) - \sum_{j=1}^r \Xi_{\bm{n};\bm{m},j}(z) F^{(\infty)}_j(z) +\delta_{\bm{n};\bm{m}} \sum_{j = 1}^r L_j[\Xi_{\bm{n};\bm{m},j}(w)].
    \end{align}
    
    Combining this with~\eqref{eq:R0 I} and~\eqref{eq:Rinfty I}, we see that
%     \begin{align}\label{eq:two point hermite pade eq tilde}
%     & \varphi_{\bm{n}}(z)F_j(z) + \widetilde\psi_{\bm{n},j}(z) = \mathcal{O}(z^{\ell_j}),  &z& \rightarrow 0, \\\label{eq:two point hermite pade eq 2 tilde}
%     & \varphi_{\bm{n}}(z)F_j(z) + \widetilde\psi_{\bm{n},j}(z) = \mathcal{O}(z^{-\ell_j-1}),  &z& \rightarrow \infty.
% \end{align}
\eqref{eq:HPtypeI1swapped}--\eqref{eq:HPtypeI2swapped} hold with the same $\bm{\Xi}_{\bm{n};\bm{m}}$ but with $\Upsilon_{\bm{n};\bm{m}}$ replaced by $\widetilde\Upsilon_{\bm{n};\bm{m}}$.  
    But the proof showed that this uniquely determines all the coefficients of $\widetilde\Upsilon_{\bm{n};\bm{m}}$. This shows that $\widetilde\Upsilon_{\bm{n};\bm{m}} = \Upsilon_{\bm{n};\bm{m}}$ and proves~\eqref{eq:upsilon}.
    
    Conversely, given $\bm{\Xi}_{\bm{n};\bm{m}}$ and the corresponding $\Upsilon_{\bm{n};\bm{m}}$ as in~\eqref{eq:upsilon}, define $R_{\bm{n};\bm{m}}$ as in~\eqref{eq:secondKindFunctionTypeI 1}--\eqref{eq:secondKindFunctionTypeI 2}. As above, we show that it satisfies~\eqref{eq:R0 I} and ~\eqref{eq:Rinfty I}, which then becomes ~\eqref{eq:HPtypeI1swapped}, \eqref{eq:HPtypeI2swapped}.
\end{proof}

The corresponding type I problem for the ``starred'' polynomials looks as follows. We are looking for the Laurent polynomials 
\begin{equation}
    \label{eq:HPtypeI*spanSwapped}
     \Xi^*_{\bm{n};\bm{m},j}(z) \in \operatorname{span}\big\{z^{k}\big\}_{k=-n_j+1}^{m_j},
    \qquad \Upsilon^*_{\bm{n};\bm{m}}(z) \in \operatorname{span}\big\{z^{k}\big\}_{k=-N+1}^{M},
\end{equation}
%$\Xi_{\bm{n};\bm{m},j}^*$ will be in $\operatorname{span}\set{z^k}_{k=-m_j}^{n_j-1}$, solving the Hermite-Pad\'e problem
that satisfy
\begin{alignat}{3}
    \label{eq:HPtypeI*1swapped}
    & \sum_{j = 1}^r \Xi^*_{\bm{n};\bm{m},j}(z)F_j^{(0)}(z) - \Upsilon^*_{\bm{n};\bm{m}}(z) = \mathcal{O}(z^{\abs{\bm{m}}}), \qquad &&z \rightarrow 0, \\
    \label{eq:HPtypeI*2swapped}
    &\sum_{j = 1}^r \Xi^*_{\bm{n};\bm{m},j}(z)F_j^{(\infty)}(z) -\Upsilon^*_{\bm{n};\bm{m}}(z) = \mathcal{O}(z^{-\abs{\bm{n}}}),  \qquad &&z \rightarrow \infty.
\end{alignat}
%for some Laurent polynomial $\Upsilon_{\bm{n};\bm{m}}^*$. 

\begin{thm}
    Any vector $\bm{\Xi}_{\bm{n};\bm{m}}^* = (\Xi_{\bm{n};\bm{m},1}^*,\dots,\Xi_{\bm{n};\bm{m},r}^*)$ that solves ~\eqref{eq:HPtypeI*spanSwapped}, \eqref{eq:HPtypeI*1swapped}, \eqref{eq:HPtypeI*2swapped} must satisfy the orthogonality relations \eqref{eq:moprlI*}, and $\Upsilon_{\bm{n};\bm{m}}^*$ is given by
    \begin{equation}\label{eq:upsilon star}
        \Upsilon_{\bm{n};\bm{m}}^*(z) = \sum_{j = 1}^r \Big( L_j\Big[\frac{w+z}{w-z}\big( \Xi_{\bm{n};\bm{m},j}^*(z) - \Xi_{\bm{n};\bm{m},j}^*(w) \big)\Big] + \delta_{\bm{n};\bm{m}} L_j[\Xi_{\bm{n};\bm{m},j}^*(w)] \Big),
    \end{equation}
    {where $\delta_{\bm{n};\bm{m}}$ is the same as in \eqref{eq:upsilon}.}

    Conversely, any vector $\bm{\Xi}_{\bm{n};\bm{m}}^*$ with $\Xi^*_{\bm{n};\bm{m},j} \in \operatorname{span}\big\{z^{k}\big\}_{k=-n_j+1}^{m_j}$ that satisfies  \eqref{eq:moprlI*} solves  ~\eqref{eq:HPtypeI*spanSwapped}, \eqref{eq:HPtypeI*1swapped}, \eqref{eq:HPtypeI*2swapped} together with~\eqref{eq:upsilon star}.
\end{thm}

%Similarly to Proposition \ref{prop:laurent normality type I}, there 



%\begin{rem}
%    It is easy to verify that the type I polynomial $\bm{\Lambda}_{\bm{n}}(z)$ from~\cite{MOPUC1,KVMOPUC} is $z^{-1} \bm{\Xi}_{\bm{0};\bm{n}}(z)$. While the extra factor of $z^{-1}$ is a mere convenience here, observe that the order of indices $(\bm{0};\bm{n})$ is reversed here compared with $\Phi_{\bm{n}} = \varphi_{\bm{n};\bm{0}}$ (see Remark~\ref{rem:typeIIvsHP}). By Proposition~\ref{prop:laurent normality type I} working with $\bm{\Xi}_{\bm{n};\bm{0}}$ and $\bm{\Xi}^*_{\bm{n};\bm{0}}$ here is more natural than with $\bm{\Xi}_{\bm{0};\bm{n}}$ and $\bm{\Xi}^*_{\bm{0};\bm{n}}$.    
%    Of course, if $L_j$ correspond to positive measures  on $\T$ (as is the setting in~\cite{MOPUC1,KVMOPUC}), then normalities of $(\bm{0};\bm{n})$ and $(\bm{n};\bm{0})$ are equivalent by Proposition~\ref{thm:reversal vs star polynomials}.  
    
%    Also, the type I function $\xi_{\bm{n}}$ (defined in Proposition~\ref{prop:typeI}) is a special case $\Xi_{\bm{n}/2;\bm{n}/2}$ (defined in Proposition~\ref{prop:laurent normality type I}) but only if each $n_j$ is even.
%\end{rem}
  
%     \begin{rem}
%     In particular, if all $n_j$'s are even, then the type I vector of functions $\bm{\xi}_{\bm{n}}$ that we defined in~\eqref{eq:type I def0}--\eqref{eq:type I def} coincided with the vector $\bm{\Xi}_{\bm{n}/2;\bm{n}/2}$.
% \end{rem}


%\begin{rem}
%    The type I polynomials we worked with in \cite{KVMOPUC}. 
%\end{rem}

%%%%%%%%%%%%%%%%%%%%%%%%%%%%%%%%%%%%%%%%%%%%%%%%%%%%%%%%%%%%%%%%%%%%%%%%%%%%%

\section{Szeg\H{o} Recurrence Relations}\label{ss:SzegoNEW}
Let us again allow $(\bm{n};\bm{m})\in\mathfrak{C}_{2r}$ as in Section~\ref{ss:normality}.
Here we establish that the Laurent polynomials $\Phi_{\bm{n};\bm{m}}$ satisfy a family of recurrence relations, which follow directly from our earlier results in~\cite{KVMOPUC}. That work was restricted to orthogonality with respect to positive measures as in~\eqref{eq:functionals given by measures} and to {\it non-Laurent} polynomials. The underlying arguments, however, extend with only minor modifications to the general case of moment functionals~\eqref{eq:moments linear functional}. This yields the following generalization.
%By following the exact same prove, we obtain 

\begin{thm}\label{thm:Szego1}
    %If $(\bm{n};\bm{m})\in\mathfrak{C}_{2r}$ is normal, along with all the neighbouring indices that belong to $\Z_+^r$, then
    Assuming all the $\mathfrak{C}_{2r}$-indices that appear in the equations below are normal, we have the Szeg\H{o} recurrence relations 
    \begin{align}
         \label{eq:SzegoII1}
        \vphantom{\sum_{j = 1}^{r}} 
        & \Phi_{\bm{n};\bm{m}}^*(z) = \Phi_{\bm{n}-\bm{e}_k;\bm{m}}^*(z) + \beta_{\bm{n};\bm{m}}z\Phi_{\bm{n}-\bm{e}_k;\bm{m}}(z), 
        \\ \label{eq:SzegoII2}
        & \Phi_{\bm{n};\bm{m}}(z) = \alpha_{\bm{n};\bm{m}}\Phi_{\bm{n};\bm{m}}^*(z) + \sum_{j = 1}^r \rho_{\bm{n};\bm{m},j}z\Phi_{\bm{n}-\bm{e}_j;\bm{m}}(z), 
        % \\ 
        % \label{eq:SzegoII3}
        % & \Phi_{\bm{n};\bm{m}}^*(z)= \beta_{\bm{n};\bm{m}}\Phi_{\bm{n};\bm{m}}(z) + \sum_{j = 1}^r \rho_{\bm{n};\bm{m},j}\Phi_{\bm{n}-\bm{e}_j;\bm{m}}^*(z),
    \end{align}
    and
    \begin{align}
        \label{eq:SzegoI1}
        & \vphantom{\sum_{j = 1}^{r}} 
        \bm{\Xi}^*_{\bm{n};\bm{m}}(z) = \bm{\Xi}^*_{\bm{n}+\bm{e}_k;\bm{m}}(z) -\alpha_{\bm{n};\bm{m}}z\bm{\Xi}_{\bm{n}+\bm{e}_k;\bm{m}}(z),
        \\
        \label{eq:SzegoI2}
        & \bm{{\Xi}}_{\bm{n};\bm{m}}(z) = -\beta_{\bm{n};\bm{m}}\bm{\Xi}^*_{\bm{n};\bm{m}}(z) + \sum_{j = 1}^{r}\rho_{\bm{n};\bm{m},j}z\bm{\Xi}_{\bm{n}+\bm{e}_j;\bm{m}}(z),
        % \\
        % \label{eq:SzegoI3}
        % & \bm{\Xi}_{\bm{n};\bm{m}}^*(z)= -\alpha_{\bm{n};\bm{m}}\bm{\Xi}_{\bm{n};\bm{m}}(z) + \sum_{j = 1}^r \rho_{\bm{n};\bm{m},j}\bm{\Xi}^*_{\bm{n}+\bm{e}_j;\bm{m}}(z),
\end{align}
for some complex numbers $\rho_{\bm{n};\bm{m},j}$.
\end{thm}


If we keep the first index constant, then we obtain another family of relations.

\begin{thm}\label{thm:Szego2}
Assuming the indices that appear in the equations below are normal, we have the Szeg\H{o} recurrence relations 
\begin{align}
         \label{eq:SzegoII4}
        \vphantom{\sum_{j = 1}^{r}} 
        & \Phi_{\bm{n};\bm{m}}(z) = \Phi_{\bm{n},\bm{m}-\bm{e}_k}(z) + \alpha_{\bm{n};\bm{m}}z^{-1}\Phi^*_{\bm{n},\bm{m}-\bm{e}_k}(z), 
        \\ \label{eq:SzegoII5}
        & \Phi^*_{\bm{n};\bm{m}}(z) = \beta_{\bm{n};\bm{m}}\Phi_{\bm{n};\bm{m}}(z) + \sum_{j = 1}^r \sigma_{\bm{n};\bm{m},j}z^{-1} \Phi_{\bm{n},\bm{m}-\bm{e}_j}^*(z), 
        % \\ 
        % \label{eq:SzegoII6}
        % & \Phi_{\bm{n};\bm{m}}(z)= \alpha_{\bm{n};\bm{m}}\Phi^*_{\bm{n};\bm{m}}(z) + \sum_{j = 1}^r \sigma_{\bm{n};\bm{m},j}\Phi_{\bm{n};\bm{m}-\bm{e}_j}(z),
    \end{align}
    and
    \begin{align}
        \label{eq:SzegoI4}
        & \vphantom{\sum_{j = 1}^{r}} 
        \bm{\Xi}_{\bm{n};\bm{m}}(z) = \bm{\Xi}_{\bm{n};\bm{m}+\bm{e}_k}(z) -\beta_{\bm{n};\bm{m}}z^{-1} \bm{\Xi}^*_{\bm{n};\bm{m}+\bm{e}_k}(z),
        \\
        \label{eq:SzegoI5}
        & \bm{{\Xi}}^*_{\bm{n};\bm{m}}(z) = -\alpha_{\bm{n};\bm{m}}\bm{\Xi}_{\bm{n};\bm{m}}(z) + \sum_{j = 1}^{r}\sigma_{\bm{n};\bm{m},j}z^{-1} \bm{\Xi}^*_{\bm{n};\bm{m}+\bm{e}_j}(z),
        % \\
        % \label{eq:SzegoI6}
        % & \bm{\Xi}_{\bm{n};\bm{m}}(z)= -\beta_{\bm{n};\bm{m}}\bm{\Xi}^*_{\bm{n};\bm{m}}(z) + \sum_{j = 1}^r \sigma_{\bm{n};\bm{m},j}\bm{\Xi}_{\bm{n};\bm{m}+\bm{e}_j}(z),
\end{align}
for some complex numbers $\sigma_{\bm{n};\bm{m},j}$.
\end{thm}
\begin{rem}
     One can formally put $\Phi_{\bm{n}-\bm{e}_j;\bm{m}} = 0$ and $\rho_{\bm{n};\bm{m},j} = 0$ when $(\bm{n}-\bm{e}_j;\bm{m})\notin \mathfrak{C}_{2r}$. With such a choice, equality~\eqref{eq:SzegoII2} 
     %and \eqref{eq:SzegoII3} 
     holds even if some of the indices $(\bm{n}-\bm{e}_j;\bm{m})$ 
     %, $(\bm{n};\bm{m}-\bm{e}_j)$
     fall outside of $\mathfrak{C}_{2r}$. Similar modifications apply to \eqref{eq:SzegoII5}. % and \eqref{eq:SzegoII6}.
     %One can formally put $\Phi_{\bm{n};\bm{m}} = 0$ for those indices $(\bm{n};\bm{m})\notin \mathfrak{C}_{2r}$. With such a choice, equalities~\eqref{eq:SzegoII2}, \eqref{eq:SzegoII3}, \eqref{eq:SzegoII5}, \eqref{eq:SzegoII6} hold even if some of the indices $(\bm{n}-\bm{e}_j;\bm{m})$, $(\bm{n};\bm{m}-\bm{e}_j)$ fall outside of $\mathfrak{C}_{2r}$. 
\end{rem}
\begin{proof}[Proof of Theorems~\ref{thm:Szego1} and~\ref{thm:Szego2}]
        The proof of recurrences~\eqref{eq:SzegoII1}and \eqref{eq:SzegoII2}
        %, \eqref{eq:SzegoII3} 
        goes through the exact same arguments as~\cite[Thm 3.1 and Cor 3.3]{KVMOPUC}. 

        Relations ~\eqref{eq:SzegoII4} and \eqref{eq:SzegoII5}
        % \eqref{eq:SzegoII6} 
        follow then from Theorem~\ref{thm:reversal vs star polynomials}. 
        %The fact that the recurrence coefficients $\alpha_{\bm{n};\bm{m}}$ in~\eqref{eq:SzegoII1} and~\eqref{eq:SzegoII4} coincide follows by examining the $z^{-m_j}$-coefficients.
        % Observe that $\alpha_{\bm{n};\bm{m}}$ and $\beta_{\bm{n};\bm{m}}$ coefficients in ~\eqref{eq:SzegoII1}--\eqref{eq:SzegoII3} and in ~\eqref{eq:SzegoII4}--\eqref{eq:SzegoII6}  coincide (just check the leading the trailing coefficients in the recurrences), the $\rho_{\bm{n};\bm{m},j}$ and $\sigma_{\bm{n};\bm{m},j}$ coefficients are different in general. 
        % \begin{prop}
        %     In the setting of 
        % Theorem~\ref{thm:reversal vs star polynomials}, we have
        % \begin{equation}
        %     \rho_{\bm{n};\bm{m},j}(\bm{L}^\sharp) = \bar\sigma_{\bm{m};\bm{n},j}(\bm{L}), \qquad \sigma_{\bm{n};\bm{m},j}(\bm{L}^\sharp) = \bar\rho_{\bm{m};\bm{n},j}(\bm{L}).
        % \end{equation}
        % \end{prop}
        % In fact in the setting of 
        %  Theorem~\ref{thm:reversal vs star polynomials}, we have
        %  \begin{equation}
        %      \rho_{\bm{n};\bm{m},j}(\bm{L}^\sharp) = \bar\sigma_{\bm{m};\bm{n},j}(\bm{L}), \qquad \sigma_{\bm{n};\bm{m},j}(\bm{L}^\sharp) = \bar\rho_{\bm{m};\bm{n},j}(\bm{L}).
        % \end{equation}
        
        The proof of recurrences~\eqref{eq:SzegoI1},and \eqref{eq:SzegoI2}
        %, \eqref{eq:SzegoI3} 
        follow similar lines to the arguments in~\cite[Thm 4.3]{KVMOPUC} with modifications related to the discussion in Remarks~\ref{rem:annoying1} and \ref{rem:annoying2}. Proposition~\ref{prop:general hermite pade measures type I} below summarizes the main modifications that are needed in order to adapt the proofs from ~\cite[Sect. 4]{KVMOPUC}.
    %     Finally,  relations ~\eqref{eq:SzegoI4} and \eqref{eq:SzegoI5}
    % %, \eqref{eq:SzegoI6} 
    % follow again from Theorem~\ref{thm:reversal vs star polynomials}. 
\end{proof}

\begin{prop}\label{prop:general hermite pade measures type I}Assume the indices that appear in the respective relations below are normal. 
    %Suppose $L_j[w^n] = \int w^n d\mu_j(w)$ for some probability measure supported on the unit circle, for $n \in \Z$ and $j = 1,\dots,r$. Then we have $\bm{\Xi}_{\bm{n};\bm{m}}^\sharp = \bm{\Xi}_{\bm{m};\bm{n}}^*$, and
    %$\overline{\bm{\Xi}_{\bm{n};\bm{m}}(1/\bar{z})} = \bm{\Xi_{\bm{m};\bm{n}}^*}(z)$, and
    \noindent\hfill
    \begin{itemize}
        \item[(i)] Let $\kappa_{\bm{n};\bm{m},j}$ be the $z^{-n_j}$ coefficient of $\Xi_{\bm{n};\bm{m},j}(z)$ and let $\ell_{\bm{n};\bm{m},j}$ be the $z^{m_j}$ coefficient of $\Xi^*_{\bm{n};\bm{m},j}(z)$. Then
        \begin{align}
            \label{eq:kappa}
            \kappa_{\bm{n}+\bm{e}_j;\bm{m},j} & = \frac{1}{L_j[\Phi_{\bm{n};\bm{m}}(w)w^{-n_j}]},
            \\
            \label{eq:l}
            \ell_{\bm{n}+\bm{e}_j;\bm{m},j} &= \frac{1}{L_j[\Phi^*_{\bm{n};\bm{m}}(w)w^{m_j}]}.
        \end{align}
        
        % \item[(ii)] . Then
        % \begin{equation}
        %     \label{eq:l}
        %     \ell_{\bm{n}+\bm{e}_j;\bm{m},j} = \frac{1}{L_j[\Phi^*_{\bm{n};\bm{m}}(w)w^{m_j}]}.
        % \end{equation}

        \item[(ii)] We have the relations
        \begin{alignat}{1}
        \label{eq:rho}
            & \rho_{\bm{n};\bm{m},j}= \frac{L_j[\Phi_{\bm{n};\bm{m}}(w)w^{-n_j}]}{L_j[\Phi_{\bm{n}-\bm{e}_j;\bm{m}}(w)w^{-n_j+1}]},
            \\
            \label{eq:sigma}
            & \sigma_{\bm{n};\bm{m},j}= \frac{L_j[\Phi^*_{\bm{n};\bm{m}}(w)w^{m_j}]}{L_j[\Phi^*_{\bm{n};\bm{m}-\bm{e}_j}(w)w^{m_j-1}]}.
        \end{alignat}
        Furthermore, $\rho_{\bm{n};\bm{m},j} \neq 0$ if and only if $(\bm{n}+\bm{e}_j;\bm{m})$ is normal. Similarly, $\sigma_{\bm{n};\bm{m},j} \neq 0$ if and only if $(\bm{n};\bm{m}+\bm{e}_j)$ is normal.
        
        \item[(iii)] The following relations hold:
        \begin{alignat}{3}
            \label{eq:biorthogonality 1}
            & \sum_{j = 1}^r L_j \big[ \Xi_{\bm{n};\bm{m}{,j}}(w)w^{-\abs{\bm{m}}} \big] = -{\beta}_{\bm{n};\bm{m}}, 
            \\ \label{eq:biorthogonality 2}
            & \sum_{j = 1}^r L_j \big[ \Xi_{\bm{n};\bm{m}{,j}}^*(w) w^{\abs{\bm{n}}} \big] = -{\alpha}_{\bm{n};\bm{m}}.
        \end{alignat}
    \end{itemize}
\end{prop}        
\begin{rem}
    Note that the right hand side of \eqref{eq:rho} only requires normality of $(\bm{n};\bm{m})$ and $(\bm{n}-\bm{e}_j;\bm{m})$. It is possible to derive \eqref{eq:SzegoII2} without appealing to normality of indices of the form $(\bm{n}-\bm{e}_{\ell};\bm{m})$, $\ell\ne j$. In that case any choice of polynomials $\Phi_{\bm{n}-\bm{e}_{\ell};\bm{m}}$,  $\ell\ne j$, will work for the relation to hold. When $(\bm{n}-\bm{e}_j;\bm{m})$ is normal for a given $j$ then $\rho_{\bm{n};\bm{m},j}$ is uniquely defined and can be computed to \eqref{eq:rho}. Similarly, $\sigma_{\bm{n};\bm{m}}$ is defined as long as $(\bm{n};\bm{m})$ and $(\bm{n};\bm{m}-\bm{e}_j)$ are normal. 
\end{rem}

    
\begin{rem}\label{rem:reversal}
    It is easy to see that in the setting of 
         Theorem~\ref{thm:reversal vs star polynomials}, we have
         \begin{equation}
             \rho_{\bm{n};\bm{m},j}(\bm{L}^\sharp) = \bar\sigma_{\bm{m};\bm{n},j}(\bm{L}), \qquad \sigma_{\bm{n};\bm{m},j}(\bm{L}^\sharp) = \bar\rho_{\bm{m};\bm{n},j}(\bm{L}),
        \end{equation}
        and if each $L_j$ is associated with a probability measure then
        \begin{equation}
             \rho_{\bm{n};\bm{m},j}= \bar\sigma_{\bm{m};\bm{n},j} .
        \end{equation}
\end{rem}

%In the last section we will need to use the following important result. Note that it allows to express $L_j[\Phi_{\bm{n};\bm{m}}(w) w^{-n_j}]$ in terms of $\alpha$'s and $\beta$'s, similarly as in the classical $r=1$ theory.

% \begin{prop}
%     Suppose all the indices in the following equation and all of its neighbours are normal. Then
%     \begin{align}\label{eq:kappaAlphaBeta}
%         \frac{\kappa_{\bm{n};\bm{m}-\bm{e}_j,j}}{\kappa_{\bm{n};\bm{m},j}} & = 1-\alpha_{\bm{n}-\bm{e}_j;\bm{m}}  \beta_{\bm{n};\bm{m}-\bm{e}_j},
%         \\
%         \frac{\ell_{\bm{n}-\bm{e}_j;\bm{m},j}}{\ell_{\bm{n};\bm{m},j}} & = 1-\alpha_{\bm{n};\bm{m}-\bm{e}_j}  \beta_{\bm{n}-\bm{e}_j;\bm{m}}.
%     \end{align}

%     In particular, $\alpha_{\bm{n}-\bm{e}_j;\bm{m}}  \beta_{\bm{n};\bm{m}-\bm{e}_j} \ne 1$ and $\alpha_{\bm{n};\bm{m}-\bm{e}_j}  \beta_{\bm{n}-\bm{e}_j;\bm{m}}\ne 1$.
% \end{prop}
% \begin{proof}
%     Denote the $z^{-n_j+1}$ coefficient of $\Xi_{\bm{n};\bm{m},j}$ by $d_{\bm{n};\bm{m},j}$. 
%     Consider the $j$-th component of the vector equation~\eqref{eq:SzegoI4} with $k=j$. Computing the     
%     $z^{-n_j}$  coefficients of each term, we get
%     \begin{equation}\label{eq:tempkappas}
%         \kappa_{\bm{n};\bm{m},j} = \kappa_{\bm{n};\bm{m}+\bm{e}_j,j} - \beta_{\bm{n};\bm{m}}  d_{\bm{n};\bm{m}+\bm{e}_j,j}.
%     \end{equation}

%     Now consider the $j$-th component of the vector equation~\eqref{eq:SzegoI1} with $k=j$. Computing the     
%     $z^{-n_j}$  coefficients of each term, we get
%     \begin{equation}\label{eq:tempkappas2}
%         0 =  d_{\bm{n}+\bm{e}_j;\bm{m},j} - \alpha_{\bm{n};\bm{m}} \kappa_{\bm{n}+\bm{e}_j;\bm{m},j} .
%     \end{equation}

%     Combining ~\eqref{eq:tempkappas} with $(\bm{n};\bm{m})$ replaced with $(\bm{n};\bm{m}-\bm{e}_j)$ and~\eqref{eq:tempkappas2} with $(\bm{n};\bm{m})$ replaced with $(\bm{n}-\bm{e}_j;\bm{m})$ produces
%     \begin{align}
%         \kappa_{\bm{n};\bm{m}-\bm{e}_j,j} &= \kappa_{\bm{n};\bm{m},j} - \beta_{\bm{n};\bm{m}-\bm{e}_j}  d_{\bm{n};\bm{m},j}
%         \\
%         &= 
%         \kappa_{\bm{n};\bm{m},j} - \beta_{\bm{n};\bm{m}-\bm{e}_j}  \alpha_{\bm{n}-\bm{e}_j;\bm{m}} \kappa_{\bm{n};\bm{m},j},
%     \end{align}
%     completing the proof of~\eqref{eq:kappaAlphaBeta}. The other equality can be proved in a similar fashion.
% \end{proof}

%%%%%%%%%%%%%%%%%%%%%%%%%%%%%%%%%%%%%%

\section{Compatibility relations}\label{ss:CC}

\begin{prop}\label{pr:compPoly1}
Assuming the indices that appear in the equations below are normal, and $k \neq \ell$, then there is a complex number $\gamma_{\bm{n};\bm{m}}^{k\ell}$ such that
\begin{align}\label{eq:CCII}
    & \Phi_{\bm{n}+\bm{e}_k;\bm{m}}-\Phi_{\bm{n}+\bm{e}_\ell;\bm{m}} = \gamma_{\bm{n};\bm{m}}^{k\ell}\Phi_{\bm{n};\bm{m}},
\\
\label{eq:CCI}
    & \bm\Xi_{\bm{n}-\bm{e}_k;\bm{m}}-\bm\Xi_{\bm{n}-\bm{e}_\ell;\bm{m}} = \gamma_{\bm{n}-\bm{e}_k-\bm{e}_\ell;\bm{m}}^{k\ell}\bm\Xi_{\bm{n};\bm{m}}.
\end{align}

Similarly, for some complex number $\eta_{\bm{n};\bm{m}}^{k\ell}$,
\begin{align}\label{eq:CCII2}
    & \Phi^*_{\bm{n};\bm{m}+\bm{e}_k}-\Phi^*_{\bm{n};\bm{m}+\bm{e}_\ell} = \eta_{\bm{n};\bm{m}}^{k\ell}\Phi^*_{\bm{n};\bm{m}},
    \\
    \label{eq:CCIInew}
    & \bm\Xi^*_{\bm{n};\bm{m}-\bm{e}_k}-\bm\Xi^*_{\bm{n};\bm{m}-\bm{e}_\ell} = \eta_{\bm{n};\bm{m}-\bm{e}_k-\bm{e}_\ell}^{k\ell}\bm\Xi^*_{\bm{n};\bm{m}}.
\end{align}
\end{prop}

\begin{proof}
    Check orthogonality relations and degree restrictions of $\Phi_{\bm{n}+\bm{e}_k;\bm{m}}-\Phi_{\bm{n}+\bm{e}_\ell;\bm{m}}$, and compare with $\Phi_{\bm{n};\bm{m}}$. The first relation then follows by normality of $(\bm{n};\bm{m})$, and same arguments prove the other relations. 
\end{proof}

\begin{rem}
    Similarly to the previous section, we have, for $k\ne \ell$,
    \begin{align}\label{eq:gamma}
    \gamma^{k\ell}_{\bm{n};\bm{m}} &= \frac{L_\ell[\Phi_{\bm{n}+\bm{e}_k;\bm{m}}(w)w^{-n_\ell}]}{L_\ell[\Phi_{\bm{n};\bm{m}}(w)w^{-n_\ell}]},
    \\\label{eq:eta}
    \eta^{k\ell}_{\bm{n};\bm{m}} &= \frac{L_\ell[\Phi^*_{\bm{n};\bm{m}+\bm{e}_k}(w)w^{m_\ell}]}{L_\ell[\Phi^*_{\bm{n};\bm{m}}(w)w^{m_\ell}]},
    \end{align}
    and  
    \begin{equation}
             \gamma_{\bm{n};\bm{m}}^{k\ell}(\bm{L}^\sharp) = \bar\eta_{\bm{m};\bm{n}}^{k\ell}(\bm{L}), \qquad \eta_{\bm{n};\bm{m}}^{k\ell}(\bm{L}^\sharp) = \bar\gamma_{\bm{m};\bm{n},j}^{k\ell}(\bm{L}).
    \end{equation}
    Furthermore, $\gamma^{k\ell}_{\bm{n};\bm{m}} \neq 0$ if and only if $(\bm{n}+\bm{e}_k+\bm{e}_\ell;\bm{m})$ is normal, and $\eta_{\bm{n};\bm{m}}^{k\ell} \neq 0$ if and only if $(\bm{n};\bm{m}+\bm{e}_k+\bm{e}_\ell)$ is normal. The constants in \eqref{eq:CCI} and~\eqref{eq:CCIInew} exist even if $(\bm{n}-\bm{e}_k-\bm{e}_\ell;\bm{m})$, resp. $(\bm{n};\bm{m}-\bm{e}_k-\bm{e}_\ell)$, is not normal (see \cite{KVMOPUC} for the details).
\end{rem}

\begin{thm}
%Assuming all the $\N^r\times \N^r$ multi-indices are normal, we have
%If $(\bm{n};\bm{m})\in\mathfrak{C}_{2r}$ is normal, along with all of its neighbouring indices, then
%that belong to $\Z_+^r$, then
Assuming normality of all indices required for the recurrence coefficients below to be defined, we have the partial difference equations
     \begin{align} 
    \label{eq:compatibility condition 2}
    & \alpha_{\bm{n};\bm{m}}\beta_{\bm{n};\bm{m}} + \sum_{j = 1}^r \rho_{\bm{n};\bm{m},j} = 1,
    \\
    \label{eq:compatibility condition 1}
    & 
    \vphantom{\sum_{j = 1}^r }
    \alpha_{\bm{n}+\bm{e}_k;\bm{m}} - \alpha_{\bm{n}+\bm{e}_\ell;\bm{m}}
    = \alpha_{\bm{n};\bm{m}} \gamma_{\bm{n};\bm{m}}^{k\ell},
    \\
    \label{eq:another comp}
    \vphantom{\sum_{j = 1}^r }
    &\beta_{\bm{n}+\bm{e}_\ell;\bm{m}} - \beta_{\bm{n}+\bm{e}_k;\bm{m}}
    = \beta_{\bm{n}+\bm{e}_\ell+\bm{e}_k;\bm{m}} \gamma_{\bm{n};\bm{m}}^{k\ell},
    \\
    \vphantom{\sum_{j = 1}^r }
    \label{eq:compatibility condition 3.1}
    & \rho_{\bm{n};\bm{m},k} \gamma_{\bm{n};\bm{m}}^{k\ell}= \rho_{\bm{n}+\bm{e}_\ell;\bm{m},k}, \gamma_{\bm{n}-\bm{e}_k;\bm{m}}^{k\ell} .
    \end{align}


Similarly,
    \begin{align} 
    \label{eq:compatibility condition 2_another}
    & \alpha_{\bm{n};\bm{m}}\beta_{\bm{n};\bm{m}} + \sum_{j = 1}^r \sigma_{\bm{n};\bm{m},j} = 1,
    \\
    \label{eq:compatibility condition 1_another}
    & 
    \vphantom{\sum_{j = 1}^r }
    \alpha_{\bm{n};\bm{m}+\bm{e}_\ell} - \alpha_{\bm{n};\bm{m}+\bm{e}_k}
    = \alpha_{\bm{n};\bm{m}+\bm{e}_\ell+\bm{e}_k} \eta_{\bm{n};\bm{m}}^{k\ell},
    \\
    \vphantom{\sum_{j = 1}^r }
    &\beta_{\bm{n};\bm{m}+\bm{e}_k} - \beta_{\bm{n};\bm{m}+\bm{e}_\ell}
    = \beta_{\bm{n};\bm{m}} \eta_{\bm{n};\bm{m}}^{k\ell},
    \\
    \vphantom{\sum_{j = 1}^r }
    \label{eq:compatibility condition 3.1_another}
    & \sigma_{\bm{n};\bm{m},k} \eta_{\bm{n};\bm{m}}^{k\ell}= \sigma_{\bm{n};\bm{m}+\bm{e}_\ell,k}, \eta_{\bm{n};\bm{m}-\bm{e}_k}^{k\ell} .
    \end{align}
\end{thm}

\begin{proof}
    \eqref{eq:compatibility condition 2} follows by comparing leading coefficients in \eqref{eq:SzegoII2}, and \eqref{eq:compatibility condition 1} follows by comparing the lowest coefficient in \eqref{eq:CCII}. \eqref{eq:another comp} can be obtained from \eqref{eq:CCI} and \eqref{eq:biorthogonality 1}. \eqref{eq:compatibility condition 3.1} is immediate from \eqref{eq:rho} and \eqref{eq:gamma}. Similar arguments then imply the dual relations \eqref{eq:compatibility condition 2_another}-\eqref{eq:compatibility condition 3.1_another}.
\end{proof}


%%%%%%%%%%%%%%%%%%%%%%%%%%%%%%%%%%%


\section{Consequences of the Szeg\H{o} Relations}\label{ss:extra}

Aside from the main recurrence relations of the previous sections, we want to mention some other recurrence relation that may be useful. We use two of them in the proof of the Geronimus relations for the Szeg\H{o} mapping in Section \ref{ss:SzegoMap}.

\begin{prop}
Assuming all the indices appearing in the equations below are normal, we have recurrence relations
    \begin{align}\label{eq:SzegoII7}
        & \Phi_{\bm{n}+\bm{e}_k;\bm{m}}^*(z) = (1 - \alpha_{\bm{n};\bm{m}+\bm{e}_k}\beta_{\bm{n}+\bm{e}_k;\bm{m}})\Phi_{\bm{n};\bm{m}}^*(z) + \beta_{\bm{n}+\bm{e}_k;\bm{m}}z\Phi_{\bm{n};\bm{m}+\bm{e}_k}(z), \\
        & \label{eq:SzegoII8}\Phi_{\bm{n};\bm{m}+\bm{e}_k}(z) = (1 - \alpha_{\bm{n};\bm{m}+\bm{e}_k}\beta_{\bm{n}+\bm{e}_k;\bm{m}})\Phi_{\bm{n};\bm{m}}(z) + \alpha_{\bm{n};\bm{m}+\bm{e}_k}z^{-1}\Phi_{\bm{n}+\bm{e}_k;\bm{m}}^*(z), 
        \end{align}
        and 
        \begin{align}& \bm{\Xi}_{\bm{n}-\bm{e}_k;\bm{m}}^*(z) = (1-\alpha_{\bm{n}-\bm{e}_k;\bm{m}}\beta_{\bm{n};\bm{m}-\bm{e}_k})\bm{\Xi}_{\bm{n};\bm{m}}^*(z) - \alpha_{\bm{n}-\bm{e}_k;\bm{m}}z\bm{\Xi}_{\bm{n};\bm{m}-\bm{e}_k}(z), \\
        & \bm{\Xi}_{\bm{n};\bm{m}-\bm{e}_k}(z) = (1-\alpha_{\bm{n}-\bm{e}_k;\bm{m}}\beta_{\bm{n};\bm{m}-\bm{e}_k})\bm{\Xi}_{\bm{n};\bm{m}}(z) - \beta_{\bm{n};\bm{m}-\bm{e}_k}z^{-1}\bm{\Xi}_{\bm{n}-\bm{e}_k;\bm{m}}^*(z).
    \end{align}
\end{prop}
\begin{proof}
    Multiply \eqref{eq:SzegoII4} for the index $(\bm{n};\bm{m}+\bm{e}_k)$ by $\beta_{\bm{n}+\bm{e}_k;\bm{m}}z$ and substitute the term $\beta_{\bm{n}+\bm{e}_k;\bm{m}}z\Phi_{\bm{n};\bm{m}}(z)$ using \eqref{eq:SzegoII1} to get \eqref{eq:SzegoII7}. The other equations follow similarly. 
\end{proof}

\begin{cor}Assuming normality of the indices that appear below, we have 
\begin{equation}
    \label{eq:one minus alphabeta}
      \frac{L_k[\Phi_{\bm{n},\bm{m}+\bm{e}_k}(w)w^{-n_k}]}{L_k[\Phi_{\bm{n},\bm{m}}(w)w^{-n_k}]}
      =
        \frac{L_k[\Phi_{\bm{n}+\bm{e}_k,\bm{m}}^*(w)w^{m_k}]}{L_k[\Phi_{\bm{n},\bm{m}}^*(w)w^{m_k}]} = 1-\alpha_{\bm{n};\bm{m}+\bm{e}_k}\beta_{\bm{n}+\bm{e}_k;\bm{m}}.
    \end{equation}
    % \begin{align}\label{eq:another one minus alphabeta}
    %     \frac{L_k[\Phi_{\bm{n}+\bm{e}_k,\bm{m}}^*(w)w^{m_k}]}{L_k[\Phi_{\bm{n},\bm{m}}^*(w)w^{m_k}]} &= 1-\alpha_{\bm{n};\bm{m}+\bm{e}_k}\beta_{\bm{n}+\bm{e}_k;\bm{m}},
    %     \\
    %     \label{eq:one minus alphabeta}
    %     \frac{L_k[\Phi_{\bm{n},\bm{m}+\bm{e}_k}(w)w^{-n_k}]}{L_k[\Phi_{\bm{n},\bm{m}}(w)w^{-n_k}]} &= 1-\alpha_{\bm{n};\bm{m}+\bm{e}_k}\beta_{\bm{n}+\bm{e}_k;\bm{m}}.
    % \end{align}
    %Also, $\alpha_{\bm{n};\bm{m}+\bm{e}_k}\beta_{\bm{n}+\bm{e}_k;\bm{m}} \neq 1$ if $(\bm{n}+\bm{e}_k;\bm{m}+\bm{e}_k)$ is normal.
    Furthermore, 
    %that $(\bm{n};\bm{m})$, $(\bm{n}+\bm{e}_k;\bm{m})$, and $(\bm{n};\bm{m}+\bm{e}_k)$ is normal, then 
    $\alpha_{\bm{n};\bm{m}+\bm{e}_k}\beta_{\bm{n}+\bm{e}_k;\bm{m}} \neq 1$ if and only if $(\bm{n}+\bm{e}_k;\bm{m}+\bm{e}_k)$ is normal.
\end{cor}
\begin{rem}
    % If each functional is given by integration with respect to a unit circle measure, then the condition $\alpha_{\bm{n};\bm{m}+\bm{e}_k}\beta_{\bm{n}+\bm{e}_k;\bm{m}} \neq 1$ turns into $\abs{\alpha_{\bm{n};\bm{m}+\bm{e}_k}} \neq 1$.
    If each functional is given by integration with respect to a probability measure on $\T$, then this together with Proposition~\ref{thm:reversal vs star polynomials} shows that $\abs{\alpha_{\bm{n};\bm{n}+\bm{e}_k}} \neq 1$ and $\abs{\beta_{\bm{n}+\bm{e}_k;\bm{n}}} \neq 1$.
\end{rem}

\begin{proof}
    Multiply both sides of \eqref{eq:SzegoII7} by $z^{m_k}$, and multiply \eqref{eq:SzegoII8} by $z^{-n_k}$, then apply $L_k$ to get \eqref{eq:one minus alphabeta}. If these expressions vanish then $\Phi_{\bm{n},\bm{m}+\bm{e}_k}$ satisfies all the orthogonality relations of $(\bm{n}+\bm{e}_k;\bm{m}+\bm{e}_k)$, which contradicts normality. 
\end{proof}

%\begin{prop}
%    \begin{align}\label{eq:SzegoII9}
%        & \alpha_{\bm{n};\bm{m}}z\Phi_{\bm{n};\bm{m}}(z) = \alpha_{\bm{n};\bm{m}}\Phi_{\bm{n}+\bm{e}_k;\bm{m}}(z) - \alpha_{\bm{n}+\bm{e}_k;\bm{m}}\Phi_{\bm{n};\bm{m}}(z) + \sum_{j = 1}^r \alpha_{\bm{n}+\bm{e}_j;\bm{m}}\rho_{\bm{n};\bm{m},j}z\Phi_{\bm{n}-\bm{e}_j;\bm{m}}(z), \\
%        & \label{eq:SzegoII10}\beta_{\bm{n};\bm{m}}\Phi_{\bm{n};\bm{m}}^*(z) = \beta_{\bm{n};\bm{m}}\Phi_{\bm{n}+\bm{e}_k;\bm{m}}^*(z) - \beta_{\bm{n}+\bm{e}_k;\bm{m}}z\Phi_{\bm{n};\bm{m}}^*(z) + \sum_{j = 1}^r\beta_{\bm{n}+\bm{e}_k;\bm{m}}\rho_{\bm{n};\bm{m},j}z\Phi_{\bm{n}-\bm{e}_j;\bm{m}}^*(z),
%    \end{align}
%    and \red{type I analogues}
%\end{prop}
%
%\begin{proof}
%    
%\end{proof}

\begin{prop} Assuming normality of all indices that appear in the respective equations below, %along with their neighbouring indices, 
we have the recurrence relations
    \begin{align}
        \label{eq:SzegoII9}
         z&\Phi_{\bm{n};\bm{m}}(z) = \Phi_{\bm{n}+\bm{e}_k;\bm{m}}(z) - \alpha_{\bm{n}+\bm{e}_k;\bm{m}}\Phi_{\bm{n};\bm{m}}^*(z) + \sum_{j = 1}^r \rho_{\bm{n};\bm{m},j}\gamma_{\bm{n};\bm{m}}^{jk}z\Phi_{\bm{n}-\bm{e}_j;\bm{m}}(z), \\
        \label{eq:SzegoII10}
         z^{-1}&\Phi_{\bm{n};\bm{m}}^*(z) = \Phi_{\bm{n};\bm{m}+\bm{e}_k}^*(z) - \beta_{\bm{n};\bm{m}+\bm{e}_k}\Phi_{\bm{n};\bm{m}}(z) + \sum_{j = 1}^r \sigma_{\bm{n};\bm{m},j}\eta_{\bm{n};\bm{m}}^{jk}z^{-1}\Phi_{\bm{n};\bm{m}-\bm{e}_j}(z),
     %z^{-1}&\Phi_{\bm{n};\bm{m}}^*(z) = \Phi_{\bm{n};\bm{m}+\bm{e}_k}^*(z) - \beta_{\bm{n};\bm{m}+\bm{e}_k}\Phi_{\bm{n};\bm{m}}(z) + \sum_{j = 1}^r \sigma_{\bm{n};\bm{m},j}\eta_{\bm{n};\bm{m}}^{jk}z^{-1}\Phi_{\bm{n};\bm{m}-\bm{e}_j}(z),
    \end{align}
    and 
    \begin{align*}
        z &\bm{\Xi}_{\bm{n};\bm{m}}(z) = \bm{\Xi}_{\bm{n}-\bm{e}_k;\bm{m}}(z) + \beta_{\bm{n}-\bm{e}_k;\bm{m}}\bm{\Xi}_{\bm{n};\bm{m}}^*(z) + \sum_{j=1}^r \rho_{\bm{n};\bm{m},j}\gamma_{\bm{n}-\bm{e}_k+\bm{e}_j;\bm{m}}^{jk}z\bm{\Xi}_{\bm{n}+\bm{e}_j;\bm{m}}(z),\\
        z^{-1}&\bm{\Xi}_{\bm{n};\bm{m}}^*(z) = \bm{\Xi}_{\bm{n};\bm{m}-\bm{e}_k}^*(z) + \alpha_{\bm{n};\bm{m}-\bm{e}_k}\bm{\Xi}_{\bm{n};\bm{m}}(z)+\sum_{j=1}^r \sigma_{\bm{n};\bm{m},j}\eta_{\bm{n};\bm{m}-\bm{e}_k+\bm{e}_j}^{jk}z^{-1}\bm{\Xi}_{\bm{n};\bm{m}+\bm{e}_j}(z)
    \end{align*}
\end{prop}

\begin{proof}
    By \eqref{eq:SzegoII2} and \eqref{eq:CCII} we have
    \begin{align*}
        & \Phi_{\bm{n}+\bm{e}_k;\bm{m}}(z) = \alpha_{\bm{n}+\bm{e}_k;\bm{m}}\Phi_{\bm{n}+\bm{e}_k;\bm{m}}^*(z) + \sum_{j = 1}^r \rho_{\bm{n}+\bm{e}_k;\bm{m},j}z\Phi_{\bm{n}+\bm{e}_k-\bm{e}_j;\bm{m}}(z) \\ & 
        = \alpha_{\bm{n}+\bm{e}_k;\bm{m}}\Phi_{\bm{n}+\bm{e}_k;\bm{m}}^*(z) + \sum_{j \neq k}\rho_{\bm{n}+\bm{e}_k;\bm{m},j}z(\Phi_{\bm{n};\bm{m}}(z) + \gamma_{\bm{n}-\bm{e}_j;\bm{m}}^{kj}\Phi_{\bm{n}-\bm{e}_j;\bm{m}}(z)) \\ & \hspace{5cm}+ \rho_{\bm{n}+\bm{e}_k;\bm{m},k}z\Phi_{\bm{n};\bm{m}}(z).
    \end{align*}
    By separating the two terms in the sum and then using \eqref{eq:compatibility condition 2} and \eqref{eq:compatibility condition 3.1} we end up with
    \begin{multline*}
        \Phi_{\bm{n}+\bm{e}_k;\bm{m}}(z) = \alpha_{\bm{n}+\bm{e}_k;\bm{m}}\Phi_{\bm{n}+\bm{e}_k;\bm{m}}^*(z) + (1 - \alpha_{\bm{n}+\bm{e}_k;\bm{m}}\beta_{\bm{n}+\bm{e}_k;\bm{m}})z\Phi_{\bm{n};\bm{m}}(z) \\ - \sum_{j = 1}^r\rho_{\bm{n};\bm{m},j}\gamma_{\bm{n};\bm{m}}^{jk}z\Phi_{\bm{n}-\bm{e}_j;\bm{m}}(z).
    \end{multline*}
    This together with \eqref{eq:SzegoII1} implies the first equality. %rom here \eqref{eq:SzegoII9} follows by . \eqref{eq:SzegoII10} 
    The second equality then follows by Proposition \ref{thm:reversal vs star polynomials}. The type I formulas can be derived analogously. 
\end{proof}

\begin{rem}
    The recurrence relations can alternatively be derived directly from orthogonality relations, e.g., the first relation follows by checking orthogonality relations and degree restrictions of $\Phi_{\bm{n};\bm{m}}(z) - z^{-1}\Phi_{\bm{n}+\bm{e}_k;\bm{m}}+\alpha_{\bm{n}+\bm{e}_k;\bm{m}}z^{-1}\Phi_{\bm{n};\bm{m}}^*(z)$ (using similar arguments as in the proofs of the Szeg\H{o} recurrence relations of \cite{KVMOPUC}), and then applying \eqref{eq:rho} and \eqref{eq:gamma}. This approach does not need normality of indices of the form $(\bm{n}+\bm{e}_k-\bm{e}_j;\bm{m})$, $j \neq k$. 
\end{rem}

We can eliminate the coefficients $\gamma_{\bm{n};\bm{m}}^{jk}$ in \eqref{eq:SzegoII9} if we multiply both sides by $\alpha_{\bm{n};\bm{m}}$. By applying \eqref{eq:compatibility condition 1} and then using \eqref{eq:SzegoII2} we get
\begin{multline}\label{eq:generalized three term}
    \alpha_{\bm{n};\bm{m}}z\Phi_{\bm{n};\bm{m}}(z) = \alpha_{\bm{n};\bm{m}}\Phi_{\bm{n}+\bm{e}_k;\bm{m}}(z) - \alpha_{\bm{n}+\bm{e}_k;\bm{m}}\Phi_{\bm{n};\bm{m}}(z) \\ + \sum_{j = 1}^r \alpha_{\bm{n}+\bm{e}_j;\bm{m}}\rho_{\bm{n};\bm{m},j}z\Phi_{\bm{n}-\bm{e}_j;\bm{m}}(z).
    \end{multline}
Hence the end result expresses $z\Phi_{\bm{n};\bm{m}}(z)$ in terms of the nearest neighbours, without appealing to any $\Phi^*$-polynomials or $\beta$-coefficients. In the case $\bm{m} = \bm{0}$ this is the main result of \cite{MOPUC2}. In terms of our recurrence coefficients, in the case $r=1$ this turns into exactly the three-term recurrence relation of OPUC \cite[Sect. 1.5]{OPUC1}. In the degenerate case $\alpha_{\bm{n};\bm{m}} = 0$ we get essentially no information from this relation. 

Of course, an analogous formula can be derived for the type I polynomials, as well as reversed versions through Proposition \ref{thm:reversal vs star polynomials}.
    
% Note that we can get a similar formula to \eqref{eq:generalized three term} by combining \eqref{eq:SzegoII1} and \eqref{eq:SzegoII3} into
%   \begin{multline}
%         \beta_{\bm{n};\bm{m}}\Phi_{\bm{n};\bm{m}}^*(z) = \beta_{\bm{n};\bm{m}}\Phi_{\bm{n}+\bm{e}_k;\bm{m}}^*(z) - \beta_{\bm{n}+\bm{e}_k;\bm{m}}z\Phi_{\bm{n};\bm{m}}^*(z) \\ + \sum_{j = 1}^r \beta_{\bm{n}+\bm{e}_k;\bm{m}}\rho_{\bm{n};\bm{m},j}z\Phi_{\bm{n}-\bm{e}_j;\bm{m}}^*(z)
%     \end{multline}


% Another recurrence relation for $z\Phi_{\bm{n};\bm{m}}(z) - \Phi_{\bm{n}+\bm{e}_k;\bm{m}}(z)$ was obtained in \cite[Cor. 3.4(i)]{KVMOPUC}, also by combining \eqref{eq:SzegoII1} and \eqref{eq:SzegoII3}.  





%%%%%%%%%%%%%%%%%%%%%%%%%%%%%%%%%%%


\section{Heine formulas}\label{ss:Heine}

Assuming $(\bm{n};\bm{m})$ is normal, it is easy to verify the determinantal formula
    \begin{equation}\label{eq:HeineII}
        \Phi_{\bm{n};\bm{m}}(z) = \frac{1}{\det T_{\bm{n};\bm{m}}}
        \det
        %\begin{pmatrix}
          \left(\begin{array}{c c c}
        c_{\abs{\bm{m}}-m_1,1} & \cdots & c_{-\abs{\bm{n}}-m_1,1} \\
        c_{\abs{\bm{m}}-m_1+1,1} & \cdots & c_{-\abs{\bm{n}}-m_1+1,1} \\
        \vdots & \ddots & \vdots \\
        c_{\abs{\bm{m}}+n_1-1,1} & \cdots & c_{-\abs{\bm{n}}+n_1-1,1} 
        % \hline
        % & \vdots & \\
        % \hline
        \\[0.2em]
        %\hline
        \hdashline[0.5pt/2pt]
        \\[-1.2em]
        & \vdots \\[0.2em]
        %\hline
        \hdashline[0.5pt/2pt]
        \\[-1.0em]
        c_{\abs{\bm{m}}-m_r,r} & \cdots & c_{-\abs{\bm{n}}-m_r,r} \\
        c_{\abs{\bm{m}}-m_r+1,r} & \cdots & c_{-\abs{\bm{n}}-m_r+1,r} \\
        \vdots & \ddots & \vdots \\
        c_{\abs{\bm{m}}+n_r-1,r} & \cdots & c_{-\abs{\bm{n}}+n_r-1,r}
        \\[0.2em]
        %\hline
        \hdashline[0.5pt/2pt]
        \\[-1.0em]
        z^{-|\bm{m}|} & \cdots & z^{|\bm{n}|}
    %\end{pmatrix}},
    \end{array}\right),
    \end{equation}
    where $T_{\bm{n};\bm{m}}$ is given in~\eqref{eq:T}. Indeed, it is easy to see that the right-hand side of~\eqref{eq:HeineII} is %indeed 
    of the form~\eqref{eq:moprlIIspan}. Then multiplying it with $z^{-k}$ for a given $k=-m_j,\ldots,n_j-1$ and applying $L_j$, one obtains a determinant with two identical rows, implying~\eqref{eq:moprlII}.

    Similarly,
    \begin{equation}\label{eq:HeineII*}
        \Phi^*_{\bm{n};\bm{m}}(z) = \frac{(-1)^{|\bm{n}|+|\bm{m}|}}{\det T_{\bm{n};\bm{m}}}
        \det
        %\begin{pmatrix}
          \left(\begin{array}{c c c}
        c_{\abs{\bm{m}}-m_1+1,1} & \cdots & c_{-\abs{\bm{n}}-m_1+1,1} \\
        c_{\abs{\bm{m}}-m_1+2,1} & \cdots & c_{-\abs{\bm{n}}-m_1+2,1} \\
        \vdots & \ddots & \vdots \\
        c_{\abs{\bm{m}}+n_1,1} & \cdots & c_{-\abs{\bm{n}}+n_1,1} % \hline
        % & \vdots & \\
        % \hline
        \\[0.2em]
        %\hline
        \hdashline[0.5pt/2pt]
        \\[-1.2em]
        & \vdots \\[0.2em]
        %\hline
        \hdashline[0.5pt/2pt]
        \\[-1.0em]
        c_{\abs{\bm{m}}-m_r+1,r} & \cdots & c_{-\abs{\bm{n}}-m_r+1,r} \\
        c_{\abs{\bm{m}}-m_r+2,r} & \cdots & c_{-\abs{\bm{n}}-m_r+2,r} \\
        \vdots & \ddots & \vdots \\
        c_{\abs{\bm{m}}+n_r,r} & \cdots & c_{-\abs{\bm{n}}+n_r,r}
        \\[0.2em]
        %\hline
        \hdashline[0.5pt/2pt]
        \\[-1.0em]
        z^{-|\bm{m}|} & \cdots & z^{|\bm{n}|}
    %\end{pmatrix}}.
    \end{array}\right).
    \end{equation}
    Analogous formulas can be established for the type I polynomials $\bm{\Xi}_{\bm{m};\bm{n}}$ and $\bm{\Xi}^*_{\bm{m};\bm{n}}$ in the same vein as in~\cite{Kui}.

    This leads to the following determinantal expressions for all the recurrence coefficients from Section~\ref{ss:SzegoNEW} and Section \ref{ss:CC}. In what follows, we denote $wL$ to be linear functional (Christoffel transform) defined by $(wL)[w^k]:=L[w^{k+1}]$, $w\bm{L} = (wL_1,\ldots,wL_r)$, and similarly for $w^{-1}L$ and $w^{-1} \bm{L}$. Finally, we denote $T_{\bm{n};\bm{m}}(\bm{M})$ to be the matrix in~\eqref{eq:T} for any system of linear functionals $\bm{M}$.

    \begin{thm}
    \noindent\hfill
        \begin{itemize}
            \item[(i)] If $(\bm{n};\bm{m})\in\mathfrak{C}_{2r}$ is normal, then
            \begin{alignat}{2}
                \alpha_{\bm{n};\bm{m}} &= (-1)^{|\bm{n}|+|\bm{m}|} \frac{\det T_{\bm{n};\bm{m}}(w\bm{L})}{\det T_{\bm{n};\bm{m}}(\bm{L})},
                \\
                \beta_{\bm{n};\bm{m}} &=(-1)^{|\bm{n}|+|\bm{m}|} \frac{\det T_{\bm{n};\bm{m}}(w^{-1}\bm{L})}{\det T_{\bm{n};\bm{m}}(\bm{L})}.
            \end{alignat}

            \item[(ii)] If $(\bm{n};\bm{m}),(\bm{n}-\bm{e}_j;\bm{m})\in\mathfrak{C}_{2r}$ are normal, then
            \begin{equation}
                \rho_{\bm{n};\bm{m},j} =  \frac{\det T_{\bm{n}+\bm{e}_j;\bm{m}} \det T_{\bm{n}-\bm{e}_j;\bm{m}} }{(\det T_{\bm{n};\bm{m}})^2} ,
            \end{equation}
            and if $(\bm{n};\bm{m}),(\bm{n};\bm{m}-\bm{e}_j)\in\mathfrak{C}_{2r}$ are normal, then
            \begin{equation}
                \sigma_{\bm{n};\bm{m},j} = \frac{\det T_{\bm{n};\bm{m}+\bm{e}_j} \det T_{\bm{n};\bm{m}-\bm{e}_j} }{(\det T_{\bm{n};\bm{m}})^2}.
            \end{equation}
            % \begin{alignat}{2}
            %     \rho_{\bm{n};\bm{m},j} &=  \frac{\det T_{\bm{n}+\bm{e}_j;\bm{m}} \det T_{\bm{n}-\bm{e}_j;\bm{m}} }{(\det T_{\bm{n};\bm{m}})^2} ,
            %     \\
            %     \sigma_{\bm{n};\bm{m},j} &= \frac{\det T_{\bm{n};\bm{m}+\bm{e}_j} \det T_{\bm{n};\bm{m}-\bm{e}_j} }{(\det T_{\bm{n};\bm{m}})^2} .
            % \end{alignat}

            \item[(iii)] If $(\bm{n};\bm{m}),(\bm{n}+\bm{e}_k;\bm{m}),(\bm{n}+\bm{e}_\ell;\bm{m})\in\mathfrak{C}_{2r}$  are normal, and $k < \ell$, then
            \begin{equation}
                \gamma^{k\ell}_{\bm{n};\bm{m}} = \frac{\det T_{\bm{n}+\bm{e}_\ell+\bm{e}_k;\bm{m}} \det T_{\bm{n};\bm{m}} }{\det T_{\bm{n}+\bm{e}_\ell;\bm{m}} \det T_{\bm{n}+\bm{e}_k;\bm{m}}} ,
            \end{equation}
            and if $(\bm{n};\bm{m}),(\bm{n};\bm{m}+\bm{e}_k),(\bm{n};\bm{m}+\bm{e}_\ell)\in\mathfrak{C}_{2r}$  are normal, and $k < \ell$, then
            \begin{equation}
                \eta^{k\ell}_{\bm{n};\bm{m}} = \frac{\det T_{\bm{n};\bm{m}+\bm{e}_\ell+\bm{e}_k} \det T_{\bm{n};\bm{m}} }{\det T_{\bm{n};\bm{m}+\bm{e}_\ell} \det T_{\bm{n};\bm{m}+\bm{e}_k}} .
            \end{equation}
            \item[(iv)] If $(\bm{n};\bm{m}),(\bm{n}+\bm{e}_k;\bm{m}),(\bm{n};\bm{m}+\bm{e}_k)\in\mathfrak{C}_{2r}$  are normal
            \begin{equation}
                1-\alpha_{\bm{n};\bm{m}+\bm{e}_k}\beta_{\bm{n}+\bm{e}_k;\bm{m}} = \frac{\det{T_{\bm{n}+\bm{e}_k;\bm{m}+\bm{e}_k}}\det{T_{\bm{n};\bm{m}}}}{\det{T_{\bm{n}+\bm{e}_k;\bm{m}}}\det{T_{\bm{n};\bm{m}+\bm{e}_k}}}.
            \end{equation}
            
            % \begin{alignat}{2}
            %     \gamma^{k\ell}_{\bm{n};\bm{m}} &= \frac{\det T_{\bm{n}+\bm{e}_\ell+\bm{e}_k;\bm{m}} \det T_{\bm{n};\bm{m}} }{\det T_{\bm{n}+\bm{e}_\ell;\bm{m}} \det T_{\bm{n}+\bm{e}_k;\bm{m}}} ,
            %     \\
            %     \eta^{k\ell}_{\bm{n};\bm{m}} &= \frac{\det T_{\bm{n};\bm{m}+\bm{e}_\ell+\bm{e}_k} \det T_{\bm{n};\bm{m}} }{\det T_{\bm{n};\bm{m}+\bm{e}_\ell} \det T_{\bm{n};\bm{m}+\bm{e}_k}} .
            % \end{alignat}
        \end{itemize}
    \end{thm}
\begin{proof}
    For (i) and (ii), we simply need to find the $z^{-|\bm{m}|}$ coefficient of~\eqref{eq:HeineII} and $z^{|\bm{n}|}$ coefficient of~\eqref{eq:HeineII*}.
    
    Now note that from~\eqref{eq:HeineII} and~\eqref{eq:HeineII*} we obtain
    \begin{align}
    \label{eq:norm1}
        L_j[\Phi_{\bm{n};\bm{m}}(w)w^{-n_j}]
        & = (-1)^{\sum_{s=j+1}^r (n_s+m_s)} \frac{\det T_{\bm{n}+\bm{e}_j;\bm{m}}}{\det T_{\bm{n};\bm{m}}},
        \\
        \label{eq:norm2}
        L_j[\Phi^*_{\bm{n};\bm{m}}(w)w^{m_j}]
        & =  (-1)^{\sum_{s=1}^{j-1} (n_s+m_s)} \frac{\det T_{\bm{n};\bm{m}+\bm{e}_j}}{\det T_{\bm{n};\bm{m}}}.
    \end{align}
    Then (ii) follows by combining this with~\eqref{eq:rho} and~\eqref{eq:sigma}, (iii) follows by combining with \eqref{eq:gamma} and \eqref{eq:eta}, and (iv) follows by combining with \eqref{eq:one minus alphabeta}.
    %For (iii), observe that from~\eqref{eq:CCII} and ~\eqref{eq:CCII2} we get
    %\begin{align}\label{eq:gamma}
    %\gamma^{kl}_{\bm{n};\bm{m}} &= \frac{L_l[\Phi_{\bm{n}+\bm{e}_k;\bm{m}}(w)w^{-n_l}]}{L_l[\Phi_{\bm{n};\bm{m}}(w)w^{-n_l}]},
    %\\
    %\eta^{kl}_{\bm{n};\bm{m}} &= \frac{L_l[\Phi^*_{\bm{n};\bm{m}+\bm{e}_k}(w)w^{m_l}]}{L_l[\Phi^*_{\bm{n};\bm{m}}(w)w^{m_l}]},
    %\end{align}
    %which, together with~\eqref{eq:norm1}--\eqref{eq:norm2}, imply the statement.
\end{proof}



%%%%%%%%%%%%%%%%%%%%%%%%%%%%%%%%%%%%%

\section{The Christoffel--Darboux formula}\label{ss:CD}

We can establish a generalized Christoffel--Darboux formula by following the same argument as in~\cite[Thm 6.1]{KVMOPUC}. In this formulation, the paths may originate from points other than $\bm{0}$, and the resulting expression takes a more streamlined form due to our choice of definition of $\bm\Xi_{\bm{n};\bm{m}}$, as discussed in Remarks~\ref{rem:annoying1}, \ref{rem:annoying2}.
\begin{thm} 
    Let $\bm{m}\in\N^r$ be fixed, and consider $(\bm{n}_k;\bm{m})_{k = 0}^{N}$ to be a path of $\mathfrak{C}_{2r}$-multi-indices, %with $\bm{n}_k+\bm{m} \ge \bm{0}$,
    such that $\bm{n}_0 = \bm{-\bm{m}}$, 
    %$|\bm{n}_k|=k$,
    %, $\bm{n}_{N} = \bm{n}$, 
    and $\bm{n}_{k+1}-\bm{n}_k = \bm{e}_{l_k}$ for some $1\le l_k\le r$. %, in particular $|\bm{n}_k|=k$.
    Assume all multi-indices on the path are normal, along with all the neighbouring indices that belong to $\mathfrak{C}_{2r}$. Then we have the Christoffel--Darboux formulas
    \begin{equation}\label{eq:cd formula1}
        \begin{aligned}
    (\xi - z)& \sum_{k = 0}^{N-1} \Phi_{\bm{n}_k;\bm{m}}(z)  \bm\Xi_{\bm{n}_{k+1};\bm{m}}(\xi)
    \\
    &= \Phi^*_{\bm{n}_N;\bm{m}}(z) \bm\Xi^*_{\bm{n}_N;\bm{m}}(\xi) - z\xi \sum_{j = 1}^{r}\rho_{\bm{n}_N;\bm{m},j}\Phi_{\bm{n}_N-\bm{e}_j;\bm{m}}(z) \bm\Xi_{\bm{n}_N+\bm{e}_j;\bm{m}}(\xi)
    \end{aligned}
    \end{equation}
    and
    \begin{equation}\label{eq:cd formula2}
    \begin{aligned}
    (z-\xi)& \sum_{k = 0}^{N-1} \Phi^*_{\bm{m};\bm{n}_k}(z)  \bm\Xi^*_{\bm{m};\bm{n}_{k+1}}(\xi) 
    \\
    &= z\xi \Phi_{\bm{m};\bm{n}_N}(z)\bm\Xi_{\bm{m};\bm{n}_N}(\xi) -  \sum_{j = 1}^{r}\sigma_{\bm{m};\bm{n}_N,j}\Phi^*_{\bm{m};\bm{n}_N-\bm{e}_j}(z) \bm\Xi^*_{\bm{m};\bm{n}_N+\bm{e}_j}(\xi).
    \end{aligned}
    \end{equation}
    \begin{proof}
        In the expression
        $$
         \Phi_{\bm{n}_k;\bm{m}}(z)  \bm\Xi_{\bm{n}_{k+1};\bm{m}}(\xi) - z\xi^{-1} \Phi_{\bm{n}_k;\bm{m}}(z)  \bm\Xi_{\bm{n}_{k+1};\bm{m}}(\xi)
        $$
        use the Szeg\H{o} recurrence~\eqref{eq:SzegoII2} followed by~\eqref{eq:SzegoI1} on the first term and~\eqref{eq:SzegoI2} followed by~\eqref{eq:SzegoII1}. Then the arguments from the proof ~\cite[Thm 6.1]{KVMOPUC} go through in a straightforward manner.
        
        The second formula~\eqref{eq:cd formula2} follows from~\eqref{eq:cd formula1} after applying Proposition~\ref{thm:reversal vs star polynomials} and Remark~\ref{rem:reversal}.
    \end{proof}
\end{thm}


% \begin{thm} (WRONG)
%     Let $\bm{m}\in\N^r$ be fixed, and consider $(\bm{n}_k;\bm{m})_{k = 0}^{N}$ to be a path of multi-indices such that $\bm{n}_0 = \bm{0}$, 
%     %$|\bm{n}_k|=k$,
%     %, $\bm{n}_{N} = \bm{n}$, 
%     and $\bm{n}_{k+1}-\bm{n}_k = \bm{e}_{l_k}$ for some $1\le l_k\le r$, in particular $|\bm{n}_k|=k$. Assume all multi-indices on the path are normal, along with all the neighbouring indices that belong to $\N^r\times \N^r$. Then we have the Christoffel--Darboux formulas
%     \begin{align}\label{eq:cd formula1}
%     (\xi - z)& \sum_{k = 0}^{N-1} \Phi_{\bm{n}_k;\bm{m}}(z)  \bm\Xi_{\bm{n}_{k+1};\bm{m}}(\xi)
%     \\
%     &= \Phi^*_{\bm{n}_N;\bm{m}}(z) \bm\Xi^*_{\bm{n}_N;\bm{m}}(\xi) - z\xi \sum_{j = 1}^{r}\rho_{\bm{n}_N;\bm{m},j}\Phi_{\bm{n}_N-\bm{e}_j;\bm{m}}(z) \bm\Xi_{\bm{n}_N+\bm{e}_j;\bm{m}}(\xi)
%     \end{align}
%     and
%     \begin{align}\label{eq:cd formula2}
%     (z-\xi)& \sum_{k = 0}^{N-1} \Phi^*_{\bm{m};\bm{n}_k}(z)  \bm\Xi^*_{\bm{m};\bm{n}_{k+1}}(\xi) 
%     \\
%     &= z\xi \Phi_{\bm{m};\bm{n}_N}(z)\bm\Xi_{\bm{m};\bm{n}_N}(\xi) -  \sum_{j = 1}^{r}\sigma_{\bm{m};\bm{n}_N,j}\Phi^*_{\bm{m};\bm{n}_N-\bm{e}_j}(z) \bm\Xi^*_{\bm{m};\bm{n}_N+\bm{e}_j}(\xi).
%     \end{align}
%     \begin{proof}
%         In the expression
%         $$
%          \Phi_{\bm{n}_k;\bm{m}}(z)  \bm\Xi_{\bm{n}_{k+1};\bm{m}}(\xi) - z\xi^{-1} \Phi_{\bm{n}_k;\bm{m}}(z)  \bm\Xi_{\bm{n}_{k+1};\bm{m}}(\xi)
%         $$
%         use the Szeg\H{o} recurrence~\eqref{eq:SzegoII2} followed by~\eqref{eq:SzegoI1} on the first term and~\eqref{eq:SzegoI2} followed by~\eqref{eq:SzegoII1}. Then the arguments from the proof ~\cite[Thm 6.1]{KVMOPUC} go through in a straightforward manner.

%         NO! The telescoping will not disappear at $\bm{n}_0=\bm{0}$! So there will be extra terms of the form
%         $$
%         -\xi^{-1} \Phi^*_{\bm{0};\bm{m}}(z)\bm\Xi^*_{\bm{0};\bm{m}}+ \sum_{j=1}
%         ^r\rho_{\bm{0};\bm{m},j} z \Phi_{-\bm{e}_j;\bm{m}}\bm\Xi_{\bm{e}_jm}(\xi)
%         $$
        
%         The second formula~\eqref{eq:cd formula2} follows from~\eqref{eq:cd formula1} after applying Theorem~\ref{thm:reversal vs star polynomials} and Remark~\ref{rem:reversal}.
%     \end{proof}
% \end{thm}
% % \begin{rem}
% %     If $r=1$, then (recall Remark~\ref{rem:r=1}), then after some algebraic manipulation, one obtains that formulas~\eqref{eq:cd formula1} and~\eqref{eq:cd formula2} reduce to
% %     \begin{equation}\label{eq:CDr=1}
% %         \sum_{k=0}^{N-1} \varphi_{k+m}(z) \overline{\varphi_{k+m}(\zeta)}
% %         = \frac{1}{1-z\bar\zeta} ...
% %     \end{equation}
% % \end{rem}

%%%%%%%%%%%%%%%%%%%%%%%%%%%%%%%%%%%%%%%%%%%%%%%%%%%%%%%%%%%%%%%%%%%%%%%%%%%%%

\section{Szeg\H{o} Mapping}\label{ss:SzegoMap}

Let us recall the main notions from the theory of orthogonal and multiple orthogonal polynomials on the real line with respect to moment functionals acting on the space of complex polynomials 
(with nonnegative powers, in contrast to the Laurent setting~\eqref{eq:moments linear functional}). To this end, let $M$ be defined by
\begin{equation}\label{eq:moments linear functional real}
         M[x^{k}] = m_{k} , \qquad k\in\N, %\qquad j = 1,\dots,r,
\end{equation}
where $m_{k}$ are arbitrary complex numbers. 

Often $M$ is induced by a probability measure $\gamma$ on the real line $\R$ with all finite moments
\begin{equation}\label{eq:functionals given by measures real}
        M[x^{k}] = \int_{\R} x^{k} d\gamma(x), \qquad 
        k \in \N. % \qquad j = 1,\dots,r, 
\end{equation}
%For brevity, we will refer to  $M$ as a functional on $\R$, and to $L$ in~\eqref{eq:moments linear functional} as a functional on $\T$.

Orthogonal polynomials with respect to $M$ are defined by $\deg P_n = n$ and 
%Recall~\cite{Chihara} that a moment functional $M$ is called quasi-definite if monic orthogonal polynomial $P_n$, $\deg P_n = n$,
$$
M[P_n(x) x^k] = 0, \qquad k=0,1,\ldots,n-1,
$$
%exists and is unique for every $n\in\N$. 
Recall~\cite{Chihara} that a moment functional $M$ is called quasi-definite if there is a unique monic $P_n$ and $\deg{P_n} = n$, for each positive integer $n$. This is equivalent to all Hankel matrices $(m_{j+\ell})_{j,\ell=0}^{n-1}$ being invertible. This is automatic if $M$ is associated with a positive measure on $\R$ with finite moments and infinite support, as in~\eqref{eq:functionals given by measures real}.

Similarly, a moment functional $L$ on the space of Laurent polynomials~\eqref{eq:moments linear functional} is called quasi-definite if a monic polynomial $\Phi_n$ with $\deg \Phi_n = n$ and
$$
L[\Phi_n(z) z^{-k}] = 0, \qquad k=0,1,\ldots,n-1,
$$ 
exists and is unique for every $n\in\N$. This is equivalent to all Toeplitz matrices $(c_{j-\ell})_{j,\ell=0}^{n-1}$ being invertible and is automatic if $L$ is associated with a positive measure with infinite support on $\T$ as in~\eqref{eq:functionals given by measures}.

Given a system $\bm{M} = (M_1,\dots,M_r)$ of moment functionals~\eqref{eq:moments linear functional real}, the type II multiple orthogonal polynomials $P_{\bm{n}}$ with respect to the multi-index $\bm{n}$ and the system $\bm{M}$ are non-zero polynomials of at most degree $\abs{\bm{n}}$, satisfying the orthogonality relations
\begin{equation}\label{eq:MOPRLIIreal}
    M_j[P_{\bm{n}}(x)x^k] = 0, \qquad k = 0,\dots,n_j-1, \qquad j = 1,\dots,r.
\end{equation}
Such polynomials always exists, and we say that the index $\bm{n}$ is normal if there exists a unique $P_{\bm{n}}$ with $x^{\abs{\bm{n}}}$-coefficient equal to $1$. %, and $\bm{\gamma}$ is perfect if every $\bm{n} \in \N^r$ is normal. 

The type I polynomials are non-zero vectors $\bm{A}_{\bm{n}} = (A_{\bm{n},1},\dots,A_{\bm{n},r})$, except for $\bm{A}_{\bm{0}} = \bm{0}$, where $A_{\bm{n},j}$ are polynomials of degree at most $n_j-1$ for each $j = 1,\dots,r$, and
\begin{equation}\label{eq:MOPRLIreal}
    \sum_{j = 1}^r M_j[A_{\bm{n},j}(x)x^k ] = 
        0, \qquad k = 0,\dots,\abs{\bm{n}}-2.
\end{equation}
$\bm{n} \neq \bm{0}$ is normal if and only if there exists a unique $\bm{A}_{\bm{n}}$ with
\begin{equation}\label{eq:normalization type I moprl}
    \sum_{j = 1}^r M_j[A_{\bm{n},j}(x) x^{\abs{\bm{n}}-1}] = 1. 
\end{equation}
When $\bm{n}$ is normal we always work with the above normalizations for the type I and type II polynomials. %See~\cite{Aptekarev, Ismail, Nikishin} for more about multiple orthogonal polynomials on the real line. 

Assuming all the indices are normal (we then say that the system is perfect), there exist~\cite{NNRR,Ismail} coefficients $a_{\bm{n},j}, b_{\bm{n},j}$, called the nearest neighbour recurrence coefficients, so that 
\begin{equation}\label{eq:one nnr}
xP_{\bm{n}}(x) = P_{\bm{n} + \bm{e}_k}(x) + b_{\bm{n},k}P_{\bm{n}}(x) + \sum_{j = 1}^{r}a_{\bm{n},j}P_{\bm{n} - \bm{e}_j}(x),
\end{equation}
as well as
\begin{equation}\label{eq:type 1 recurrence relation}
x\bm{A}_{\bm{n}}(x) = \bm{A}_{\bm{n} - \bm{e}_k}(x) + b_{\bm{n}-\bm{e}_k,k}\bm{A}_{\bm{n}}(x) + \sum_{j = 1}^{r}a_{\bm{n},j}\bm{A}_{\bm{n} + \bm{e}_j}(x).
\end{equation}


% Given any probability measure $\gamma$ on $[-2,2]$, we define its image under the Szeg\H{o} mapping $\mu=\operatorname{Sz}(\gamma)$ to be the probability measure on $\T$ that is invariant under the reflection $e^{i\theta}\mapsto e^{-i\theta}$ and satisfies
% \begin{equation}\label{eq:szego map}
%     \int_{\T} g(2\cos\theta) d\mu(e^{i\theta}) = \int_{-2}^2 g(x) d\gamma(x),
% \end{equation}
% for all measurable functions $g$ on $[-2,2]$. This is called the Szeg\H{o} mapping. Conversely,  any probability measure $\mu$ on $\T$ that is invariant under $e^{i\theta}\mapsto e^{-i\theta}$ arises in this way, i.e.,  $\mu=\operatorname{Sz}(\gamma)$ for a uniquely determined probability measure $\gamma$ on $[-2,2]$  determined by~\eqref{eq:szego map}.

Recall the definition of the Szeg\H{o} mapping in Section~\ref{ss:intro3} for the case of measures. It is easy to extend this construction to linear functionals. Indeed, the equalities
\begin{equation}
    \label{eq:SzegoMapFunctionals}
    L[(w+w^{-1})^k] = M[x^k], \qquad k\in\N
\end{equation}
set up one-to-one correspondence between all linear functionals $M$ on the space of polynomials (i.e.,~\eqref{eq:moments linear functional real}), and all linear functionals $L$ on the space of Laurent polynomials (i.e., ~\eqref{eq:moments linear functional}) with the symmetry
\begin{equation}\label{eq:symmetricFunctionals}
    L[w^k]=L[w^{-k}], \qquad k\in\N.
\end{equation}
We denote this correspondence by $L = \operatorname{Sz}(M)$ and its inverse by $M = \operatorname{Sz}^{-1}(L)$.



% \begin{prop}
%     Let $L=\operatorname{Sz}(M)$. If $L$ is quasi-definite then $M$ is quasi-definite, in which case
%     \begin{align}
%     P_n(z+z^{-1}) & = \frac{1}{1+\alpha_{2n}} \big( z^{-n} \Phi_{2n}(z) + z^{-n} \Phi_{2n}(1/z) \big) \\
%     &  =  \big( z^{-n+1} \Phi_{2n-1}(z) + z^{n-1} \Phi_{2n-1}(1/z) \big) 
%     \end{align}
% \end{prop}

%Assume that $\bm{L} = (L_1,\dots,L_r)$ is a system of linear functionals on $\T$ satisfying~\eqref{eq:symmetricFunctionals} and let $\bm{M}$ be defined by $M_j=\operatorname{Sz}^{-1}(L_j)$, $j = 1,\dots,r$.
%Assume that $\bm{M}$ is a system of linear functionals on $\R$ and let $\bm{L} = (L_1,\dots,L_r) = \operatorname{Sz}(\bm{M})$ be defined by $L_j=\operatorname{Sz}(M_j)$, $j = 1,\dots,r$. %Let us also assume that $\bm{\mu}$ is perfect on $\T$ (which is not automatically true). 


Functionals on Laurent polynomials with the symmetry~\eqref{eq:symmetricFunctionals} exhibit a number of properties similar to those that we observed in Proposition~\ref{thm:reversal vs star polynomials} and Remark~\ref{rem:reversal}. 

\begin{prop}\label{rem:Rsymmetry}
    %Given moment functionals~\eqref{eq:moments linear functional}, define $\widetilde{L}_j$ to be $\widetilde{L}_j[w^{-k}]=L_j[w^k]$ and $\widetilde{\bm{L}}=(\widetilde{L}_1,\ldots,\widetilde{L}_r)$. 
    Let $L_1,\ldots,L_r$ be linear functionals on the space of Laurent polynomials, each satisfying the symmetry condition~\eqref{eq:symmetricFunctionals}.
    
    Then $(\bm{n};\bm{m})$ is normal if and only if $(\bm{m};\bm{n})$ is normal, and  
    \begin{alignat}{1}
        %\label{eq:symmReal}
        % & \det T_{\bm{n};\bm{m}} = \pm \det T_{\bm{m};\bm{n}}, 
        % \\        
        & \Phi_{\bm{n};\bm{m}}^*(z)=\Phi_{\bm{m};\bm{n}}(1/z),
        \qquad
        \bm{\Xi}^*_{\bm{n};\bm{m}}(z) = \bm{\Xi}_{\bm{m};\bm{n}}(1/z), 
        \\
        & \label{eq:szego symmetry}
        \alpha_{\bm{n};\bm{m}} = {\beta}_{\bm{m};\bm{n}},
        \qquad
        \rho_{\bm{n};\bm{m},j}= \sigma_{\bm{m};\bm{n},j},
        \qquad
        \gamma^{kl}_{\bm{n};\bm{m},j}= \eta^{kl}_{\bm{m};\bm{n},j},
    \end{alignat}
    hold.
\end{prop}

In particular, in the special case when each $L_j$ corresponds to integration with respect to a positive measure on the unit circle~\eqref{eq:functionals given by measures} that is invariant under $e^{i\theta}\mapsto e^{-i\theta}$, then Propositions~\ref{thm:reversal vs star polynomials} and~\ref{rem:Rsymmetry} together imply that $\alpha_{\bm{n};\bm{m}} = \beta_{\bm{m};\bm{n}}$ and $\rho_{\bm{n};\bm{m}} = \sigma_{\bm{m};\bm{n}}$ are real.
Now we are ready to establish a relationship between multiple orthogonality on the real line and multiple orthogonality on the unit circle. For $r=1$ this relationship goes back to Szeg\H{o} \cite{Szego} (for functionals, it is effectively in~\cite{GeronimusRelations,GeronimusBook}, see also~\cite[Sect. 13]{OPUC2} and~\cite{DerSim18}). %We follow the elementary proof from ~\cite[Thm 13.1.5]{OPUC2}.

\begin{thm}\label{thm:szego mapping type II}
    Assume that $\bm{L} = (L_1,\dots,L_r)$ is a system of linear functionals on $\T$ satisfying~\eqref{eq:symmetricFunctionals} and let $\bm{M}$ be defined by $M_j=\operatorname{Sz}^{-1}(L_j)$, $j = 1,\dots,r$.

    %Assume $\bm{L} = \operatorname{Sz}(\bm{M})$ and that 
    If $(\bm{n};\bm{n})$ and $(\bm{n}+\bm{e}_j;\bm{n})$ are normal for $\bm{L}$ for all $\bm{n}\in\N^r$ and $j=1,\ldots,r$, then $\bm{M}$ is perfect (that is, every $\bm{n}\in\N^r$ is normal for $\bm{M}$), and
    \begin{align}\label{eq:SzegoPolynomialRelation}
          P_{\bm{n}}(z+z^{-1}) & = \frac{1}{(1+\alpha_{\bm{n};\bm{n}})} \Big( \Phi_{\bm{n};\bm{n}}(z) + \Phi_{\bm{n};\bm{n}}(1/z) \Big) 
          \\
          \label{eq:SzegoPolynomialRelation2}
          %& = \Phi_{\bm{n};\bm{m}}(z) + \Phi_{\bm{n};\bm{m}}(1/z)
          & = \Phi_{\bm{n};\bm{n}-\bm{e}_j}(z) + \Phi_{\bm{n};\bm{n}-\bm{e}_j}(1/z).
        \end{align}
        %for any $\bm{m}$ such that $\bm{n} - \bm{1} \leq \bm{m} \lneq \bm{n}$.
        Furthermore, $\alpha_{\bm{n};\bm{n}}$ is necessarily $\ne -1$. 
    %where $\alpha_{\bm{n};\bm{n}}$ is the $z^{-\abs{\bm{n}}}$-coefficient of $\Phi_{\bm{n};\bm{n}}$.
\end{thm}
% \begin{rem}
%     As is clear from the proof, the index $\bm{n}-\bm{e}_j$ can be replaced with any index $\bm{m}$ such that  $\bm{n} - \bm{1} \leq \bm{m} \lneq \bm{n}$.
% \end{rem}
\begin{proof}
    %We use induction on $|\bm{n}|$. Let us prove the statement for $\bm{n}$ assuming that it holds true for all indices with lower $|\bm{n}|$.
    The Laurent polynomial $f_{\bm{n}}(z)=\Phi_{\bm{n};\bm{n}}(z) + \Phi_{\bm{n};\bm{n}}(1/z)$ lies in $\operatorname{span}(z^{-|\bm{n}|},\ldots,z^{|\bm{n}|})$ and satisfies $f_{\bm{n}}(z) = f_{\bm{n}}(1/z)$. Therefore (see, e.g., ~\cite[Lem 13.1.4]{OPUC2}) there exists some polynomial $Q_{\bm{n}}(z)$ of degree at most $|\bm{n}|$ such that $Q_{\bm{n}}(z+z^{-1})= \Phi_{\bm{n};\bm{n}}(z) + \Phi_{\bm{n};\bm{n}}(1/z)$.


    Since each $L_j[w^k] = L_j[w^{-k}]$, %all the moments in \eqref{eq:moments linear functional} are real. 
    we are in the setting of Proposition~\ref{rem:Rsymmetry}, so that $\Phi_{\bm{n};\bm{n}}(1/z)=\Phi_{\bm{n};\bm{n}}^*(z)$. 
    Then 
    \begin{align}
    \label{eq:proofInner1}
     M_j[Q_{\bm{n}}(x)x^k]& =  L_j[ Q_{\bm{n}}(w+w^{-1})(w+w^{-1})^k]
    \\
    \label{eq:proofInner2}
    &= 
    L_j[ \Phi_{\bm{n};\bm{n}}(w)(w+w^{-1})^k] 
    +L_j[\Phi^*_{\bm{n};\bm{n}}(w)(w+w^{-1})^k] 
    \\
    \label{eq:proofInner3}
    & 
    =0, \qquad k = 0,\dots,n_j-1,
    \end{align}
    by ~\eqref{eq:moprlII}, and~\eqref{eq:moprlII*}.
    Hence $Q_{\bm{n}}$ satisfies the type II orthogonality conditions~\eqref{eq:MOPRLIIreal} for $\bm{M}$ at the location $\bm{n}$. 

    Now let us consider~\eqref{eq:proofInner1}--\eqref{eq:proofInner2} with $k=n_j$. We get 
    \begin{align}
     M_j[Q_{\bm{n}}(x)x^{n_j}]
     %&=     L_j[ \Phi_{\bm{n};\bm{n}}(w)(w+w^{-1})^{n_j}]     +L_j[\Phi^*_{\bm{n};\bm{n}}(w)(w+w^{-1})^{n_j}] \\
    & 
    =L_j[ \Phi_{\bm{n};\bm{n}}(w) w^{-n_j}] 
    +L_j[\Phi_{\bm{n};\bm{n}}(1/w) w^{n_j}] 
    \\&\label{eq:computation for a}
    = 2L_j[ \Phi_{\bm{n};\bm{n}}(w) w^{-n_j}] ,
    \\&
    \label{eq:lastQnorm}
    = \frac{2\det T_{\bm{n}+\bm{e}_j;\bm{n}}}{\det T_{\bm{n};\bm{n}}},
    \end{align}
    where on the last two steps we used symmetry~\eqref{eq:symmetricFunctionals} and then~\eqref{eq:norm1}. By Proposition~\ref{prop:normality}{(i)} and normality of $(\bm{n}+\bm{e}_j;\bm{n})$, we obtain that the last expression is non-zero.
    Now we are in the position to apply the criterion  
    %~\cite[Lem~3.4]{BCVA} 
    \cite[Thm 2.19]{KVChristoffel} 
    to conclude that $\bm{M}$ is perfect.  
    Note that the $z^{|\bm{n}|}$ coefficient of $Q_{\bm{n}}$ is equal to $1+\alpha_{\bm{n};\bm{n}}$. This proves that $1+\alpha_{\bm{n};\bm{n}}\ne 0$ and that ~\eqref{eq:SzegoPolynomialRelation} holds.  

    Finally, applying the same argument to $f(z)=\Phi_{\bm{n};\bm{n}-\bm{e}_j}(z) + \Phi_{\bm{n};\bm{n}-\bm{e}_j}(1/z)$ (note  that $(\bm{n};\bm{n}-\bm{e}_j)= ((\bm{n}-\bm{e}_j)+\bm{e}_j;\bm{n}-\bm{e}_j)$ is normal by assumption), one can see that ~\eqref{eq:SzegoPolynomialRelation} holds by repeating the argument in~\eqref{eq:proofInner1}--\eqref{eq:proofInner3} and noting that the $z^{|\bm{n}|}$ coefficient of $f$ is 1.
\end{proof}

In the next result we establish the generalization of the Geronimus relations that connect Verblunsky coefficients of $\mu$ and the Jacobi coefficients of $\gamma$, compare with the original~\cite{GeronimusRelations} Geronimus relations~\eqref{eq:GerIntro1}--\eqref{eq:GerIntro2}. %, see also~\cite[Thm 13.1.7]{OPUC2}.

\begin{thm}\label{thm:Geronimus}
    In the setting of the previous theorem, we have
    \begin{alignat}{1}
    \label{eq:Geronimus1}
        % a_{\bm{n},j} & =
        % \frac{(1+\alpha_{\bm{n}-\bm{e}_j;\bm{n}-\bm{e}_j})(1-\alpha_{\bm{n}-\bm{e}_j;\bm{n}}  \beta_{\bm{n};\bm{n}-\bm{e}_j})\rho_{\bm{n};\bm{n},j} }{1+\alpha_{\bm{n};\bm{n}}},
        a_{\bm{n},j} & =
        \frac{(1+\alpha_{\bm{n}-\bm{e}_j;\bm{n}-\bm{e}_j})(1-\alpha_{\bm{n}-\bm{e}_j;\bm{n}}^2
        %\beta_{\bm{n};\bm{n}-\bm{e}_j}
        )\rho_{\bm{n};\bm{n},j} 
        }{1+\alpha_{\bm{n};\bm{n}}},
        \\
        \label{eq:Geronimus2}
        b_{\bm{n},j} & =
        \sum_{\ell=1}^r \rho_{\bm{n};\bm{n},\ell}  \gamma_{\bm{n};\bm{n}}^{\ell j}+
        \alpha_{\bm{n};\bm{n}-\bm{e}_j} - \alpha_{\bm{n}+\bm{e}_j;\bm{n}}
        - \alpha_{\bm{n};\bm{n}} \alpha_{\bm{n}-\bm{e}_j;\bm{n}}
        - \alpha_{\bm{n};\bm{n}} \alpha_{\bm{n}+\bm{e}_j;\bm{n}} .
        % - \alpha_{\bm{n};\bm{n}} \beta_{\bm{n};\bm{n}-\bm{e}_j}
        % - \beta_{\bm{n};\bm{n}} \alpha_{\bm{n}+\bm{e}_j;\bm{n}} .
    \end{alignat}
\end{thm}
\begin{rem}
    When $\alpha_{\bm{n};\bm{n}} \neq 0$ we can eliminate $\gamma_{\bm{n};\bm{n}}^{\ell j}$ to get
    \begin{equation}
        \label{eq:Geronimus2.1}
        b_{\bm{n},j}  =
        \frac{1}{\alpha_{\bm{n};\bm{n}}}\bigg(-\alpha_{\bm{n}+\bm{e}_j;\bm{n}}+\sum_{\ell=1}^r \alpha_{\bm{n}+\bm{e}_{\ell};\bm{n}}\rho_{\bm{n};\bm{n},\ell}\bigg)  +
        \alpha_{\bm{n},\bm{n}-\bm{e}_j} - \alpha_{\bm{n}+\bm{e}_j,\bm{n}}
        - \alpha_{\bm{n};\bm{n}} \alpha_{\bm{n}-\bm{e}_j;\bm{n}} .
        %- \alpha_{\bm{n};\bm{n}} \beta_{\bm{n};\bm{n}-\bm{e}_j} .
    \end{equation}
\end{rem}
\begin{proof}
    From~\eqref{eq:one nnr} and~\eqref{eq:MOPRLIIreal}, and then from~\eqref{eq:type 1 recurrence relation} and~\eqref{eq:MOPRLIreal}--\eqref{eq:normalization type I moprl}, one gets
    \begin{align}
    \label{eq:avsM}
        a_{\bm{n},j} 
        &= \frac{ M_j[P_{\bm{n}}(x)x^{n_j} ]}{ M_j[P_{\bm{n}-\bm{e}_j}(x) x^{n_j - 1}] },
        \\
    \label{eq:bvsM}
        b_{\bm{n},j}
        % &= \sum_{\ell=1}^r M_\ell[A_{\bm{n},\ell}(x) x^{\abs{\bm{n}}}] - \sum_{\ell=1}^r M_\ell[A_{\bm{n}-\bm{e}_j,\ell}(x) x^{\abs{\bm{n}}-1}].
        & = k_{|\bm{n}|-1}(P_{\bm{n}}) - k_{|\bm{n}|}(P_{\bm{n}+\bm{e}_j}),
    \end{align}
    where $k_j(P)$ denotes the $z^j$ coefficient of a polynomial $P$.

    {Using ~\eqref{eq:avsM}, \eqref{eq:SzegoPolynomialRelation}, and \eqref{eq:computation for a}, we obtain
    \begin{equation}
        a_{\bm{n},j} = \frac{1+\alpha_{\bm{n}-\bm{e}_j;\bm{n}-\bm{e}_j}}{1+\alpha_{\bm{n};\bm{n}}}\frac{2L_j[\Phi_{\bm{n};\bm{n}}(z)z^{-n_j}]}{L_j[\Phi_{\bm{n}-\bm{e}_j;\bm{n}}(z)z^{-n_j+1}]}\frac{L_j[\Phi_{\bm{n}-\bm{e}_j;\bm{n}}(z)z^{-n_j+1}]}{2L_j[\Phi_{\bm{n}-\bm{e}_j;\bm{n}-\bm{e}_j}(z)z^{-n_j+1}]}.
    \end{equation}
    Now \eqref{eq:Geronimus1} follows from \eqref{eq:rho} and \eqref{eq:one minus alphabeta}, and the symmetry \eqref{eq:szego symmetry}.
    }

    % \red{Using~\eqref{eq:avsM}, \eqref{eq:SzegoPolynomialRelation}, and \eqref{eq:lastQnorm}, we obtain
    % \begin{equation}
    %     a_{\bm{n},j} =\frac{1+\alpha_{\bm{n}-\bm{e}_j;\bm{n}-\bm{e}_j}}{1+\alpha_{\bm{n};\bm{n}}}
    %     \frac{2\det T_{\bm{n}+\bm{e}_j;\bm{n}}}{\det T_{\bm{n};\bm{n}}}
    %     \frac{\det T_{\bm{n}-\bm{e}_j;\bm{n}-\bm{e}_j}}{2\det T_{\bm{n};\bm{n}-\bm{e}_j}}.
    % \end{equation}
    % Applying the relations~\eqref{eq:rho},~\eqref{eq:kappa}, and~\eqref{eq:norm1}, this can be rewritten as
    % \begin{equation}
    %     a_{\bm{n},j} =\frac{1+\alpha_{\bm{n}-\bm{e}_j;\bm{n}-\bm{e}_j}}{1+\alpha_{\bm{n};\bm{n}}} \rho_{\bm{n};\bm{n},j} \frac{\kappa_{\bm{n};\bm{n}-\bm{e}_j,j}}{\kappa_{\bm{n};\bm{n},j}}.
    % \end{equation}
    % Now we use~\eqref{eq:kappaAlphaBeta} to obtain~\eqref{eq:Geronimus1}.}
    
    
    % Applying the relations~\eqref{eq:symmReal} produces
    % \begin{alignat*}{1}
    %     a_{\bm{n},j}
    %     &=
    %     \frac{1+\alpha_{\bm{n}-\bm{e}_j;\bm{n}-\bm{e}_j}}{1+\alpha_{\bm{n};\bm{n}}}
    %     \frac{\det T_{\bm{n}+\bm{e}_j;\bm{n}} \det T_{\bm{n}-\bm{e}_j;\bm{n}} }{(\det T_{\bm{n};\bm{n}})^2}
    %     \frac{\det T_{\bm{n};\bm{n}} \det T_{\bm{n}-\bm{e}_j;\bm{n}-\bm{e}_j}}{\det T_{\bm{n}-\bm{e}_j;\bm{n}} \det T_{\bm{n};\bm{n}-\bm{e}_j}}
    %     % \\
    %     % &= 
    %     % \frac{1+\alpha_{\bm{n}-\bm{e}_j;\bm{n}-\bm{e}_j}}{1+\alpha_{\bm{n};\bm{n}}}
    %     % \rho_{\bm{n};\bm{n},j}
    %     % %\rho_{\bm{n};\bm{n}-\bm{e}_j,j}
    %     % \frac{\det T_{\bm{n};\bm{n}} \det T_{\bm{n}-\bm{e}_j;\bm{n}-\bm{e}_j}}{\det T_{\bm{n}-\bm{e}_j;\bm{n}} \det T_{\bm{n};\bm{n}-\bm{e}_j}}
    %     \\
    %     &= 
    %     \pm\frac{1+\alpha_{\bm{n}-\bm{e}_j;\bm{n}-\bm{e}_j}}{1+\alpha_{\bm{n};\bm{n}}}
    %     \rho_{\bm{n};\bm{n},j}
    %     %\rho_{\bm{n};\bm{n}-\bm{e}_j,j}
    %     \frac{\det T_{\bm{n};\bm{n}} \det T_{\bm{n}-\bm{e}_j;\bm{n}-\bm{e}_j}}{(\det T_{\bm{n}-\bm{e}_j;\bm{n}})^2}
    % \end{alignat*}
    % Ideally, the last ratio should simplify to something like $\rho_{\bm{n}-\bm{e}_j;\bm{n},j}$ but it's not quite right...

    % In other words, see~\eqref{eq:norm1}, we need to be able to express
    % $$
    % \frac{L_j[\Phi_{\bm{n}-\bm{e}_j;\bm{n}}(z)z^{-n_j+1}]}{L_j[\Phi_{\bm{n}-\bm{e}_j;\bm{n}-\bm{e}_j}(z)z^{-n_j+1}]}
    % $$
    % in terms of $\alpha$'s, $\beta$'s, $\rho$'s, $\sigma$'s, $\gamma$'s.

    % Actually, using~\eqref{eq:kappa}, we can write
    % $$
    % a_{\bm{n},j} = \frac{1+\alpha_{\bm{n}-\bm{e}_j;\bm{n}-\bm{e}_j}}{1+\alpha_{\bm{n};\bm{n}}} \rho_{\bm{n};\bm{n},j} \frac{\kappa_{\bm{n};\bm{n}-\bm{e}_j,j}}{\kappa_{\bm{n};\bm{n},j}},
    % $$
    % which is not hard to reduce to the correct expression if $r=1$. Of course this is not what we want since $\kappa$'s are not our recurrence coefficients, but if we cannot get anything better, then we should keep this form, I think. I expect $b_{\bm{n},j}$ to be expressable in terms of $\kappa$'s as well.

    %Here we follow the trick from the proof of~\cite[Thm 13.1.7]{OPUC2}. 
    For the second relation we follow the trick from the proof of~\cite[Thm 13.1.7]{OPUC2}. From~\eqref{eq:bvsM} and \eqref{eq:SzegoPolynomialRelation2}, we get
    \begin{align}
        b_{\bm{n},j} & =  (k_{|\bm{n}|-1}(\Phi_{\bm{n};\bm{n}-\bm{e}_j}) +\alpha_{\bm{n},\bm{n}-\bm{e}_j}) - (k_{|\bm{n}|}(\Phi_{\bm{n}+\bm{e}_j,\bm{n}})+\alpha_{\bm{n}+\bm{e}_j,\bm{n}})
        \\
        & =
        (k_{|\bm{n}|-1}(\Phi_{\bm{n};\bm{n}}) - k_{|\bm{n}|}(\Phi_{\bm{n}+\bm{e}_j,\bm{n}}) )
        +(\alpha_{\bm{n},\bm{n}-\bm{e}_j} - \alpha_{\bm{n}+\bm{e}_j,\bm{n}})
        - \alpha_{\bm{n};\bm{n}} \beta_{\bm{n};\bm{n}-\bm{e}_j},
    \end{align}
    where we used~\eqref{eq:SzegoII4} in the last line. The result now follows since we can compute $k_{|\bm{n}|-1}(\Phi_{\bm{n};\bm{n}}) - k_{|\bm{n}|}(\Phi_{\bm{n}+\bm{e}_j;\bm{n}})$ directly from \eqref{eq:SzegoII9} to get \eqref{eq:Geronimus2}, and then use symmetry \eqref{eq:szego symmetry}. For the remark, we instead compute $k_{|\bm{n}|-1}(\Phi_{\bm{n};\bm{n}}) - k_{|\bm{n}|}(\Phi_{\bm{n}+\bm{e}_j;\bm{n}})$ through \eqref{eq:generalized three term}.

    %Now we use~\cite[Cor. 3.4(i)]{KVMOPUC} to see that
    %\begin{equation}
    %    k_{|\bm{n}|-1}(\Phi_{\bm{n};\bm{n}}) - k_{|\bm{n}|}(\Phi_{\bm{n}+\bm{e}_j,\bm{n}}) 
    %    =
    %    -
    %    \frac{\beta_{\bm{n};\bm{n}}-\sum_{\ell=1}^r \rho_{\bm{n}+\bm{e}_j;\bm{n},\ell} \beta_{\bm{n}+\bm{e}_j-\bm{e}_\ell;\bm{n}}}{\beta_{\bm{n}+\bm{e}_j;\bm{n}}}
    %\end{equation}
    %But
    %\begin{equation}
    %    \beta_{\bm{n}+\bm{e}_j-\bm{e}_\ell;\bm{n}}
    %    =
    %    \beta_{\bm{n};\bm{n}} + \beta_{\bm{n}+\bm{e}_j;\bm{n}}
    %    \gamma_{\bm{n}-\bm{e}_\ell;\bm{n}}^{j\ell},
    %\end{equation}
    % by~\eqref{eq:another comp}, so that
    % \begin{equation}
    %    k_{|\bm{n}|-1}(\Phi_{\bm{n};\bm{n}}) - k_{|\bm{n}|}(\Phi_{\bm{n}+\bm{e}_j,\bm{n}}) 
    %    =
    %    -
    %    \frac{\beta_{\bm{n};\bm{n}}-\sum_{\ell=1}^r \rho_{\bm{n}+\bm{e}_j;\bm{n},\ell} (\beta_{\bm{n};\bm{n}} + \beta_{\bm{n}+\bm{e}_j;\bm{n}}
    %    \gamma_{\bm{n}-\bm{e}_\ell;\bm{n}}^{j\ell})}{\beta_{\bm{n}+\bm{e}_j;\bm{n}}}.
    %\end{equation}

    %Now use~\eqref{eq:compatibility condition 2} to get
    %\begin{multline}
    %    -
    %    \frac{\beta_{\bm{n};\bm{n}}
    %    \alpha_{\bm{n}+\bm{e}_j;\bm{n}} 
    %    \beta_{\bm{n}+\bm{e}_j;\bm{n}} 
    %    -\sum_{\ell=1}^r \rho_{\bm{n}+\bm{e}_j;\bm{n},\ell}  \beta_{\bm{n}+\bm{e}_j;\bm{n}}
    %    \gamma_{\bm{n}-\bm{e}_\ell;\bm{n}}^{j\ell}}{\beta_{\bm{n}+\bm{e}_j;\bm{n}}}
    %    \\
    %    =
    %    -
    %    \beta_{\bm{n};\bm{n}}
    %    \alpha_{\bm{n}+\bm{e}_j;\bm{n}} 
    %    +\sum_{\ell=1}^r \rho_{\bm{n}+\bm{e}_j;\bm{n},\ell}  \gamma_{\bm{n}-\bm{e}_\ell;\bm{n}}^{j\ell}.
    %\end{multline}
\end{proof}

% \begin{prop}
%     If $(\bm{n};\bm{n})$ and $(\bm{m};\bm{n})$ are normal for $\bm{\mu}$, for some index $\bm{m}$ such that $\bm{n} \lneq \bm{m} \leq \bm{n} + \bm{1}$, then $1 + \alpha_{\bm{n};\bm{n}} \neq 0$.
% \end{prop}
% \begin{proof}
%     If $\alpha_{\bm{n};\bm{n}} = -1$ then we have $\Phi_{\bm{n};\bm{n}}^* = -\Phi_{\bm{n};\bm{n}}$, by \eqref{eq:SzegoPolynomialRelation}. Then $\Phi_{\bm{n};\bm{n}}$ would satisfy all the orthogonality conditions for the index $(\bm{m};\bm{n})$, which would imply that $(\bm{m};\bm{n})$ is not normal.
% \end{proof}



% For the type I polynomials we get similar relations. 
% %, but $\bm{\Xi}_{\bm{n};\bm{n}} + \bm{\Xi}_{\bm{n};\bm{n}}^\sharp$ will not be a constant multiple of $\bm{A}_{\bm{n}}$, but rather a linear combination of the vectors $\set{\bm{A}_{\bm{n}+\bm{e}_j}}_{j = 1}^r$.

% \begin{thm}
%      Assume that $\bm{L} = (L_1,\dots,L_r)$ is a system of linear functionals on $\T$ satisfying~\eqref{eq:symmetricFunctionals} and let $\bm{M}$ be defined by $M_j=\operatorname{Sz}^{-1}(L_j)$, $j = 1,\dots,r$.

%     %Assume $\bm{L} = \operatorname{Sz}(\bm{M})$ and that 
%     If $(\bm{n};\bm{n})$ and $(\bm{n}-\bm{e}_j;\bm{n})$ are normal for $\bm{L}$ for all $\bm{n}\in\N^r$ and $j=1,\ldots,r$, then $\bm{M}$ is perfect (that is, every $\bm{n}\in\N^r$ is normal for $\bm{M}$), and
%     \begin{align}\label{eq:SzegoPolynomialRelationI}
%           \bm{A}_{\bm{n}}(z+z^{-1}) & = \frac{1}{(1-\beta_{\bm{n};\bm{n}})} \Big( \bm\Xi_{\bm{n};\bm{n}}(z) + \bm\Xi_{\bm{n};\bm{n}}(1/z) \Big) 
%           \\
%           \label{eq:SzegoPolynomialRelation2I}
%           %& = \Phi_{\bm{n};\bm{m}}(z) + \Phi_{\bm{n};\bm{m}}(1/z)
%           & ?= \Phi_{\bm{n};\bm{n}-\bm{e}_j}(z) + \Phi_{\bm{n};\bm{n}-\bm{e}_j}(1/z).
%         \end{align}
%         %for any $\bm{m}$ such that $\bm{n} - \bm{1} \leq \bm{m} \lneq \bm{n}$.
%         Furthermore, $\alpha_{\bm{n};\bm{n}}$ is necessarily $\ne -1$. 
%     %where $\alpha_{\bm{n};\bm{n}}$ is the $z^{-\abs{\bm{n}}}$-coefficient of $\Phi_{\bm{n};\bm{n}}$.
% \end{thm}
% \begin{proof}
%     The vector $\bm{B}_{\bm{n}}(z)=\bm\Xi_{\bm{n};\bm{n}}(z) + \bm{\Xi}_{\bm{n};\bm{n}}(1/z)$ is a vector of Laurent polynomials with $B_{\bm{n},j}(z) \in \operatorname{span}(z^{-n_j},\ldots,z^{n_j})$ that satisfy $B_{\bm{n},j}(z) = B_{\bm{n},j}(1/z)$. Therefore there exists a vector $\bm{C}_{\bm{n}}(z) = (C_{\bm{n},1}(z),\ldots,C_{\bm{n},r}(z))$ of polynomials such that $\deg C_{\bm{n},j}(z) \le n_j$ and $\bm{C}_{\bm{n}}(z+z^{-1})= \bm\Xi_{\bm{n};\bm{n}}(z) + \bm{\Xi}_{\bm{n};\bm{n}}(1/z)$.


%     Since each $L_j[w^k] = L_j[w^{-k}]$, %all the moments in \eqref{eq:moments linear functional} are real. 
%     we are in the setting of Remark~\ref{rem:Rsymmetry}, so that $\Xi_{\bm{n};\bm{n}}(1/z)=\Xi_{\bm{n};\bm{n}}^*(z)$, and, therefore,
%     \begin{align}
%     \label{eq:proofInner1I}
%      \sum_{j=1}^r M_j[C_{\bm{n},j}(x) x^k]& =  \sum_{j=1}^r L_j[ C_{\bm{n},j}(w+w^{-1})(w+w^{-1})^k]
%     \\
%     \label{eq:proofInner2I}
%     &= 
%     \sum_{j=1}^r L_j[ \bm\Xi_{\bm{n};\bm{n},j}(w)(w+w^{-1})^k]
%     +
%     \sum_{j=1}^r L_j[ \bm\Xi^*_{\bm{n};\bm{n},j}(w)(w+w^{-1})^k]
%     \\
%     \label{eq:proofInner3I}
%     & 
%     =0, \qquad k = 0,\dots,|\bm{n}|-1,
%     \end{align}
%     by ~\eqref{eq:moprlI}, and~\eqref{eq:moprlI*}.
%     Hence $\bm{C}_{\bm{n}}$ satisfies the type I orthogonality conditions~\eqref{eq:MOPRLIreal} for $\bm{M}$ at the location $\bm{n}+\bm{1}$, where $\bm{1} = (1,\ldots,1)$. 

%     Now let us consider~\eqref{eq:proofInner1}--\eqref{eq:proofInner2} with $k=n_j$. We get 
%     \begin{align}
%      M_j[Q_{\bm{n}}(x)x^{n_j}]
%      %&=     L_j[ \Phi_{\bm{n};\bm{n}}(w)(w+w^{-1})^{n_j}]     +L_j[\Phi^*_{\bm{n};\bm{n}}(w)(w+w^{-1})^{n_j}] \\
%     & 
%     =L_j[ \Phi_{\bm{n};\bm{n}}(w) w^{-n_j}] 
%     +L_j[\Phi_{\bm{n};\bm{n}}(1/w) w^{n_j}] 
%     \\&
%     = 2L_j[ \Phi_{\bm{n};\bm{n}}(w) w^{-n_j}] ,
%     \\&
%     = \frac{2\det T_{\bm{n}+\bm{e}_j;\bm{n}}}{\det T_{\bm{n};\bm{n}}},
%     \end{align}
%     where on the last two steps we used symmetry~\eqref{eq:symmetricFunctionals} and then~\eqref{eq:norm1}. By Proposition~\ref{prop:normality}{\it (i)} and normality of $(\bm{n}+\bm{e}_j;\bm{n})$, we obtain that the last expression is non-zero.
%     Now we are in the position to apply the criterion from~\cite[Lem~3.4]{BCVA} %(or \cite[Thm 2.19]{KVChristoffel}) 
%     to conclude that $\bm{M}$ is perfect.  
%     Note that the $z^{|\bm{n}|}$ coefficient of $Q_{\bm{n}}$ is equal to $1+\alpha_{\bm{n};\bm{n}}$. This proves that $1+\alpha_{\bm{n};\bm{n}}\ne 0$ and that ~\eqref{eq:SzegoPolynomialRelation} holds.  

%     Finally, applying the same argument to $f(z)=\Phi_{\bm{n};\bm{n}-\bm{e}_j}(z) + \Phi_{\bm{n};\bm{n}-\bm{e}_j}(1/z)$ (note  that $(\bm{n};\bm{n}-\bm{e}_j)= ((\bm{n}-\bm{e}_j)+\bm{e}_j;(\bm{n}-\bm{e}_j))$ is normal by assumption), one can see that ~\eqref{eq:SzegoPolynomialRelation} holds by repeating the argument in~\eqref{eq:proofInner1}--\eqref{eq:proofInner3} and noting that the $z^{|\bm{n}|}$ coefficient of $f$ is 1.
% \end{proof}



% \begin{prop}\label{thm:type I szego mapping}
%     Assume $\bm{n}$ is normal for $\bm{\gamma}$ and let $\bm{\mu} = \operatorname{Sz}(\bm{\gamma})$. If $(\bm{m};\bm{n})$ is normal for $\bm{\gamma}$, and $\bm{n} \leq \bm{m} \leq \bm{n} + \bm{1}$, then we have
%     \begin{equation}\label{eq:szego relation type I}
%         \bm{\Xi}_{\bm{m};\bm{n}}(z) + \bm{\Xi}_{\bm{m};\bm{n}}^\sharp(z) \in \operatorname{span}\set{\bm{A}_{\bm{n}+\bm{e}_k}(z+z^{-1})}_{k = 1}^r.
%     \end{equation}
% \end{prop}

% \begin{proof}
%     By similar arguments to the proof of Theorem \ref{thm:szego mapping type II}, the polynomials $Q_j$ defined by $Q_j(z+z^{-1}) = \Xi_{\bm{m};\bm{n},j}(z) + \Xi_{\bm{m};\bm{n},j}^\sharp(z)$ have the property
%     \begin{equation}\label{eq:type I szego mapping proof}
%         \sum_{j = 1}^r \int_{-2}^{2} Q_j(x)x^{-k}d\gamma_j(x) = 0, \qquad k = 0,\dots,\abs{\bm{n}}-1.
%     \end{equation}
%     $Q_j$ has degree at most $n_j$, so we end up with a system of equations for the coefficients of $\bm{Q} = (Q_1,\dots,Q_r)$, with $\abs{\bm{n}}$ equations and $\abs{\bm{n}}+r$ unknowns. The system has coefficient matrix equal to the coefficient matrix of \eqref{eq:MOPRLIreal} with $r$ columns added. {Hence the solution space has dimension $r$. }
%     %Hence $\bm{Q}$ %\rkno{sits in a}
%     %\rk{belongs to the}  solution space of dimension $r$. 
%     The vectors $\bm{A}_{\bm{n}+\bm{e}_k}$ 
%     %\rkno{have the same orthogonality relations and degree constraints}
%     {belong to the same subspace}, for each $k = 1,\dots,r$. They are linearly independent since $A_{\bm{n}+\bm{e}_k,k}$ has degree $n_k$ {(see \cite[Cor 23.1.1]{Ismail})} and every other $A_{\bm{n}+\bm{e}_j,k}$, $j\ne k$, has smaller degree. %\rkno{$\deg{A_{\bm{n}+\bm{e}_k,k}} = n_k$ follows since otherwise $\bm{Q} = \bm{A}_{\bm{n}+\bm{e}_k}$ would satisfy \eqref{eq:type I szego mapping proof} along with all the degree constraints of $\bm{A}_{\bm{n}}$, which gives a system of equations with the same coefficients matrix as \eqref{eq:MOPRLIreal}. By normality \rk{of $\bm{n}$ for $\bm{\gamma}$} we then must have $\bm{A}_{\bm{n}+\bm{e}_k} = \bm{0}$, which contradicts Remark}
%     Therefore~\eqref{eq:szego relation type I} follows.
%     %\ref{rem:type I moprl no normalization}. 
%     %we would have $\bm{A}_{\bm{n}+\bm{e}_j} = \bm{A}_{\bm{n}}$, which contradicts normality of $\bm{n}$.
% \end{proof}

% \begin{rem}\label{rem:type I moprl no normalization}
%     Any choice of vectors $\bm{A}_{\bm{n}+\bm{e}_1},\dots,\bm{A}_{\bm{n}+\bm{e}_r}$ works in \eqref{eq:szego relation type I}, but each choice leads to a different set of $c_j$'s, which is why normality of the indices $\bm{n}+\bm{e}_1,\dots,\bm{n}+\bm{e}_r$ does not need to be assumed. Normality of $\bm{n}$ is needed for the linear independence of these vectors.
% \end{rem}


\bibsection

\begin{biblist}[\small]

\bib{Aptekarev}{article}{
   author={Aptekarev, A.I.},
   title={Multiple orthogonal polynomials},
   journal={J. Comput. Appl. Math.},
   volume={99},
   year={1998},
   pages={423--447},
}

\bib{Baker}{book}{
    AUTHOR = {Baker, G.A.},
    AUTHOR = {Graves-Morris, P.},
     TITLE = {Pad\'e{} approximants},
    SERIES = {Encyclopedia of Mathematics and its Applications},
    VOLUME = {59},
   EDITION = {Second},
 PUBLISHER = {Cambridge University Press, Cambridge},
      YEAR = {1996},
     PAGES = {xiv+746},
}


\bib{BCVA}{article}{
   author={Beckermann, B.},
   author={Coussement, J.},
   author={Van Assche, W.},
   title={Multiple Wilson and Jacobi-Pi\~neiro polynomials},
   journal={J. Approx. Theory},
   volume={132},
   date={2005},
   number={2},
   pages={155--181}
}

\bib{Bultheel}{book}{
    AUTHOR = {Bultheel, A.},
     TITLE = {Laurent series and their {P}ad\'e{} approximations},
    SERIES = {Operator Theory: Advances and Applications},
    VOLUME = {27},
 PUBLISHER = {Birkh\"auser Verlag, Basel},
      YEAR = {1987},
     PAGES = {xii+274},
}

\bib{Chihara}{book}{
   author={Chihara, T.S.},
   title={An Introduction to Orthogonal Polynomials},
   isbn={9780486479293},
   series={Mathematics and Its Applications},
   volume={13},
   publisher={Gordon and Breach Science Publishers, Inc.},
   year={1978},
}

\bib{MOPUC2}{article}{
   author={Cruz-Barroso, R.},
   author={D\'iaz Mendoza, C.},
   author={Orive, R.},
   title={Multiple orthogonal polynomials on the unit circle. Normality and recurrence relations},
   journal={J. Comput. Appl. Math.},
   volume={284},
   year={2015},
   pages={115--132},
}

%\bib{DaeKui}{article}{
%    AUTHOR = {Daems, E.},
%    AUTHOR = {Kuijlaars, A.B.J.},
%     TITLE = {A {C}hristoffel-{D}arboux formula for multiple orthogonal
%              polynomials},
%   JOURNAL = {J. Approx. Theory},
%  FJOURNAL = {Journal of Approximation Theory},
%    VOLUME = {130},
%      YEAR = {2004},
%    NUMBER = {2},
%     PAGES = {190--202},
%      ISSN = {0021-9045,1096-0430},
   %MRCLASS = {42C05},
  %MRNUMBER = {2100703},
%MRREVIEWER = {K.\ Pan},
       %DOI = {10.1016/j.jat.2004.07.003},
       %URL = {https://doi.org/10.1016/j.jat.2004.07.003},
%}

\bib{DerSim18}{article}{
    AUTHOR = {Derevyagin, M.}
    AUTHOR = {Simanek, B.},
     TITLE = {Asymptotics for polynomials orthogonal in an indefinite
              metric},
   JOURNAL = {J. Math. Anal. Appl.},
  FJOURNAL = {Journal of Mathematical Analysis and Applications},
    VOLUME = {460},
      YEAR = {2018},
    NUMBER = {2},
     PAGES = {777--793},
}

\bib{GeronimusRelations}{article}{
    AUTHOR = {Geronimus, Ya.L.},
     TITLE = {On the trigonometric moment problem},
   JOURNAL = {Ann. of Math. (2)},
  FJOURNAL = {Annals of Mathematics. Second Series},
    VOLUME = {47},
      YEAR = {1946},
     PAGES = {742--761},
  MRNUMBER = {18265},
}

\bib{GeronimusBook}{book}{
    AUTHOR = {Geronimus, Ya.L.},
     TITLE = {Teoriya ortogonal\cprime nyh mnogo\v clenov},
 PUBLISHER = {Gosudarstv. Izdat. Tehn.-Teor. Lit., Moscow-Leningrad},
      YEAR = {1950},
     PAGES = {164},
}


\bib{HueMan}{article}{
    AUTHOR={Huertas, E.J.},
    AUTHOR={Ma\~{n}as, M.},
     TITLE = {Mixed-type multiple orthogonal {L}aurent polynomials on the
              unit circle},
   JOURNAL = {J. Comput. Appl. Math.},
  FJOURNAL = {Journal of Computational and Applied Mathematics},
    VOLUME = {475},
      YEAR = {2026},
     PAGES = {Paper No. 117037},
}


\bib{Ismail}{book}{
   author={Ismail, M.E.H.},
   title={Classical and Quantum Orthogonal
Polynomials in One Variable},
   isbn={9780521782012},
   series={Encyclopedia of Mathematics and its Applications},
   Volume={98},
   publisher={Cambridge University Press},
   year={2005},
}

\bib{JNT}{article}{
    author={Jones, W.B.},
    author={Nj\r astadt, O.},
    author={Thron, W.J.},
    title={Moment Theory, Orthogonal Polynomials, Quadrature, and Continued Fractions Associated with the unit Circle},
    journal={Bulletin of the London Mathematical Society},
    volume={21},
    number={2},
    year={1989},
    pages={113--152}
}

\bib{KNik}{article}{
    AUTHOR={Kozhan, R.},
    TITLE={Nikishin Systems on the Unit Circle},
  journal={},
   volume={},
   date={},
   number={},
   pages={under submission, arXiv:2410.20813},
}


\bib{KVMLOPUC1}{article}{
    AUTHOR={Kozhan, R.},
    AUTHOR={Vaktn\"as, M.},
    TITLE={Angelesco and AT Systems on the Unit Circle},
    journal={},
    volume={},
    date={},
    number={},
    pages={under submission, arXiv:2410.12094},
}

\bib{KVChristoffel}{article}{
   author={Kozhan, R.},
   author={Vaktn\"{a}s, M.},
   title={Christoffel transform and multiple orthogonal polynomials},
      JOURNAL = {J. Comput. Appl. Math.},
    VOLUME = {476},
      YEAR = {2026},
     PAGES = {Paper No. 117121},
}

\bib{KVMOPUC}{article}{
    AUTHOR={Kozhan, R.},
    AUTHOR={Vaktn\"as, M.},
    TITLE={Szeg\H{o} recurrence for multiple orthogonal polynomials on the unit circle},
    JOURNAL={Proc. Amer. Math. Soc.},
    VOLUME={152},
    NUMBER={11},
    YEAR={2024},
    PAGES={2983-2997},
    ISSN={1088-6826,0002-9939},
}

\bib{KVInterlacing}{article}{
    AUTHOR={Kozhan, R.},
    AUTHOR={Vaktn\"as, M.},
    TITLE={Zeros of multiple orthogonal polynomials: location and interlacing},
    journal={Bull. of London Math. Soc.},
    volume={},
    date={},
    number={},
    pages={to appear, arXiv:2503.15122},
}


 
\bib{Kui}{article}{
   AUTHOR = {Kuijlaars, A.B.J.},
     TITLE = {Multiple orthogonal polynomial ensembles},
    JOURNAL = {Recent trends in orthogonal polynomials and approximation
              theory, Contemp. Math., Amer. Math. Soc., Providence, RI},
    VOLUME = {507},
     PAGES = {155--176},
      YEAR = {2010},
      ISBN = {978-0-8218-4803-6},
}

\bib{Applications}{article}{
   author={Mart\'inez-Finkelshtein, A.},
   author={Van Assche, W.},
   title={WHAT IS...A Multiple Orthogonal Polynomial?},
   journal={Not. Am. Math. Soc.},
   volume={63},
   year={2016},
   pages={1029--1031},
}

\bib{MOPUC1}{article}{
   author={M\'inguez Ceniceros, J.},
   author={Van Assche, W.},
   title={Multiple orthogonal polynomials on the unit circle},
   journal={Constr. Approx.},
   volume={28},
   year={2008},
   pages={173--197},
}

\bib{Nikishin}{book}{
    AUTHOR = {Nikishin, E.M.},
    AUTHOR = {Sorokin, V.N.},
     TITLE = {Rational approximations and orthogonality},
    SERIES = {Translations of Mathematical Monographs},
    VOLUME = {92},
      NOTE = {Translated from the Russian by Ralph P. Boas},
 PUBLISHER = {American Mathematical Society, Providence, RI},
      YEAR = {1991},
     PAGES = {viii+221},
      ISBN = {0-8218-4545-4},
   %MRCLASS = {30E10 (31A15 41A21 42C05)},
  %MRNUMBER = {1130396},
       %DOI = {10.1090/mmono/092},
       %URL = {https://doi.org/10.1090/mmono/092},
}

\bib{PeherstorferSteinbauer}{article}{
    author={Peherstorfer, F.},
    author={Steinbauer, R.},
    title={Characterization of orthogonal polynomials with respect to a functional},
    journal={Journal of Computational and Applied Mathematics},
    volume={65},
    year={1995},
    pages={339-355},
}

\bib{OPUC1}{book}{
   author={Simon, B.},
   title={Orthogonal Polynomials on the Unit Circle, Part 1: Classical Theory},
   isbn={0-8218-3446-0},
   series={Colloquium Lectures},
   Volume={54},
   publisher={American Mathematical Society},
   year={2004},
}

\bib{OPUC2}{book}{
   author={Simon, B.},
   title={Orthogonal Polynomials on the Unit Circle, Part 2: Spectral Theory},
   isbn={978-0-8218-4864-7},
   series={Colloquium Lectures},
   Volume={54},
   publisher={American Mathematical Society},
   year={2005},
}


\bib{SimonL2}{book}{
   author={Simon, B.},
   title={Szeg\H{o}'s Theorem and Its Descendants: Spectral Theory for $L^2$ Perturbations of Orthogonal Polynomials},
   isbn={9780691147048},
   series={Porter Lectures},
   publisher={Princeton University Press},
   year={2011},
}

\bib{SzegoBook}{book}{
    AUTHOR = {Szeg\H o, G.},
     TITLE = {Orthogonal polynomials},
    SERIES = {American Mathematical Society Colloquium Publications},
    VOLUME = {Vol. XXIII, 4th ed},
   %EDITION = {Fourth},
 PUBLISHER = {American Mathematical Society, Providence, RI},
      YEAR = {1975},
     PAGES = {xiii+432},
}

\bib{Szego}{article}{
    author={Szeg\H{o}, G.},
    title={\"Uber die Entwickelung einer analytischen Funktion nach den Polynomen eines Orthogonalsystems},
    issn={0025-5831},
    journal={Mathematische Annalen},
    volume={82},
    year={1921},
    pages={188-212},
}
	
\bib{NNRR}{article}{
   author={Van Assche, W.},
   title={Nearest neighbor recurrence relations for multiple
orthogonal polynomials},
   journal={J. Approx. Theory},
   volume={163},
   year={2011},
   pages={1427--1448},
}


\end{biblist}
%\end{thebibliography}


%%%%%%%%%%%%%%%%%%%%%%%%%%%%%%%%%%%%%%%%%%%%%%%%%%%%%%%%


\end{document}
